\documentclass[11pt,a4paper,oneside]{book}
\usepackage[margin=20mm,top=22mm,bottom=22mm]{geometry}
\usepackage{fontspec}
\usepackage{amsmath}
\usepackage[math-style=ISO]{unicode-math}
\usepackage{polyglossia}
\setdefaultlanguage{ukrainian}
\setotherlanguage{english}
\directlua{
  luaotfload.add_fallback("fallbacks", {"DejaVu Sans:mode=harf;", "FreeSerif:mode=harf;", "DejaVu Math TeX Gyre:mode=harf;"})
}
\setmainfont{DejaVu Serif}[RawFeature={fallback=fallbacks}]
\setsansfont{DejaVu Sans}[RawFeature={fallback=fallbacks}]
\setmonofont{DejaVu Sans Mono}[Scale=0.85,RawFeature={fallback=fallbacks}]
\setmathfont{DejaVu Math TeX Gyre}
\usepackage{microtype}
\usepackage{needspace}
\usepackage{booktabs,longtable,array}
\usepackage{enumitem}
\setlist{nosep,leftmargin=*}
\usepackage{xcolor}
\definecolor{accent}{HTML}{1F4E79}
\definecolor{warncol}{HTML}{B45309}
\usepackage{titlesec}
\titleformat{\chapter}[hang]{\Large\bfseries\sffamily\color{accent}}{\thechapter.}{0.6em}{}
\titlespacing*{\chapter}{0pt}{-10pt}{10pt}
\titleformat{\section}{\normalsize\bfseries\sffamily\color{accent}}{}{0pt}{}
\titlespacing*{\section}{0pt}{8pt}{3pt}
\titleformat{\part}[display]{\centering\Huge\bfseries\sffamily\color{accent}}{\partname~\thepart}{12pt}{}
\usepackage{fancyhdr}
\pagestyle{fancy}\fancyhf{}
\fancyhead[L]{\small\sffamily\leftmark}\fancyhead[R]{\small\sffamily\thepage}
\renewcommand{\headrulewidth}{0.3pt}
\usepackage{fancyvrb}
\DefineVerbatimEnvironment{layout}{Verbatim}{fontsize=\footnotesize,frame=single,framesep=2mm,rulecolor=\color{accent!40}}
\usepackage[hidelinks,unicode]{hyperref}
\newcommand{\pen}{{\sffamily\color{accent}✎}}
\newcommand{\warn}{{\sffamily\color{warncol}⚠}}
\newcommand{\ok}{{\sffamily\color{green!45!black}✔}}
\newcommand{\spark}{{\sffamily\color{warncol}⚡}}
\newcommand{\book}{}
\newcommand{\srcline}[1]{\par\noindent{\small\color{black!70}\textit{Джерело:} #1}\par\medskip}
\setlength{\parindent}{0pt}\setlength{\parskip}{3pt}
% Mixed Ukrainian/English text: Ukrainian hyphenation patterns can't break English words,
% so give the line breaker room instead of producing overfull lines.
\setlength{\emergencystretch}{3em}\tolerance=2000
\setlength{\tabcolsep}{3.5pt}
\raggedbottom


\begin{document}

\frontmatter


\begin{titlepage}
\centering\sffamily
\vspace*{40mm}
{\Huge\bfseries\color{accent} Загальний курс фізики\par}
\vspace{6mm}
{\LARGE Каркаси конспектів лекцій\par}
\vspace{12mm}
{\large Електростатика і постійний струм\\[2mm] Молекулярна фізика і термодинаміка\par}
\vfill
{\normalsize За підручником: Кучерук І.\,М., Горбачук І.\,Т., Луцик П.\,П.\\
\textit{Загальний курс фізики}, т.~1 (1999), т.~2 (2001). — К.: Техніка.\par}
\vspace{8mm}
{\small Згенеровано з вікі проєкту, \today\par}
\end{titlepage}

\chapter*{Як користуватися}
\addcontentsline{toc}{chapter}{Як користуватися}\markboth{Як користуватися}{}
Кожен розділ — це \textbf{каркас} однієї рукописної сторінки конспекту, а не готовий конспект. Він містить:
\begin{itemize}
\item \textbf{Layout} — план сторінки у вигляді рамок [A], [B], …;
\item точні терміни й означення з підручника (у «лапках»), формули з номерами рівнянь книги, сторінки;
\item \pen{} — місця, де треба намалювати схему чи сформулювати думку власними словами;
\item \warn{} — розбіжності між підручником і програмою, помилки в книзі, неперевірені місця;
\item \ok{} — формули звірено зі сканом сторінки; \spark{} — важливий або парадоксальний висновок;
\item \textbf{Self-check} — 2–4 запитання для самоперевірки.
\end{itemize}
Пояснення до структури подано англійською, терміни й цитати — українською, як у підручнику. Посилання «т.~2, с.~38» — друковані сторінки підручника.


\tableofcontents

\mainmatter

\part{Електростатика і постійний струм}

\chapter*{Програма: 23 питання}\addcontentsline{toc}{chapter}{Програма: 23 питання}\markboth{Програма: 23 питання}{Програма: 23 питання}




{\small
\begin{longtable}{>{\raggedright\arraybackslash}p{0.044\linewidth}>{\raggedright\arraybackslash}p{0.661\linewidth}>{\raggedright\arraybackslash}p{0.099\linewidth}>{\raggedright\arraybackslash}p{0.132\linewidth}}
\toprule
\textbf{№} & \textbf{Питання (as worded in the syllabus)} & \textbf{Wiki page} & \textbf{§} \\
\midrule
\endhead
1 & Електричний заряд та його властивості & 01 & 1.1 \\
2 & Закон Кулона. Диполь в неоднорідному полі & 02 & 1.2–1.3, 1.6 \\
3 & Напруженість електричного поля E & 03 & 1.4–1.5 \\
4 & Потенціал електричного поля & 04 & 1.10–1.11 \\
5 & Диполь в однорідному електричному полі & 05 & 1.6 \\
6 & Потік вектора напруженості електричного поля E. Теорема Гауса & 06 & 1.7 \\
7 & Дивергенція вектора E & 07 & 1.8 \\
8 & Зв'язок між напруженістю та потенціалом & 08 & 1.11 \\
9 & Теорема Стокса & 09 & 1.10 \\
10 & Поле зарядженої площини & 10–12 & 1.7 \\
11 & Поле двох заряджених площин (теорема Гауса) & 10–12 & 1.7 \\
12 & Поле сферичної поверхні & 10–12 & 1.7, 1.11 \\
13 & Теорема Остроградського-Гауса & 13 & 1.8 \\
14 & Провідник у зовнішньому електричному полі & 14 & 1.12–1.13 \\
15 & Конденсатори та їх ємність. Ємність плоского, циліндричного та сферичного конденсаторів & 15 & 1.14 \\
16 & З'єднання конденсаторів & 16 & 1.14 \\
17 & Енергія зарядженого провідника & 17 & 1.25–1.26 \\
18 & Енергія електростатичного поля & 18 & 1.26 \\
19 & Електричний струм. Постійний електричний струм. Закони Ома & 19 & 2.1–2.3 \\
20 & Сила струму, густина струму. Електрорушійна сила & 20 & 2.1, 2.3 \\
21 & Поняття опору. Потужність & 21 & 2.2, 2.4 \\
22 & Перший та другий закон Кірхгофа & 22 & 2.6 \\
23 & Закон Джоуля-Ленца & 23 & 2.5 \\
\bottomrule
\end{longtable}
}

\needspace{5\baselineskip}\section*{Overlaps to keep in mind}
\begin{itemize}
\item 2 and 5 share the dipole. Define it once (in 5) and refer to it from 2.
\item 6, 7, 13 form one chain: Gauss (integral) → div (differential) → Ostrogradsky–Gauss (the mathematical bridge).
\item 19, 20, 21 overlap. Proposed split: 19 = Ohm's law in every form, 20 = I, j, ЕРС, 21 = R, ρ, P.
\item 17 and 18 give the same W, read once through the charges and once through the field.
\end{itemize}

\chapter*{1. Електричний заряд та його властивості}\addcontentsline{toc}{chapter}{1. Електричний заряд та його властивості}\markboth{1. Електричний заряд та його властивості}{1. Електричний заряд та його властивості}


\srcline{\book{} Кучерук т. 2, \textbf{§ 1.1 «Електричний заряд», с. 14–21} (spreads 8–11). All formulas checked against the scan.}
This is a scaffold for the handwritten page: exact terms and formulas plus the structure. The wording on the page is yours.

\par\noindent{\color{accent!30}\rule{\linewidth}{0.4pt}}\par

\needspace{5\baselineskip}\section*{Suggested page layout (one A4 side)}

\begin{layout}
┌──────────────────────────────────────────────────────┐
│ 1. ЕЛЕКТРИЧНИЙ ЗАРЯД                                  │
│ [Box 0] What is it? (definition + 1.1)                │
├──────────────────────────┬───────────────────────────┤
│ [1] Два види (±)         │ [2] Квантування            │
│                          │     + Millikan sketch      │
├──────────────────────────┼───────────────────────────┤
│ [3] Інваріантність       │ [4] Закон збереження (1.2)│
├──────────────────────────┴───────────────────────────┤
│ [5] Неперервний розподіл: ρ, σ, τ  (1.3)–(1.5)        │
│ [footer] Носії заряду: e⁻, p⁺, іони · електризація     │
└──────────────────────────────────────────────────────┘
\end{layout}

\par\noindent{\color{accent!30}\rule{\linewidth}{0.4pt}}\par

\needspace{5\baselineskip}\section*{Box 0 — Що таке заряд (p. 15)}

\begin{itemize}
\item Book's definition: «Електричний заряд є \textbf{внутрішньою характеристикою} деяких елементарних матеріальних частинок, яка проявляється в їхній взаємодії на відстані або під час дії на них зовнішніх електромагнітних полів.»
\item Key idea to write in your own words: a charge is \textbf{both the source and the object} of the field («як джерело, так і об'єкт дії»).
\item How a charge is measured (p. 14): put two particles in turn at the same point M of a field:
\end{itemize}

\[\frac{F_1}{F_2} = \frac{q_1}{q_2} \qquad (1.1)\]

The ratio depends on neither the field nor the point M, only on the particles.
\begin{itemize}
\item Contrast with mass (p. 15): mass has one sign and only attracts; charge has two signs, so charges can \textbf{compensate} each other.
\end{itemize}

\needspace{5\baselineskip}\section*{Box 1 — Два види зарядів (p. 15)}

\begin{itemize}
\item Like charges repel and unlike charges attract (однойменні / різнойменні).
\item Sign convention: \textbf{позитивний} = скло, натерте шовком; \textbf{негативний} = ебоніт (бурштин), натертий хутром.
\item The convention is arbitrary. What matters is that there are two kinds.
\item \pen{} Your line: why are bodies usually neutral?
\end{itemize}

\needspace{5\baselineskip}\section*{Box 2 — Дискретність: закон квантування електричного заряду (pp. 16–19)}

\begin{itemize}
\item Book's term (italic): \textbf{\textit{закон квантування електричного заряду}}. Free charges consist of equal discrete portions equal to the electron's charge.
\end{itemize}

\[q = n|e|,\quad n = \pm1, \pm2, \ldots \qquad |e| = 1{,}602\cdot10^{-19}\ \text{\text{Кл}}\]

\begin{itemize}
\item \textbf{Дослід Міллікена (рис. 1.1, p. 18)}: sketch two horizontal plates, an oil droplet, and arrows for qE↑, $F_{A}$↑, mg↓, av↕.
\begin{itemize}
\item No field: the droplet falls uniformly, so mg = av + $F_{A}$, which gives \textit{a}.
\item With field E: the droplet rises uniformly, so \textbf{qE + $F_{A}$ = mg + av′}, which gives \textit{q}.
\item X-rays change the charge: Δq = q′ − q = n|e| \textbf{always}.
\end{itemize}
\item The quantization itself is a \textbf{постулат}: theory doesn't explain it (p. 19).
\item Side facts, optional: e/m = 1,7588·10¹¹ Кл/кг (Thomson 1897); $m_{e}$ = 9,109·10⁻³¹ кг; quarks ±e/3, 2e/3, never found free.
\end{itemize}

\needspace{5\baselineskip}\section*{Box 3 — Інваріантність заряду (p. 20)}

\begin{itemize}
\item The value does \textbf{not} depend on the inertial frame or the speed of the carrier.
\item Evidence: atoms are neutral. Electrons in He move about 2× faster than in H₂, yet both are neutral, so the charge doesn't depend on speed. Checked up to v ≈ 0,5c.
\item \pen{} Contrast it yourself: mass \textit{does} change with speed (relativity), charge doesn't.
\end{itemize}

\needspace{5\baselineskip}\section*{Box 4 — Закон збереження заряду (pp. 20–21)}

\begin{itemize}
\item \textbf{\textit{Замкнена система}} is one where no charged particles cross its boundary.
\item «Повний електричний заряд замкненої системи є величиною постійною і дорівнює алгебричній сумі її позитивних та негативних зарядів»:
\end{itemize}

\[q = q_1 + q_2 + \ldots + q_n = \sum_{i=1}^{n} q_i = \text{const} \qquad (1.2)\]

\begin{itemize}
\item Consequence: charges are born and die \textbf{only in ± pairs} (e⁻e⁺ pair creation from γ, annihilation).
\item Electrification by friction only \textbf{redistributes} charge; the total doesn't change (p. 16). This is a good example to draw.
\end{itemize}

\needspace{5\baselineskip}\section*{Box 5 — Неперервний розподіл: густини заряду (p. 21)}

\begin{itemize}
\item Why: 1 А ≈ 6·10¹⁸ elementary charges per second, so the discreteness is invisible and we treat charge as continuous.
\end{itemize}

{\small
\begin{longtable}{>{\raggedright\arraybackslash}p{0.306\linewidth}>{\raggedright\arraybackslash}p{0.594\linewidth}>{\raggedright\arraybackslash}p{0.052\linewidth}}
\toprule
\textbf{Term (book's italic)} & \textbf{Formula} & \textbf{Unit} \\
\midrule
\endhead
\textbf{\textit{об'ємна густина заряду}} & $\rho = \lim\limits_{\Delta V\to0}\dfrac{\Delta q}{\Delta V} = \dfrac{dq}{dV}$ (1.3) & Кл/м³ \\
\textbf{\textit{поверхнева густина заряду}} & $\sigma = \lim\limits_{\Delta S\to0}\dfrac{\Delta q}{\Delta S} = \dfrac{dq}{dS}$ (1.4) & Кл/м² \\
\textbf{\textit{лінійна густина}} & $\tau = \lim\limits_{\Delta l\to0}\dfrac{\Delta q}{\Delta l} = \dfrac{dq}{dl}$ (1.5) & Кл/м \\
\bottomrule
\end{longtable}
}

\begin{itemize}
\item Caveat (book): ΔV, ΔS, Δl are small, but each still holds many elementary charges.
\item You'll reuse σ in topics 10–12 (charged plane, two planes, sphere) and τ for a charged thread.
\end{itemize}

\needspace{5\baselineskip}\section*{Footer — носії заряду (p. 16)}

\begin{itemize}
\item Stable carriers: \textbf{електрон} (−e), \textbf{протон} (+e, equal in magnitude to the electron's charge to a very high relative precision, p. 19). Ions form when an atom loses or gains electrons.
\end{itemize}

\par\noindent{\color{accent!30}\rule{\linewidth}{0.4pt}}\par

\needspace{5\baselineskip}\section*{Skip, digressions with no bearing on the topic}

History (Gilbert, Franklin, Dufay, Stoney), the classical electron radius r₀ = e²/(4πε₀m₀c²) = 2,82·10⁻¹⁵ м, spin, the electron's "infinite energy" paradox, and the neutron's lifetime. Include them only if your lecturer cares about history.

\needspace{5\baselineskip}\section*{Self-check before you write}

\begin{enumerate}
\item Why does an atom being neutral prove invariance?
\item What does Millikan's experiment prove, and what does it measure?
\item Why does friction electrify dielectrics but not metals? (p. 16)
\item Which property is a \textit{postulate}, and which are experimental facts?
\end{enumerate}

\chapter*{2. Закон Кулона. Диполь в неоднорідному полі}\addcontentsline{toc}{chapter}{2. Закон Кулона. Диполь в неоднорідному полі}\markboth{2. Закон Кулона. Диполь в неоднорідному полі}{2. Закон Кулона. Диполь в неоднорідному полі}


\srcline{\book{} т. 2, \textbf{§ 1.2} (с. 21–26), \textbf{§ 1.3} (с. 26), \textbf{§ 1.6} (с. 32–35). \ok{} Formulas checked against the scan.}

\needspace{5\baselineskip}\section*{Layout}
\begin{layout}
┌ [A] Закон Кулона: формулювання + (1.6)/(1.7) + рис. 1.2 крутильні ваги ┐
├ [B] Векторна форма (1.8)  │ [C] k, ε₀, одиниці                         ┤
├ [D] Принцип суперпозиції (1.10) + рис. 1.4                              ┤
├ [E] Середовище: ε (§1.3)  │ [F] Диполь у неоднорідному полі (1.26)       ┤
└ footer: межі застосування закону                                         ┘
\end{layout}

\needspace{5\baselineskip}\section*{[A] Закон Кулона (p. 22)}
\begin{itemize}
\item The book's wording (worth learning by heart): «сила взаємодії F двох \textbf{нерухомих точкових зарядів у вакуумі} прямо пропорційна добуткові кількостей електрики цих зарядів q₁ і q₂, обернено пропорційна квадратові відстані r між ними і напрямлена вздовж прямої, яка сполучає ці заряди»
\item $F = k\dfrac{q_1q_2}{r^2}$ (1.6) → in SI: $F = \dfrac{1}{4\pi\varepsilon_0}\dfrac{q_1q_2}{r^2}$ (1.7)
\item \textbf{точковий заряд}: the size of the bodies ≪ the distance between them (p. 24).
\item Sign convention: repulsion F > 0, attraction F < 0.
\item \pen{} Sketch Coulomb's torsion balance (рис. 1.2): thread, beam, balls a and d.
\end{itemize}

\needspace{5\baselineskip}\section*{[B] Векторна форма (p. 24)}
\begin{itemize}
\item $\vec F_{12} = \dfrac{1}{4\pi\varepsilon_0}\dfrac{q_1q_2}{r_{12}^3}\vec r_{12}$ (1.8), with $\vec F_{12} = -\vec F_{21}$ (1.9)
\item The forces are \textbf{центральні} (central): they act along the line joining the charges.
\end{itemize}

\needspace{5\baselineskip}\section*{[C] Коефіцієнт і одиниці (pp. 23–24)}
\begin{itemize}
\item $k = \dfrac{1}{4\pi\varepsilon_0} = 9\cdot10^9\ \dfrac{\text{\text{Н}·\text{м}}^2}{\text{\text{Кл}}^2}$; \textbf{електрична стала} $\varepsilon_0 = 8{,}85\cdot10^{-12}$ Ф/м.
\item 1 Кл = 1 А·1 с; the ampere is the base unit, so the coulomb is a derived one.
\item "Раціоналізована форма" puts the 4π in k so that it drops out of the formulas used most often.
\item Optional: in СГСЕ k = 1, and 1 Кл = 3·10⁹ СГСЕ units.
\end{itemize}

\needspace{5\baselineskip}\section*{[D] Принцип суперпозиції (p. 25)}
\begin{itemize}
\item The force between two charges \textbf{does not depend on other charges being present}; the net force is the vector sum:
$\vec F = \sum_i \vec F_i = \sum_i \dfrac{1}{4\pi\varepsilon_0}\dfrac{q_iq_0}{r_{i0}^3}\vec r_{i0}$ (1.10)
\item For extended bodies, divide them mentally into point charges.
\end{itemize}

\needspace{5\baselineskip}\section*{[E] Вплив середовища (§ 1.3, p. 26)}
\begin{itemize}
\item \textbf{діелектрична проникність} $\varepsilon = F_0/F$, how many times weaker the force is in a medium than in vacuum (dimensionless).
\item $F = \dfrac{1}{4\pi\varepsilon_0}\dfrac{q_1q_2}{\varepsilon r^2}$
\item Why it happens: the medium undergoes \textbf{поляризація} (polarization) and partly neutralizes the charges.
\end{itemize}

\needspace{5\baselineskip}\section*{[F] Диполь у неоднорідному полі (p. 35)}
\begin{itemize}
\item In a non-uniform field the two forces differ, so a net force appears as well as the torque:
$F = p\,\dfrac{\partial E}{\partial x}$ (1.26), dipole along a field line.
\item Direction: \textbf{towards increasing field strength}; the dipole is pulled into the stronger field.
\item \pen{} Example to write: why a charged comb attracts scraps of paper (induced dipoles).
\item Optional: Example 1.1, two dipoles on one line, $F = -\dfrac{6p_1p_2}{4\pi\varepsilon_0\varepsilon r^4}$.
\item The definition of a dipole itself is in topic 5 (put it there or here, but only once).
\end{itemize}

\needspace{5\baselineskip}\section*{Footer: межі застосування (p. 26)}
\begin{itemize}
\item Verified from about 10⁻¹⁷ м (electron scattering) up to about 10⁷ м (satellites).
\end{itemize}

\needspace{5\baselineskip}\section*{Self-check}
\begin{enumerate}
\item Why is the law stated for \textit{point} charges \textit{at rest}?
\item What does superposition add that Coulomb's law alone doesn't give?
\item Why does a dipole in a uniform field rotate but not move?
\end{enumerate}

\chapter*{3. Напруженість електричного поля E}\addcontentsline{toc}{chapter}{3. Напруженість електричного поля E}\markboth{3. Напруженість електричного поля E}{3. Напруженість електричного поля E}


\srcline{\book{} т. 2, \textbf{§ 1.4} (с. 27), \textbf{§ 1.5} (с. 28–31). \ok{} Formulas checked against the scan.}

\needspace{5\baselineskip}\section*{Layout}
\begin{layout}
┌ [A] Що таке поле (§1.4): близькодія vs далекодія ┐
├ [B] Означення E (1.11) + пробний заряд             ┤
├ [C] Поле точкового заряду (1.12) │ [D] Суперпозиція (1.13)–(1.14) ┤
├ [E] Неперервні заряди (1.15)–(1.18)                 ┤
└ [F] Силові лінії: рис. 1.6–1.8; однорідне поле       ┘
\end{layout}

\needspace{5\baselineskip}\section*{[A] Електричне поле (p. 27)}
\begin{itemize}
\item \textbf{Електричне поле}: «матеріальна складова електромагнітного поля, яке діє на заряд, зумовлене зарядами або змінним у часі магнітним полем».
\item Main sign of a field: \textbf{a force acts on any charge placed in it}.
\item Two concepts: \textbf{далекодія} (action at a distance: instantaneous, no carrier) was rejected; \textbf{близькодія} (local action: through the field, at finite speed) is accepted.
\item In electrostatics both give the same results; they differ only for moving charges.
\end{itemize}

\needspace{5\baselineskip}\section*{[B] Означення напруженості (p. 28)}
\begin{itemize}
\item \textbf{Пробний заряд} (test charge): a point charge, positive, small enough not to distort the field.
\item «Напруженість E у даній точці електростатичного поля чисельно дорівнює силі, з якою поле діє на \textbf{одиничний позитивний пробний заряд}, вміщений у цю точку»
\item $\vec E = \vec F / q_0$ (1.11), a \textbf{вектор} and the \textbf{силова характеристика поля} (force characteristic).
\item Units: Н/Кл = В/м.
\item E is \textbf{локальна} (local): each point has its own value.
\end{itemize}

\needspace{5\baselineskip}\section*{[C] Поле точкового заряду (p. 28)}
\begin{itemize}
\item $\vec E = \dfrac{1}{4\pi\varepsilon_0}\dfrac{q}{r^3}\vec r$ (1.12), so $E = \dfrac{q}{4\pi\varepsilon_0 r^2}$
\item The field is \textbf{центрально-симетричне} (centrally symmetric).
\item The two-step view of interaction (p. 29): q₁ creates the field E₁₂, and the field acts on q₂ with $\vec F_{12} = q_2\vec E_{12}$.
\end{itemize}

\needspace{5\baselineskip}\section*{[D] Принцип суперпозиції полів (pp. 29–30)}
\begin{itemize}
\item $\vec E = \sum_{i=1}^n \vec E_i$ (1.13), $\vec E = \sum_i \dfrac{1}{4\pi\varepsilon_0}\dfrac{q_i}{r_i^3}\vec r_i$ (1.14)
\item The \textbf{польове формулювання} (field formulation): the field of a system equals the sum of the fields of each charge «за відсутності всіх інших».
\item «Основне завдання електростатики»: find \textbf{E} at every point from a given distribution of charges (p. 29).
\end{itemize}

\needspace{5\baselineskip}\section*{[E] Неперервний розподіл (pp. 30–31)}
\begin{itemize}
\item $dq = \rho\,dV$ → $\vec E = \dfrac{1}{4\pi\varepsilon_0}\int_V \dfrac{\rho\,\vec r}{r^3}dV$ (1.15); similarly with σdS and τdl (1.18).
\item Links to topic 1: ρ, σ, τ.
\end{itemize}

\needspace{5\baselineskip}\section*{[F] Силові лінії (p. 31)}
\begin{itemize}
\item \textbf{Силова лінія} is a line whose tangent at every point points along E.
\item They start on + charges and end on − charges (or go to infinity); they \textbf{never cross}.
\item Their density (lines per unit area) ∝ E.
\item \textbf{Однорідне поле} (uniform field): E is the same everywhere, so the lines are parallel and evenly spaced.
\item \pen{} Sketch three pictures: an isolated +, an isolated −, and a +/− pair (рис. 1.6–1.8).
\end{itemize}

\needspace{5\baselineskip}\section*{Self-check}
\begin{enumerate}
\item Why must the test charge be small?
\item Why can't field lines cross?
\item Which is primary in close-range action: the force or the field?
\end{enumerate}

\chapter*{4. Потенціал електричного поля}\addcontentsline{toc}{chapter}{4. Потенціал електричного поля}\markboth{4. Потенціал електричного поля}{4. Потенціал електричного поля}


\srcline{\book{} т. 2, \textbf{§ 1.10} (с. 44–45, робота), \textbf{§ 1.11} (с. 48–56). \ok{} Formulas checked against the scan.}

\needspace{5\baselineskip}\section*{Layout}
\begin{layout}
┌ [A] Робота поля не залежить від шляху (1.44)–(1.46) ┐
├ [B] Означення φ (1.54)  │ [C] φ точкового заряду (1.55)  ┤
├ [D] Суперпозиція (1.56) │ [E] Різниця потенціалів, напруга (1.59), еВ ┤
└ [F] Еквіпотенціальні поверхні: рис., E ⟂ поверхні     ┘
\end{layout}

\needspace{5\baselineskip}\section*{[A] Робота сил поля (pp. 44–45)}
\begin{itemize}
\item Moving q₀ in the field of q from 1 to 2: $A_{12} = \dfrac{qq_0}{4\pi\varepsilon_0}\left(\dfrac1{r_1}-\dfrac1{r_2}\right)$ (1.45)
\item The work \textbf{does not depend on the path}, so the field is \textbf{потенціальне (консервативне)} (conservative).
\item Around a closed loop the work is 0: $\oint_L \vec E\,d\vec l = 0$ (1.46), the \textbf{циркуляція вектора E} (circulation). More on this in topic 9.
\end{itemize}

\needspace{5\baselineskip}\section*{[B] Означення потенціалу (pp. 48–49)}
\begin{itemize}
\item $A_{12} = W_1 - W_2$: work equals the \textbf{decrease} in potential energy (1.53).
\item $\varphi_1 = \dfrac{A_{1\infty}}{q_0} = \int_1^\infty E\,dl\cos\alpha$ (1.54)
\item The book's definition: «фізична величина, що чисельно дорівнює роботі, яку виконують сили електричного поля при переміщенні \textbf{пробного одиничного заряду} із заданої точки \textbf{в нескінченність}»
\item An equivalent version: the \textbf{ratio of the potential energy of the charge to the charge}, $\varphi = W/q_0$.
\item φ is a \textbf{скалярна} (scalar) and \textbf{енергетична характеристика} (energy characteristic). 1 В = 1 Дж/1 Кл.
\item Zero point: φ(∞) = 0 in theory and φ(Землі) = 0 in practice.
\end{itemize}

\needspace{5\baselineskip}\section*{[C] Потенціал точкового заряду (p. 50)}
\begin{itemize}
\item $\varphi = \dfrac{1}{4\pi\varepsilon_0}\dfrac{q}{r}$ (1.55), positive for +q and negative for −q.
\end{itemize}

\needspace{5\baselineskip}\section*{[D] Суперпозиція (p. 50)}
\begin{itemize}
\item $\varphi = \dfrac{1}{4\pi\varepsilon_0}\sum_i \dfrac{q_i}{r_i}$ (1.56), an \textbf{algebraic} sum, not a vector one. This is the main practical advantage of φ over E.
\item Continuous charges: $\varphi = \dfrac{1}{4\pi\varepsilon_0}\int \dfrac{\rho\,dV}{r}$ (1.57).
\end{itemize}

\needspace{5\baselineskip}\section*{[E] Різниця потенціалів = напруга (pp. 50–51)}
\begin{itemize}
\item φ is defined only up to an additive constant, so \textbf{Δφ} is the physically meaningful quantity:
$\varphi_1 - \varphi_2 = \dfrac{A_{12}}{q_0}$ (1.59), written \textbf{U}, unit: вольт.
\item \textbf{Електрон-вольт}: 1 еВ = 1,6·10⁻¹⁹ Дж.
\end{itemize}

\needspace{5\baselineskip}\section*{[F] Еквіпотенціальні поверхні (p. 51)}
\begin{itemize}
\item A surface where φ = const. Moving a charge along it does zero work, so \textbf{E ⟂ the surface}.
\item A conductor's surface is equipotential.
\item \pen{} Sketch: a point charge, with concentric circles (φ) and radial arrows (E).
\end{itemize}

\needspace{5\baselineskip}\section*{Self-check}
\begin{enumerate}
\item Why can we define φ at all? (Hint: use [A].)
\item Why is φ simpler to compute than E?
\item Why is Δφ "more physical" than φ?
\end{enumerate}

\chapter*{5. Диполь в однорідному електричному полі}\addcontentsline{toc}{chapter}{5. Диполь в однорідному електричному полі}\markboth{5. Диполь в однорідному електричному полі}{5. Диполь в однорідному електричному полі}


\srcline{\book{} т. 2, \textbf{§ 1.6} (с. 32–35); potential and energy in \textbf{§ 1.11} (с. 53–54). \ok{} Formulas checked against the scan.}

\needspace{5\baselineskip}\section*{Layout}
\begin{layout}
┌ [A] Означення диполя, p (1.19) + рис. 1.9  ┐
├ [B] Поле диполя: на осі (1.20), на перпендикулярі (1.21) ┤
├ [C] Момент сил у однорідному полі (1.23)–(1.25) + рис. 1.13 ┤
└ [D] Енергія, рівновага  │ [E] Потенціал диполя (1.66)       ┘
\end{layout}

\needspace{5\baselineskip}\section*{[A] Електричний диполь (p. 32)}
\begin{itemize}
\item «сукупність двох однакових за абсолютним значенням і протилежних за знаком точкових зарядів +q і −q, розміщених на деякій відстані l»
\item \textbf{плече диполя} l (the arm). \textbf{Жорсткий} (rigid): l is constant. \textbf{Точковий} (point dipole): l ≪ r. Example: polar molecules.
\item \textbf{електричний дипольний момент}: $\vec p = q\vec l$ (1.19), pointing \textbf{from − to +}; unit Кл·м.
\end{itemize}

\needspace{5\baselineskip}\section*{[B] Поле точкового диполя (pp. 32–34)}
\begin{itemize}
\item On the axis: $\vec E = \dfrac{1}{4\pi\varepsilon_0\varepsilon}\dfrac{2\vec p}{r^3}$ (1.20)
\item On the perpendicular bisector: $\vec E = -\dfrac{1}{4\pi\varepsilon_0\varepsilon}\dfrac{\vec p}{r^3}$ (1.21), antiparallel to p.
\item In general (1.22): falls off as \textbf{1/r³}, compared with 1/r² for a point charge.
\end{itemize}

\needspace{5\baselineskip}\section*{[C] Дія однорідного поля (p. 34)}
\begin{itemize}
\item F₊ = −F₋, so the resultant is 0, but they form a \textbf{пара сил} (couple):
$M = pE\sin\alpha$ (1.24), $\vec M = [\vec p\times\vec E]$ (1.25)
\item The couple turns \textbf{p} towards \textbf{E}.
\item \pen{} Sketch рис. 1.13: a dipole at angle α to E, with two forces qE.
\end{itemize}

\needspace{5\baselineskip}\section*{[D] Енергія і рівновага (pp. 34, 54)}
\begin{itemize}
\item $W = -pE\cos\alpha = -\vec p\cdot\vec E$ (1.67)
\item p ∥ E means \textbf{стійка рівновага} (stable, W is minimal); p antiparallel to E means \textbf{нестійка} (unstable).
\end{itemize}

\needspace{5\baselineskip}\section*{[E] Потенціал диполя (p. 53)}
\begin{itemize}
\item $\varphi = \dfrac{1}{4\pi\varepsilon_0}\dfrac{\vec p\cdot\vec r}{r^3} = \dfrac{p\cos\theta}{4\pi\varepsilon_0 r^2}$ (1.66), falls off as 1/r².
\end{itemize}

\needspace{5\baselineskip}\section*{Contrast with topic 2}
{\small
\begin{longtable}{>{\raggedright\arraybackslash}p{0.190\linewidth}>{\raggedright\arraybackslash}p{0.554\linewidth}>{\raggedright\arraybackslash}p{0.208\linewidth}}
\toprule
\textbf{Поле} & \textbf{Сила} & \textbf{Момент} \\
\midrule
\endhead
однорідне & 0 & pE sin α \\
неоднорідне & p ∂E/∂x, into the stronger field & also present \\
\bottomrule
\end{longtable}
}

\needspace{5\baselineskip}\section*{Self-check}
\begin{enumerate}
\item Why is the net force zero in a uniform field?
\item Which equilibrium is stable, and why? (Use W.)
\item Why does a dipole's field fall off faster than a charge's?
\end{enumerate}

\chapter*{6. Потік вектора напруженості. Теорема Гаусса}\addcontentsline{toc}{chapter}{6. Потік вектора напруженості. Теорема Гаусса}\markboth{6. Потік вектора напруженості. Теорема Гаусса}{6. Потік вектора напруженості. Теорема Гаусса}


\srcline{\book{} т. 2, \textbf{§ 1.7} (с. 36–38). The applications are in topics 10–12. \ok{} Formulas checked against the scan.}
Spelling: the book writes \textbf{Гаусса} (double с); the syllabus writes «Гауса».

\needspace{5\baselineskip}\section*{Layout}
\begin{layout}
┌ [A] Потік: означення (1.27) + рис. 1.16 (площадка, нормаль, кут α) ┐
├ [B] Ідея доведення: тілесний кут (1.29), рис. 1.17 а/б             ┤
├ [C] Теорема Гаусса: формулювання + (1.31), (1.33)                  ┤
└ [D] Фізичний зміст: джерела/стоки; Гаусс = Кулон в інтегральній формі ┘
\end{layout}

\needspace{5\baselineskip}\section*{[A] Потік вектора E (p. 36)}
\begin{itemize}
\item A uniform field through a flat patch: $\Delta\Phi_E = E\,\Delta S\cos\alpha = E_n\Delta S$ (1.27)
\item In general: $\Phi_E = \int_S E_n\,dS$
\item The sign comes from cos α. The \textbf{outward} normal is taken as positive.
\item Picture: the number of field lines crossing the surface.
\end{itemize}

\needspace{5\baselineskip}\section*{[B] Ідея доведення (pp. 36–38)}
\begin{itemize}
\item For a point charge, $d\Phi_E = \dfrac{q}{4\pi\varepsilon_0}d\omega$ (1.29), where dω is the \textbf{тілесний кут} (solid angle).
\item Charge \textbf{inside}: $\oint d\omega = 4\pi$, so $\Phi_E = q/\varepsilon_0$ (1.30).
\item Charge \textbf{outside}: every cone crosses the surface an even number of times, so $\Phi_E = 0$ (1.32).
\item \pen{} Draw рис. 1.17 а/б: a cone from q cutting a closed surface once vs twice.
\end{itemize}

\needspace{5\baselineskip}\section*{[C] Теорема Гаусса (p. 38)}
\begin{itemize}
\item «потік вектора напруженості електростатичного поля у вакуумі Φ\_E крізь \textbf{довільну замкнену поверхню} S зсередини назовні дорівнює \textbf{алгебричній сумі} тих точкових зарядів $q_{i}$, які охоплюються поверхнею S, поділеній на електричну сталу ε₀»
\item $\Phi_E = \oint_S E_n\,dS = \dfrac{1}{\varepsilon_0}\sum_i q_i$ (1.31)
\item Continuous charges: $\Phi_E = \dfrac{1}{\varepsilon_0}\left(\int_V\rho\,dV + \int_S\sigma\,dS\right)$ (1.33)
\end{itemize}

\needspace{5\baselineskip}\section*{[D] Фізичний зміст (p. 38)}
\begin{itemize}
\item + charges are \textbf{джерела} (sources) of the field and − charges are \textbf{стоки} (sinks).
\item The theorem holds \textit{because} E ∝ 1/r². So «теорема Гаусса є \textbf{інтегральною формою запису закону Кулона}».
\item Why it matters: it turns hard integrals (1.15) into a short symmetry argument (topics 10–12), and it leads to the differential form (topic 7).
\end{itemize}

\needspace{5\baselineskip}\section*{Self-check}
\begin{enumerate}
\item Why doesn't the flux depend on the shape of the surface?
\item Why does a charge outside give zero flux?
\item Why would the theorem fail if the force went as 1/r³?
\end{enumerate}

\chapter*{7. Дивергенція вектора E}\addcontentsline{toc}{chapter}{7. Дивергенція вектора E}\markboth{7. Дивергенція вектора E}{7. Дивергенція вектора E}


\srcline{\book{} т. 2, \textbf{§ 1.8 «Теорема Гаусса в диференціальній формі»} (с. 41–43). \ok{} Formulas checked against the scan.}
This overlaps topic 13 (Остроградський–Гаусс). Suggested split: \textbf{this page = div and its meaning}, \textbf{topic 13 = the integral↔differential bridge}.

\needspace{5\baselineskip}\section*{Layout}
\begin{layout}
┌ [A] Навіщо: інтегральна vs диференціальна форма ┐
├ [B] Виведення: паралелепіпед dx dy dz, рис. 1.22 ┤
├ [C] Означення div (через границю) + декартова форма ┤
└ [D] div E = ρ/ε₀ (1.42) + фізичний зміст        ┘
\end{layout}

\needspace{5\baselineskip}\section*{[A] Мотивація (p. 41)}
\begin{itemize}
\item The integral form links \textbf{different} points: E on the surface and q inside the volume.
\item The differential form links quantities \textbf{at one point}, which fits close-range action.
\end{itemize}

\needspace{5\baselineskip}\section*{[B] Виведення (pp. 41–42)}
\begin{itemize}
\item A tiny box dx·dy·dz around charge ρ dV. The flux through the pair of x-faces is $\dfrac{\partial E_x}{\partial x}dx\,dy\,dz$; the same goes for y and z.
\item By Gauss: $\left(\dfrac{\partial E_x}{\partial x}+\dfrac{\partial E_y}{\partial y}+\dfrac{\partial E_z}{\partial z}\right)dV = \dfrac{\rho\,dV}{\varepsilon_0}$ (1.38)
\item \pen{} Draw the box with $E_{x}$ entering on the left and $E_{x}$ + (∂$E_{x}$/∂x)dx leaving on the right.
\end{itemize}

\needspace{5\baselineskip}\section*{[C] Означення дивергенції (p. 42)}
\begin{itemize}
\item \textbf{розходження, або дивергенція} of a vector A at a point is the limit of flux per unit volume:
$\operatorname{div}\vec A = \lim\limits_{V\to0}\dfrac{\oint_S A_n\,dS}{V}$, which doesn't depend on the shape of S.
\item Cartesian form: $\operatorname{div}\vec E = \dfrac{\partial E_x}{\partial x}+\dfrac{\partial E_y}{\partial y}+\dfrac{\partial E_z}{\partial z}$ (1.43)
\end{itemize}

\needspace{5\baselineskip}\section*{[D] Результат (pp. 42–43)}
\begin{itemize}
\item $\operatorname{div}\vec E = \dfrac{\rho}{\varepsilon_0}$ (1.42), Gauss's theorem in differential form = \textbf{диференціальна форма закону Кулона}.
\item Meaning: div measures the \textbf{потужність джерел або стоків} (strength of sources or sinks).
\begin{itemize}
\item div E > 0: lines start here (+ charge).
\item div E < 0: lines end here (− charge).
\item Where ρ = 0: div E = 0, so the lines pass through.
\end{itemize}
\end{itemize}

\needspace{5\baselineskip}\section*{Self-check}
\begin{enumerate}
\item What does div E = 0 mean at a point in empty space?
\item Why is (1.42) "local" and (1.31) not?
\end{enumerate}

\chapter*{8. Зв'язок між напруженістю та потенціалом}\addcontentsline{toc}{chapter}{8. Зв'язок між напруженістю та потенціалом}\markboth{8. Зв'язок між напруженістю та потенціалом}{8. Зв'язок між напруженістю та потенціалом}


\srcline{\book{} т. 2, \textbf{§ 1.11} (с. 51–53). \ok{} Formulas checked against the scan.}

\needspace{5\baselineskip}\section*{Layout}
\begin{layout}
┌ [A] Інтегральний зв'язок: φ₁ − φ₂ = ∫E dl (1.60)        ┐
├ [B] Диференціальний: E_x = −∂φ/∂x → E = −grad φ (1.62)   ┤
├ [C] Однорідне поле: E = U/d   │ [D] Градієнт: що це        ┤
└ [E] Рівняння Пуассона (1.63) і Лапласа (1.64)            ┘
\end{layout}

\needspace{5\baselineskip}\section*{[A] φ із E (p. 51)}
\begin{itemize}
\item $\varphi_1 - \varphi_2 = \int_1^2 \vec E\,d\vec l$ (1.60), which gives φ when E is known.
\begin{itemize}
\item \warn{} \textbf{Misprint in the book} (p. 51): (1.60) is printed with a minus, φ₁ − φ₂ = −∫E dl. That contradicts its own (1.53) and (1.59), A₁₂ = q₀∫E dl = q₀(φ₁ − φ₂). The form without the minus is correct: going along E, the potential falls.
\end{itemize}
\end{itemize}

\needspace{5\baselineskip}\section*{[B] E із φ (pp. 51–52)}
\begin{itemize}
\item Over a small dx: $dA = q_0E_x\,dx = -q_0\,d\varphi$, so $E_x = -\dfrac{\partial\varphi}{\partial x}$, and likewise $E_y$, $E_z$.
\item $\vec E = -\left(\dfrac{\partial\varphi}{\partial x}\vec i + \dfrac{\partial\varphi}{\partial y}\vec j + \dfrac{\partial\varphi}{\partial z}\vec k\right) = -\operatorname{grad}\varphi$ (1.61)–(1.62)
\item \textbf{Мінус}: E points the way φ \textbf{decreases}.
\end{itemize}

\needspace{5\baselineskip}\section*{[C] Однорідне поле (p. 52)}
\begin{itemize}
\item $E = \dfrac{\varphi_1-\varphi_2}{d} = \dfrac{U}{d}$, which is why E is measured in \textbf{В/м}.
\end{itemize}

\needspace{5\baselineskip}\section*{[D] Градієнт (p. 52)}
\begin{itemize}
\item \textbf{grad φ} is the vector pointing in the direction of the \textbf{fastest increase} of φ; its length is the derivative of φ along that direction.
\item So E ⟂ the equipotentials (topic 4[F]), pointing "downhill".
\item \pen{} A picture that helps: φ as the height of a hill, E as the direction a ball rolls.
\end{itemize}

\needspace{5\baselineskip}\section*{[E] Рівняння Пуассона і Лапласа (pp. 52–53)}
\begin{itemize}
\item Put E = −grad φ into div E = ρ/ε₀:
$\nabla^2\varphi = -\dfrac{\rho}{\varepsilon_0}$ (1.63), the \textbf{рівняння Пуассона}
\item Where ρ = 0: $\nabla^2\varphi = 0$ (1.64), the \textbf{рівняння Лапласа}
\item $\nabla^2 = \dfrac{\partial^2}{\partial x^2}+\dfrac{\partial^2}{\partial y^2}+\dfrac{\partial^2}{\partial z^2}$, the \textbf{лапласіан}
\item And rot grad φ ≡ 0, so the second equation (rot E = 0) holds automatically (links to topic 9).
\end{itemize}

\needspace{5\baselineskip}\section*{Self-check}
\begin{enumerate}
\item Where is the field stronger: where equipotentials are dense or sparse?
\item Can E = 0 where φ ≠ 0? Give an example. (Inside a conductor.)
\end{enumerate}

\chapter*{9. Теорема Стокса}\addcontentsline{toc}{chapter}{9. Теорема Стокса}\markboth{9. Теорема Стокса}{9. Теорема Стокса}


\srcline{\book{} т. 2, \textbf{§ 1.10} (с. 45–48). The book doesn't give Stokes a section of its own: it appears inside «Робота сил електричного поля. Потенціальний характер електростатичного поля». \ok{} Formulas checked against the scan.}

\needspace{5\baselineskip}\section*{Layout}
\begin{layout}
┌ [A] Циркуляція E = 0 (1.46): електростатичне поле потенціальне ┐
├ [B] Ротор: означення (1.47) + рис. 1.25, 1.26                  ┤
├ [C] Декартова форма rot A (1.49)                                ┤
├ [D] Теорема Стокса: ∮A dl = ∫rot A dS                           ┤
└ [E] rot E = 0 (1.51); система div E = ρ/ε₀, rot E = 0 (1.52)    ┘
\end{layout}

\needspace{5\baselineskip}\section*{[A] Циркуляція (p. 45)}
\begin{itemize}
\item \textbf{циркуляція вектора E}: $\oint_L \vec E\,d\vec l$ around a closed loop L.
\item Electrostatics: $\oint_L \vec E\,d\vec l = 0$ (1.46), the integral form of \textbf{потенціальність} (the field being conservative).
\item Consequence: electrostatic field lines are \textbf{not closed}; they start and end on charges.
\end{itemize}

\needspace{5\baselineskip}\section*{[B] Ротор (p. 46)}
\begin{itemize}
\item From the Latin \textit{roto}, "I turn". It measures the «завихрення» (curl, or whirl) of a vector field.
\item The projection onto the direction n: $\operatorname{rot}_n\vec A = \lim\limits_{\Delta S\to0}\dfrac{\oint_L \vec A\,d\vec l}{\Delta S}$ (1.47), with the loop direction set by the \textbf{правий гвинт} (right-hand screw) rule.
\end{itemize}

\needspace{5\baselineskip}\section*{[C] Декартова форма (p. 47)}
$\operatorname{rot}\vec A = \left(\dfrac{\partial A_z}{\partial y}-\dfrac{\partial A_y}{\partial z}\right)\vec i + \left(\dfrac{\partial A_x}{\partial z}-\dfrac{\partial A_z}{\partial x}\right)\vec j + \left(\dfrac{\partial A_y}{\partial x}-\dfrac{\partial A_x}{\partial y}\right)\vec k$ (1.49)
\begin{itemize}
\item \pen{} A memory aid: a 3×3 determinant with rows (i j k), (∂x ∂y ∂z), ($A_{x}$ $A_{y}$ $A_{z}$).
\end{itemize}

\needspace{5\baselineskip}\section*{[D] Теорема Стокса (p. 47)}
\begin{itemize}
\item $\oint_L \vec A\,d\vec l = \int_S \operatorname{rot}\vec A\,d\vec S$, where S is \textbf{any} surface bounded by the loop L.
\item Words: the circulation around a loop = the flux of the curl through a surface spanning that loop.
\item \pen{} Sketch: a loop L with a "soap film" S on it and normals n.
\end{itemize}

\needspace{5\baselineskip}\section*{[E] Застосування до E (pp. 47–48)}
\begin{itemize}
\item $\int_S \operatorname{rot}\vec E\,d\vec S = \oint \vec E\,d\vec l = 0$ (1.50) for any S, so $\operatorname{rot}\vec E = 0$ (1.51).
\item The two \textbf{basic differential equations of electrostatics} (1.52):
$\operatorname{div}\vec E = \rho/\varepsilon_0,\qquad \operatorname{rot}\vec E = 0$
\item These replace Coulomb's law plus superposition, and they are a special case of \textbf{Maxwell's equations}.
\end{itemize}

\needspace{5\baselineskip}\section*{Pairing table (good for the page footer)}
{\small
\begin{longtable}{>{\raggedright\arraybackslash}p{0.091\linewidth}>{\raggedright\arraybackslash}p{0.205\linewidth}>{\raggedright\arraybackslash}p{0.242\linewidth}>{\raggedright\arraybackslash}p{0.398\linewidth}}
\toprule
\textbf{} & \textbf{інтегральна форма} & \textbf{диференціальна форма} & \textbf{теорема-міст} \\
\midrule
\endhead
джерела & ∮E dS = q/ε₀ & div E = ρ/ε₀ & Остроградського–Гаусса (topic 13) \\
вихори & ∮E dl = 0 & rot E = 0 & Стокса (this topic) \\
\bottomrule
\end{longtable}
}

\needspace{5\baselineskip}\section*{Self-check}
\begin{enumerate}
\item Why do closed field lines contradict rot E = 0?
\item What is the difference between "flux" and "circulation"?
\end{enumerate}

\chapter*{10–12. Поле зарядженої площини · двох площин · сферичної поверхні}\addcontentsline{toc}{chapter}{10–12. Поле зарядженої площини · двох площин · сферичної поверхні}\markboth{10–12. Поле зарядженої площини}{10–12. Поле зарядженої площини}


\srcline{\book{} т. 2, \textbf{§ 1.7} applications (с. 38–40); potentials in \textbf{§ 1.11} (с. 54–56). \ok{} Formulas checked against the scan.}
One method, three shapes, so they fit well on one double page with a shared "recipe" box.

\needspace{5\baselineskip}\section*{Recipe box (write it once, at the top)}
\begin{enumerate}
\item Use the symmetry to find the \textbf{direction} of E.
\item Choose a closed \textbf{Gaussian surface} on which E is constant or the flux is 0.
\item Compute the flux by definition (E·area) and set it equal to $q_{\text{inside}}/\varepsilon_0$.
\item Solve for E. Then get φ from $\varphi_1-\varphi_2 = \int E\,dl$.
\end{enumerate}

\par\noindent{\color{accent!30}\rule{\linewidth}{0.4pt}}\par

\needspace{5\baselineskip}\section*{10. Поле рівномірно зарядженої нескінченної площини (pp. 38–39, 54)}
\begin{itemize}
\item "Infinite" means the distance to the plane ≪ its size.
\item Symmetry: E ⟂ the plane. Surface: a \textbf{прямий циліндр} (right cylinder) crossing the plane, with bases S.
\item Flux: through the side = 0 and through the two bases = 2ES. Enclosed charge σS.
\item $2ES = \dfrac{\sigma S}{\varepsilon_0}$ ⟹ $E = \dfrac{\sigma}{2\varepsilon_0}$ (1.34)
\item E doesn't depend on distance, so the field on each side is \textbf{однорідне} (uniform).
\item Potential: $\varphi(x) = \varphi(0) - \dfrac{\sigma}{2\varepsilon_0}x$, decreasing linearly.
\item \pen{} Sketch рис. 1.18: the plane, a cylinder through it, and arrows out of both bases.
\end{itemize}

\needspace{5\baselineskip}\section*{11. Поле двох різнойменно заряджених паралельних площин (p. 39)}
\begin{itemize}
\item +σ and −σ. By \textbf{superposition}, the fields add between the planes and cancel outside.
\item Between: $E = \dfrac{\sigma}{\varepsilon_0}$ (1.35). Outside: E = 0.
\item Two planes with the \textbf{same} sign: E = 0 between them.
\item Potential between: $\varphi(x) = \varphi(0) - \dfrac{\sigma}{\varepsilon_0}x$, so $U = \dfrac{\sigma d}{\varepsilon_0}$, the parallel-plate capacitor of topic 15.
\item \pen{} Sketch рис. 1.19: two plates, three regions, the field arrows from each plate, and their sum.
\item Related: the field right next to a conductor is $E = \sigma/\varepsilon_0$ (1.36), p. 40 (topic 14).
\end{itemize}

\needspace{5\baselineskip}\section*{12. Поле рівномірно зарядженої сферичної поверхні (кулі) (pp. 40, 55)}
\begin{itemize}
\item Symmetry: E is radial. Surface: a \textbf{concentric sphere} of radius r.
\item Outside (r > R): $E\cdot4\pi r^2 = q/\varepsilon_0$ ⟹ $E = \dfrac{1}{4\pi\varepsilon_0}\dfrac{q}{r^2}$, the same as a point charge at the centre.
\item Inside a \textbf{сфера} (conducting ball, charge on the surface): E = 0 and φ = const = $\dfrac{q}{4\pi\varepsilon_0R}$.
\item Potential outside: $\varphi(r) = \dfrac{q}{4\pi\varepsilon_0 r}$.
\item Optional, a ball charged through its \textbf{volume} with ρ: inside, $E = \dfrac{\rho r}{3\varepsilon_0\varepsilon}$, which grows linearly.
\item \pen{} Plot E(r) and φ(r) for the shell: E jumps up at R then falls as 1/r²; φ is flat inside then falls as 1/r.
\end{itemize}

\needspace{5\baselineskip}\section*{Bonus row (book p. 40, cylinder)}
\begin{itemize}
\item An infinite charged cylinder or thread with linear density τ: $E = \dfrac{\tau}{2\pi\varepsilon_0 r}$ (1.37) outside, 0 inside a cylindrical surface.
\begin{itemize}
\item \warn{} \textbf{Misprint in the book} (p. 40): (1.37) is printed as E = η/(4πε₀r), but the line above it, E·2πrh = ηh/ε₀, gives \textbf{η/(2πε₀r)}. The book itself uses η/(2πε₀r) for the cylindrical capacitor (p. 62). The book writes the linear density as \textbf{η} here, not τ.
\end{itemize}
\end{itemize}

\needspace{5\baselineskip}\section*{Summary table (good for the page footer)}
{\small
\begin{longtable}{>{\raggedright\arraybackslash}p{0.180\linewidth}>{\raggedright\arraybackslash}p{0.240\linewidth}>{\raggedright\arraybackslash}p{0.324\linewidth}>{\raggedright\arraybackslash}p{0.192\linewidth}}
\toprule
\textbf{Розподіл} & \textbf{Гауссова поверхня} & \textbf{E} & \textbf{залежність від r} \\
\midrule
\endhead
площина σ & циліндр ⟂ & σ/2ε₀ & const \\
дві площини ±σ & — (суперпозиція) & σ/ε₀ між, 0 зовні & const \\
нитка/циліндр τ & коаксіальний циліндр & τ/2πε₀r & 1/r \\
сфера q & концентрична сфера & q/4πε₀r² зовні, 0 всередині & 1/r² \\
\bottomrule
\end{longtable}
}

\needspace{5\baselineskip}\section*{Self-check}
\begin{enumerate}
\item Why does the plane give 1/(2ε₀) but the conductor's surface 1/ε₀?
\item Why does E not depend on distance for a plane, when it does for a sphere?
\end{enumerate}

\chapter*{13. Теорема Остроградського–Гаусса}\addcontentsline{toc}{chapter}{13. Теорема Остроградського–Гаусса}\markboth{13. Теорема Остроградського–Гаусса}{13. Теорема Остроградського–Гаусса}


\srcline{\book{} т. 2, \textbf{§ 1.8} (с. 42); also used in § 2.1 (с. 109) for the continuity equation. \ok{} Formulas checked against the scan.}
Watch the naming. The syllabus says «Теорема Остроградського-Гауса». The book uses this name for the \textbf{mathematical formula} (the divergence theorem), and calls the physical law (topic 6) «теорема Гаусса». Keep them apart on the page.

\needspace{5\baselineskip}\section*{Layout}
\begin{layout}
┌ [A] Математична теорема: ∮A dS = ∫div A dV     ┐
├ [B] Застосування до E: (1.39) → (1.40) → (1.41) ┤
├ [C] Висновок: div E = ρ/ε₀ (1.42)              ┤
└ [D] Ще одне застосування: рівняння неперервності (2.12) ┘
\end{layout}

\needspace{5\baselineskip}\section*{[A] Формула Остроградського–Гаусса (p. 42)}
\begin{itemize}
\item For any vector field A, closed surface S and the volume V it bounds:
$\oint_S A_n\,dS = \int_V \operatorname{div}\vec A\,dV$
\item Words: the flux out through the boundary = the sum of the "sources" inside.
\item \pen{} Picture: a volume filled with small cubes; the fluxes across shared faces cancel, leaving only the outer surface.
\end{itemize}

\needspace{5\baselineskip}\section*{[B] Застосування до поля E (p. 42)}
\begin{enumerate}
\item Physics (Gauss, topic 6): $\oint_S E_n\,dS = \dfrac1{\varepsilon_0}\int_V\rho\,dV$ (1.39)
\item Mathematics (Ostrogradsky–Gauss): $\oint_S E_n\,dS = \int_V\operatorname{div}\vec E\,dV$ (1.40)
\item So: $\int_V\operatorname{div}\vec E\,dV = \dfrac1{\varepsilon_0}\int_V\rho\,dV$ (1.41)
\end{enumerate}

\needspace{5\baselineskip}\section*{[C] Висновок (p. 42)}
\begin{itemize}
\item V is \textbf{arbitrary}, so the integrands are equal: $\operatorname{div}\vec E = \dfrac{\rho}{\varepsilon_0}$ (1.42)
\item This is the same result as topic 7, reached differently: topic 7 used a small box, and this one uses the general theorem.
\end{itemize}

\needspace{5\baselineskip}\section*{[D] Bonus: рівняння неперервності (§ 2.1, p. 109)}
\begin{itemize}
\item The same trick on the current density j: $\dfrac{\partial\rho}{\partial t} = -\operatorname{div}\vec j$ (2.12), which is charge conservation in local form. Links to topics 1 and 20.
\end{itemize}

\needspace{5\baselineskip}\section*{Self-check}
\begin{enumerate}
\item Which step needs "V is arbitrary", and why?
\item Who proved what? (Ostrogradsky the maths, Gauss the electrostatics.)
\end{enumerate}

\chapter*{14. Провідник у зовнішньому електричному полі}\addcontentsline{toc}{chapter}{14. Провідник у зовнішньому електричному полі}\markboth{14. Провідник у зовнішньому електричному полі}{14. Провідник у зовнішньому електричному полі}


\srcline{\book{} т. 2, \textbf{§ 1.12} «Розподіл зарядів на поверхні провідника» (с. 56–59), \textbf{§ 1.13} «Провідники в електричному полі» (с. 59–60). \ok{} Formulas checked against the scan.}

\needspace{5\baselineskip}\section*{Layout}
\begin{layout}
┌ [A] Провідник: вільні носії; рівновага → E = 0 всередині ┐
├ [B] Наслідки: ρ = 0, заряд на поверхні, поверхня еквіпотенціальна, E ⟂ поверхні ┤
├ [C] Електростатична індукція (§1.13) + малюнок        ┤
├ [D] Порожнина, клітка Фарадея │ [E] Кривизна і σ, вістря ┤
└ footer: E біля поверхні = σ/ε₀ (1.36); Ван де Грааф (optional) ┘
\end{layout}

\needspace{5\baselineskip}\section*{[A] Провідник у рівновазі (p. 56)}
\begin{itemize}
\item \textbf{Провідники} have \textbf{вільні носії заряду} (free charge carriers): electrons in metals, ions in electrolytes.
\item In equilibrium the carriers don't drift, so \textbf{E = 0 inside}.
\end{itemize}

\needspace{5\baselineskip}\section*{[B] Наслідки (pp. 56–57)}
\begin{itemize}
\item From div E = ρ/ε₀ with E = 0, we get \textbf{ρ = 0 inside}: any excess charge sits \textbf{only on the outer surface}, described by σ.
\item The surface is \textbf{еквіпотенціальна}, and \textbf{E ⟂ the surface} at every point.
\item Just outside: $E = \sigma/\varepsilon_0$ (1.36), p. 40.
\end{itemize}

\needspace{5\baselineskip}\section*{[C] Електростатична індукція (pp. 59–60)}
\begin{itemize}
\item The book's term: «Явище перерозподілу електричних зарядів у провіднику під дією зовнішнього електричного поля … називають \textbf{електростатичною індукцією} або \textbf{електризацією через вплив}»
\item The electrons shift against E, so the near face (facing the incoming field) turns −, and the far face turns +.
\item The induced charges create an inner field \textbf{E′ opposite to E}. The shift stops once E′ = E, so the total inside is 0.
\item Charging by induction: put two touching conductors in the field, separate them, then remove the field. You are left with a + body and a − body.
\item \pen{} Sketch: a metal block in uniform E, with − on the left face, + on the right, and E′ inside pointing the other way.
\end{itemize}

\needspace{5\baselineskip}\section*{[D] Порожнина, захист (p. 57)}
\begin{itemize}
\item The field inside a closed conducting cavity is 0 (Faraday's experiments). This is the basis of \textbf{електростатичний захист} (electrostatic shielding).
\item A charged body inside a cavity gives up \textbf{all} its charge to the outer surface. Optional: the \textbf{електростатичний генератор} (Van de Graaff) works this way, pp. 57–58.
\end{itemize}

\needspace{5\baselineskip}\section*{[E] Кривизна поверхні (pp. 58–59)}
\begin{itemize}
\item Two balls joined by a wire end up with $\varphi_1=\varphi_2$, so $\dfrac{\sigma_1}{\sigma_2} = \dfrac{R_2}{R_1}$.
\item σ, and therefore E, is larger where the curvature is larger, so charge leaks off \textbf{вістря} (sharp points). This is how lightning rods work, and why power lines lose energy.
\end{itemize}

\needspace{5\baselineskip}\section*{Self-check}
\begin{enumerate}
\item Why E = 0 inside, from the carriers' point of view?
\item Where does excess charge go, and why not into the bulk?
\item Why does charge leak off sharp points?
\end{enumerate}

\chapter*{15. Конденсатори та їх ємність. Ємність плоского, циліндричного та сферичного конденсаторів}\addcontentsline{toc}{chapter}{15. Конденсатори та їх ємність. Ємність плоского, циліндричного та сферичного конденсаторів}\markboth{15. Конденсатори та їх ємність. Ємність плоского,…}{15. Конденсатори та їх ємність. Ємність плоского,…}


\srcline{\book{} т. 2, \textbf{§ 1.14} (с. 60–63). \ok{} Formulas checked against the scan.}

\needspace{5\baselineskip}\section*{Layout}
\begin{layout}
┌ [A] Ємність відокремленого провідника (1.68); фарад; куля C = 4πε₀εR ┐
├ [B] Конденсатор: означення, обкладки, (1.70); вплив діелектрика ε     ┤
├ [C] Плоский │ [D] Циліндричний │ [E] Сферичний  (три малюнки в ряд)  ┤
└ footer: робоча напруга / пробій                                        ┘
\end{layout}

\needspace{5\baselineskip}\section*{[A] Електрична ємність провідника (p. 60)}
\begin{itemize}
\item For an isolated conductor q/φ = const: $C = \dfrac{q}{\varphi}$ (1.68), $\Delta q = C\Delta\varphi$ (1.69)
\item Meaning: «електрична ємність провідника чисельно дорівнює зарядові, який потрібний для зміни потенціалу провідника на одиницю»
\item Unit: \textbf{фарад}, 1 Ф = 1 Кл/В. It is enormous: a ball of R = 9·10⁶ км would have 1 Ф. Hence мкФ = 10⁻⁶ Ф and пФ = 10⁻¹² Ф.
\item A sphere: $C = 4\pi\varepsilon_0\varepsilon R$.
\item C depends on the \textbf{size, the shape and the surrounding medium}, but not on the material or on hollows inside.
\end{itemize}

\needspace{5\baselineskip}\section*{[B] Конденсатор (p. 61)}
\begin{itemize}
\item «система з двох металевих електродів, розміщених на близькій відстані один від одного і розділених шаром діелектрика»; the electrodes are called \textbf{обкладки} (plates).
\item The field sits almost entirely between the plates, so C hardly depends on nearby bodies (unlike a single conductor).
\item $C = \dfrac{q}{\varphi_1-\varphi_2} = \dfrac{q}{U}$ (1.70)
\item Filling it with a dielectric: $C/C_0 = \varepsilon$.
\end{itemize}

\needspace{5\baselineskip}\section*{[C] Плоский (p. 61)}
\begin{itemize}
\item Two parallel plates of area S at distance d ≪ the plate size (ignore edge effects).
\item $U = Ed = \dfrac{\sigma}{\varepsilon_0}d$ (topic 11) ⟹ $C = \dfrac{\varepsilon_0\varepsilon S}{d}$
\end{itemize}

\needspace{5\baselineskip}\section*{[D] Циліндричний (pp. 61–62)}
\begin{itemize}
\item Two coaxial cylinders r₁ < r₂ of height h ≫ r.
\item $E = \dfrac{\tau}{2\pi\varepsilon_0\varepsilon r}$ ⟹ $U = \dfrac{\tau}{2\pi\varepsilon_0\varepsilon}\ln\dfrac{r_2}{r_1}$ ⟹ $C = \dfrac{2\pi\varepsilon_0\varepsilon h}{\ln(r_2/r_1)}$
\item A small gap d = r₂ − r₁ ≪ r₁ gives $C\approx\dfrac{2\pi\varepsilon_0\varepsilon r_1h}{d}$, which is a flat capacitor. \pen{} Worth noting why.
\end{itemize}

\needspace{5\baselineskip}\section*{[E] Сферичний (p. 62)}
\begin{itemize}
\item Two concentric spheres r₁ < r₂.
\item $U = \dfrac{q}{4\pi\varepsilon_0\varepsilon}\left(\dfrac1{r_1}-\dfrac1{r_2}\right)$ ⟹ $C = \dfrac{4\pi\varepsilon_0\varepsilon r_1r_2}{r_2-r_1}$
\item As r₂ → ∞ this becomes the single sphere $4\pi\varepsilon_0\varepsilon r_1$, a nice check.
\end{itemize}

\needspace{5\baselineskip}\section*{The same pattern three times}
Find E (Gauss) → integrate to get U → C = q/U. \pen{} Write that one line above the three drawings.

\needspace{5\baselineskip}\section*{Footer (p. 63)}
\begin{itemize}
\item \textbf{робоча напруга} (rated voltage) $U_{p}$: above it \textbf{пробій} (breakdown) occurs.
\item Optional: a two-wire line, $C = \pi\varepsilon_0\varepsilon l/\ln(d/r)$ (Example 1.3).
\end{itemize}

\needspace{5\baselineskip}\section*{Self-check}
\begin{enumerate}
\item Why do capacitors work better than isolated conductors?
\item How do the three C formulas show "C depends only on geometry and ε"?
\end{enumerate}

\chapter*{16. З'єднання конденсаторів}\addcontentsline{toc}{chapter}{16. З'єднання конденсаторів}\markboth{16. З'єднання конденсаторів}{16. З'єднання конденсаторів}


\srcline{\book{} т. 2, \textbf{§ 1.14} (с. 63–64). \ok{} Formulas checked against the scan. A short topic: half a page.}

\needspace{5\baselineskip}\section*{Layout (two columns)}
\begin{layout}
┌ Паралельне (рис. 1.31)        │ Послідовне (рис. 1.32)          ┐
│ U однакова                    │ q однаковий (електростатична індукція) │
│ q = q₁ + q₂ + …               │ U = U₁ + U₂ + …                  │
│ C = C₁ + C₂ + … (1.71)        │ 1/C = 1/C₁ + 1/C₂ + … (1.72)     │
│ n однакових: C = nC₀          │ n однакових: C = C₀/n            │
└ footer: навіщо (батареї); імпульс nU (рис. 1.33)                 ┘
\end{layout}

\needspace{5\baselineskip}\section*{Паралельне з'єднання (p. 63)}
\begin{itemize}
\item The same voltage U on every capacitor; the charges add: $q = (C_1+C_2+\ldots+C_n)U$.
\item $C = C_1 + C_2 + \ldots + C_n$ (1.71)
\end{itemize}

\needspace{5\baselineskip}\section*{Послідовне з'єднання (pp. 63–64)}
\begin{itemize}
\item \textbf{Electrostatic induction} gives every capacitor the same charge q; the voltages add.
\item $\dfrac1C = \dfrac1{C_1}+\dfrac1{C_2}+\ldots+\dfrac1{C_n}$ (1.72)
\item C is always \textbf{smaller than the smallest} Cᵢ.
\end{itemize}

\needspace{5\baselineskip}\section*{Why connect them (p. 63)}
\begin{itemize}
\item To get a needed capacitance at a needed voltage, in a \textbf{батарея конденсаторів} (capacitor bank). Series connection splits the voltage, so the bank can handle a higher one.
\item Trick (p. 64, рис. 1.33): charge n capacitors in parallel to U, then switch them to series and you get nU. Used for high-voltage pulses.
\end{itemize}

\needspace{5\baselineskip}\section*{\pen{} Your own line}
\begin{itemize}
\item Contrast with resistors (topic 22): capacitors in parallel \textbf{add up}, while resistors \textbf{add up in series}. Note why: C ∝ q/U and R ∝ U/I.
\end{itemize}

\needspace{5\baselineskip}\section*{Self-check}
\begin{enumerate}
\item Why is the charge the same on capacitors in series?
\item Two capacitors, 2 мкФ and 3 мкФ: give C for each connection. (5 and 1,2 мкФ.)
\end{enumerate}

\chapter*{17. Енергія зарядженого провідника}\addcontentsline{toc}{chapter}{17. Енергія зарядженого провідника}\markboth{17. Енергія зарядженого провідника}{17. Енергія зарядженого провідника}


\srcline{\book{} т. 2, \textbf{§ 1.25} «Енергія взаємодії електричних зарядів» (с. 101–103), \textbf{§ 1.26} start (с. 103–104). \ok{} Formulas checked against the scan.}

\needspace{5\baselineskip}\section*{Layout}
\begin{layout}
┌ [A] Енергія двох зарядів (1.128)–(1.129) ┐
├ [B] Система n зарядів (1.130): чому ½    ┤
├ [C] Заряджання провідника: dA = φ dq → (1.132) ┤
└ [D] Те саме для конденсатора (1.133)       ┘
\end{layout}

\needspace{5\baselineskip}\section*{[A] Два точкові заряди (pp. 101–102)}
\begin{itemize}
\item Bring the charges from infinity: $W = \dfrac{1}{4\pi\varepsilon_0}\dfrac{q_1q_2}{r_{12}}$ (1.128)
\item In symmetric form: $W = \tfrac12(q_1\varphi_1 + q_2\varphi_2)$ (1.129), where φ₁ is the potential that q₂ creates at q₁.
\end{itemize}

\needspace{5\baselineskip}\section*{[B] Система n зарядів (p. 102)}
\begin{itemize}
\item $W = \dfrac12\sum_{i=1}^n q_i\varphi_i$ (1.130), where φᵢ is the potential at qᵢ from \textbf{all the other} charges.
\item \textbf{Why ½}: each pair is counted twice in the sum.
\item Continuous charges: $W = \tfrac12\left(\int\rho\varphi\,dV + \int\sigma\varphi\,dS\right)$.
\end{itemize}

\needspace{5\baselineskip}\section*{[C] Енергія зарядженого провідника (pp. 103–104)}
\begin{itemize}
\item Charge the conductor in portions dq brought from infinity. Each costs $dA = \varphi\,dq = \dfrac{q}{C}dq$.
\item From 0 to q: $A = \dfrac{q^2}{2C}$, which becomes the stored energy.
\item $W = \dfrac{q\varphi}{2} = \dfrac{C\varphi^2}{2} = \dfrac{q^2}{2C}$ (1.132)
\item \pen{} Draw the graph: φ against q is a straight line through 0, and \textbf{W is the area under it}, a triangle. That is where the ½ comes from.
\end{itemize}

\needspace{5\baselineskip}\section*{[D] Конденсатор (p. 104)}
\begin{itemize}
\item The same formula with φ replaced by U: $W = \dfrac{CU^2}{2} = \dfrac{qU}{2} = \dfrac{q^2}{2C}$ (1.133)
\item Picture: carry charge from one plate to the other.
\end{itemize}

\needspace{5\baselineskip}\section*{Self-check}
\begin{enumerate}
\item Where does the ½ come from? Give two explanations: double counting and the triangle.
\item If q doubles, how does W change? (×4.)
\end{enumerate}

\chapter*{18. Енергія електростатичного поля}\addcontentsline{toc}{chapter}{18. Енергія електростатичного поля}\markboth{18. Енергія електростатичного поля}{18. Енергія електростатичного поля}


\srcline{\book{} т. 2, \textbf{§ 1.26} (с. 103–105). \ok{} Formulas checked against the scan.}

\needspace{5\baselineskip}\section*{Layout}
\begin{layout}
┌ [A] Від енергії конденсатора до енергії поля (1.133)→(1.134) ┐
├ [B] Неоднорідне поле (1.135)                                 ┤
├ [C] Об'ємна густина енергії w (1.136)                         ┤
└ [D] Де «живе» енергія: заряди чи поле? (близькодія vs далекодія) ┘
\end{layout}

\needspace{5\baselineskip}\section*{[A] Перетворення для плоского конденсатора (p. 104)}
\begin{itemize}
\item Start from $W = \tfrac12 C(\Delta\varphi)^2$ and put in $\Delta\varphi = Ed$ and $C = \dfrac{\varepsilon_0\varepsilon S}{d}$:
\item $W = \dfrac{\varepsilon_0\varepsilon E^2}{2}Sd = \dfrac{\varepsilon_0\varepsilon E^2}{2}V$ (1.134), where V = Sd is the volume between the plates.
\item The energy is now written through the \textbf{field E} and the \textbf{volume}, with no charges in it.
\end{itemize}

\needspace{5\baselineskip}\section*{[B] Загальний випадок (p. 104)}
\begin{itemize}
\item A non-uniform field: $W = \dfrac12\int_V \varepsilon_0\varepsilon E^2\,dV$ (1.135)
\end{itemize}

\needspace{5\baselineskip}\section*{[C] Густина енергії (p. 105)}
\begin{itemize}
\item \textbf{об'ємна густина енергії електричного поля}:
$w = \dfrac{W}{V} = \dfrac{\varepsilon_0\varepsilon E^2}{2} = \dfrac{ED}{2}$ (1.136)
\item Here $D = \varepsilon_0\varepsilon E$ is the \textbf{вектор електричної індукції} (electric displacement), defined in the dielectrics sections, §1.16+. One line on the page is enough.
\item Unit: Дж/м³.
\end{itemize}

\needspace{5\baselineskip}\section*{[D] Фізична інтерпретація (p. 105), worth a boxed note}
\begin{itemize}
\item \textbf{Близькодія}: the energy is carried by the \textbf{field} and spread through space with density w.
\item \textbf{Далекодія}: the energy sits on the charged conductors.
\item In \textbf{electrostatics no experiment can tell them apart}, because the field can't be separated from its charges.
\item \textbf{Changing fields} settle it: EM waves carry energy with no charges present, so the energy is in the field.
\end{itemize}

\needspace{5\baselineskip}\section*{Self-check}
\begin{enumerate}
\item Topic 17 and topic 18 give the same W. What is the difference between them?
\item Why does an EM wave prove the energy is in the field?
\end{enumerate}

\chapter*{19. Електричний струм. Постійний електричний струм. Закони Ома}\addcontentsline{toc}{chapter}{19. Електричний струм. Постійний електричний струм. Закони Ома}\markboth{19. Електричний струм. Постійний електричний струм.…}{19. Електричний струм. Постійний електричний струм.…}


\srcline{\book{} т. 2, \textbf{§ 2.1} (с. 105–110), \textbf{§ 2.2} (с. 110–115), \textbf{§ 2.3} (с. 118–120). \ok{} Formulas checked against the scan.}
Overlaps topics 20–21. Suggested split: \textbf{this page = what a current is plus all three forms of Ohm's law}; topic 20 = the quantities I, j, ЕРС; topic 21 = R, ρ, power.

\needspace{5\baselineskip}\section*{Layout}
\begin{layout}
┌ [A] Струм: означення, носії, напрям, види ┐
├ [B] Постійний струм: умови (2.3), (2.13)–(2.14) ┤
├ [C] Закон Ома для однорідної ділянки (2.16)–(2.17) ┤
├ [D] Диференціальна форма j = σE (2.23)              ┤
└ [E] Неоднорідна ділянка (2.33) і повне коло (2.34) ┘
\end{layout}

\needspace{5\baselineskip}\section*{[A] Електричний струм (pp. 105–106)}
\begin{itemize}
\item «\textbf{упорядкований рух} електричних зарядів»
\item \textbf{Носії заряду} (charge carriers): electrons in metals and semiconductors; ions in electrolytes; electrons and ions in gases.
\item \textbf{Напрям струму} (direction): the way \textbf{positive} charges move, so in metals it runs opposite to the electrons.
\item Kinds of current: \textbf{струм провідності} (conduction, inside a conductor), \textbf{конвекційний} (convection, the charged body itself moves), \textbf{струм поляризації} (polarization, in dielectrics).
\end{itemize}

\needspace{5\baselineskip}\section*{[B] Постійний струм (pp. 107, 110)}
\begin{itemize}
\item \textbf{стаціонарний або постійний} (steady): equal charges pass in equal times and the direction doesn't change. $I = \Delta q/\Delta t$ (2.3)
\item Continuity for a steady current: $\oint \vec j\,d\vec S = 0$ (2.13), $\operatorname{div}\vec j = 0$ (2.14). Current lines are \textbf{closed}, and I is the same along an unbranched wire.
\end{itemize}

\needspace{5\baselineskip}\section*{[C] Закон Ома (однорідна ділянка) (pp. 111–112)}
\begin{itemize}
\item «сила струму в однорідному провіднику прямо пропорційна напрузі на кінцях цього провідника»
\item $I = kU$ (2.16), where k is the \textbf{провідність} (conductance, unit См) → $I = \dfrac{U}{R}$ (2.17)
\item \textbf{спад напруги} (voltage drop) $U = \varphi_1-\varphi_2$ along a wire carrying current.
\end{itemize}

\needspace{5\baselineskip}\section*{[D] Диференціальна форма (p. 114)}
\begin{itemize}
\item A small cylinder dl·dS: $j\,dS = \dfrac{E\,dl}{\rho\,dl/dS}$ ⟹ $\vec j = \sigma\vec E$ (2.23)
\item «густина електричного струму в будь-якій точці фізично однорідного провідника пропорційна напруженості електричного поля»
\item σ depends \textbf{only on the material}, not on shape or size.
\item Ohm's law is \textbf{not fundamental}: it fails in strong fields, at contacts, and in diodes (нелінійні елементи), p. 115.
\end{itemize}

\needspace{5\baselineskip}\section*{[E] Неоднорідна ділянка і повне коло (pp. 119–120)}
\begin{itemize}
\item A section with a source: $IR_{12} = (\varphi_1-\varphi_2) + \mathcal E_{12}$ (2.33). Sign rule: \textbf{ЕРС > 0} if, moving along the current, you pass the source from − to +.
\item A closed circuit (φ₁ = φ₂): $I = \dfrac{\mathcal E}{R+r}$ (2.34), where r is the \textbf{внутрішній опір} (internal resistance).
\item $U = IR = \mathcal E - Ir$ (2.35). An open circuit gives U = ℰ. A short circuit gives $I_{max} = \mathcal E/r$ (2.37).
\end{itemize}

\needspace{5\baselineskip}\section*{Three forms of Ohm's law (good for the page footer)}
{\small
\begin{longtable}{>{\raggedright\arraybackslash}p{0.394\linewidth}>{\raggedright\arraybackslash}p{0.574\linewidth}}
\toprule
\textbf{ділянка} & \textbf{закон} \\
\midrule
\endhead
однорідна & I = U/R \\
неоднорідна & IR = φ₁ − φ₂ + ℰ \\
повне коло & I = ℰ/(R + r) \\
точка & j = σE \\
\bottomrule
\end{longtable}
}

\needspace{5\baselineskip}\section*{Self-check}
\begin{enumerate}
\item Why is a steady current impossible with Coulomb forces alone? (See topic 20.)
\item Which form of Ohm's law is local, and why does that matter?
\end{enumerate}

\chapter*{20. Сила струму, густина струму. Електрорушійна сила}\addcontentsline{toc}{chapter}{20. Сила струму, густина струму. Електрорушійна сила}\markboth{20. Сила струму, густина струму. Електрорушійна сила}{20. Сила струму, густина струму. Електрорушійна сила}


\srcline{\book{} т. 2, \textbf{§ 2.1} (с. 106–110), \textbf{§ 2.3} (с. 115–118). \ok{} Formulas checked against the scan.}

\needspace{5\baselineskip}\section*{Layout}
\begin{layout}
┌ [A] Сила струму I (2.1); ампер, кулон              ┐
├ [B] Густина струму j (2.4)–(2.10) + рис. 2.1 (циліндр v·dt) ┤
├ [C] Рівняння неперервності (2.11)–(2.12)          ┤
├ [D] Сторонні сили: навіщо + рис. 2.4 (кругообіг зарядів) ┤
└ [E] ЕРС: означення (2.28), циркуляція (2.29)–(2.30) ┘
\end{layout}

\needspace{5\baselineskip}\section*{[A] Сила струму (p. 106)}
\begin{itemize}
\item «скалярна величина, яка чисельно дорівнює зарядові, що переноситься крізь поперечний переріз провідника за одиницю часу»
\item $I = \dfrac{dq}{dt}$ (2.1). Unit: \textbf{ампер}, a base SI unit. 1 Кл = 1 А·1 с.
\end{itemize}

\needspace{5\baselineskip}\section*{[B] Густина струму (pp. 107–109)}
\begin{itemize}
\item $j = \dfrac{dI}{dS_\perp}$ (2.4), a \textbf{vector} pointing along the drift of positive carriers.
\item Through carriers: $\vec j = n_0 e\vec v_d$ (2.6), where n₀ is the concentration and $v_{d}$ the \textbf{дрейфова швидкість} (drift velocity).
\item For electrons: $\vec j = -n_0|e|\vec v_d$ (2.7).
\item $I = \int_S j_n\,dS$ (2.9). Uniform current: $I = n_0ev_dS = jS$ (2.10). Unit А/м².
\item Note: the average thermal velocity ⟨v⟩ = 0, so only the drift carries current.
\item \pen{} Sketch рис. 2.1а: a cylinder of base dS and length $v_{d}$·dt; everything inside it passes through in time dt.
\end{itemize}

\needspace{5\baselineskip}\section*{[C] Рівняння неперервності (pp. 109–110)}
\begin{itemize}
\item $\oint_S j_n\,dS = -\dfrac{\partial q}{\partial t}$ (2.11) ⟹ $\dfrac{\partial\rho}{\partial t} = -\operatorname{div}\vec j$ (2.12) via Ostrogradsky–Gauss (topic 13).
\item This is \textbf{charge conservation} in macroscopic electrodynamics (topic 1).
\end{itemize}

\needspace{5\baselineskip}\section*{[D] Сторонні сили (pp. 115–116)}
\begin{itemize}
\item Coulomb forces alone would equalize the potentials (E → 0) and stop the current.
\item The charges have to be pushed \textbf{against} E inside the source, which takes \textbf{сторонні сили}: «Будь-які сили неелектростатичної природи, які діють на заряджені частинки». They can be chemical, mechanical, magnetic, from light, and so on.
\item \textbf{Джерело струму} (current source): a device where external forces act (a cell, a generator, a thermocouple, a solar panel).
\item Their field $E_{\text{ст}}$ gives $\vec j = \sigma(\vec E + \vec E_{\text{ст}})$ (2.26).
\item \pen{} Sketch рис. 2.4: a loop where the wire runs φ₁→φ₂ and the source pushes φ₂→φ₁.
\end{itemize}

\needspace{5\baselineskip}\section*{[E] Електрорушійна сила (p. 118)}
\begin{itemize}
\item «Скалярну величину ℰ, яка дорівнює \textbf{роботі сторонніх сил, віднесеній до одиниці позитивного заряду}, називають електрорушійною силою (ЕРС)»
\item $\mathcal E = \dfrac{A_{\text{ст}}}{q}$ (2.28); in a closed loop $\mathcal E = \oint \vec E_{\text{ст}}d\vec l$ (2.29), the \textbf{circulation of the external-force field}. On a section 1–2: $\mathcal E_{12} = \int_1^2\vec E_{\text{ст}}d\vec l$ (2.30).
\item Despite the name, it's \textbf{not a force}: it's an energy characteristic of the source, measured in volts.
\item Contrast with topic 9: the circulation of E is 0 and the circulation of $E_{\text{ст}}$ is ℰ ≠ 0. That non-zero circulation is exactly what keeps the current going.
\end{itemize}

\needspace{5\baselineskip}\section*{Self-check}
\begin{enumerate}
\item Why is the drift velocity tiny, yet a lamp lights up instantly?
\item Why is "electromotive \textit{force}" a misleading name?
\end{enumerate}

\chapter*{21. Поняття опору. Потужність}\addcontentsline{toc}{chapter}{21. Поняття опору. Потужність}\markboth{21. Поняття опору. Потужність}{21. Поняття опору. Потужність}


\srcline{\book{} т. 2, \textbf{§ 2.2} (с. 111–114), \textbf{§ 2.4} (с. 121–123). \ok{} Formulas checked against the scan.}

\needspace{5\baselineskip}\section*{Layout}
\begin{layout}
┌ [A] Опір R: означення, Ом, звідки береться        ┐
├ [B] R = ρl/S (2.18); питомий опір; σ = 1/ρ (2.21) ┤
├ [C] Залежність від T (2.19)–(2.20) + графік ρ(T)  ┤
├ [D] Робота струму (2.38)–(2.41)                   ┤
└ [E] Потужність (2.43)–(2.46); ККД; макс. при R = r + графік P(R) ┘
\end{layout}

\needspace{5\baselineskip}\section*{[A] Електричний опір (pp. 111–112)}
\begin{itemize}
\item $R = 1/k$, the inverse of conductance. 1 Ом: 1 А flows at 1 В.
\item Where it comes from (metals): electrons \textbf{scatter} on lattice defects and \textbf{thermal vibrations} of the lattice.
\item It doesn't depend on I (for moderate currents); it depends on geometry, material and state (T, p).
\end{itemize}

\needspace{5\baselineskip}\section*{[B] Опір провідника (pp. 112–113)}
\begin{itemize}
\item $R = \rho\dfrac{l}{S}$ (2.18), where ρ is the \textbf{питомий опір} (resistivity), in Ом·м.
\item $\sigma = 1/\rho$ (2.21), the \textbf{питома електропровідність} (conductivity), in См/м.
\item Lowest ρ: silver, copper, gold, aluminium (used for wires). High ρ: nichrome (heaters), constantan and nickeline (reference resistors).
\item Classification by σ: провідники σ > 10⁶, діелектрики σ < 10⁻¹⁰, напівпровідники in between (См/м); checked against the scan.
\end{itemize}

\needspace{5\baselineskip}\section*{[C] Температурна залежність (pp. 112–113)}
\begin{itemize}
\item The \textbf{температурний коефіцієнт опору} $\alpha = \dfrac1\rho\dfrac{d\rho}{dT}$ (2.19). Metals: α > 0. Electrolytes: α < 0.
\item For metals near 0 °C, $\alpha\approx 1/273$ К⁻¹ and $\rho = \rho_0(1+\alpha t)$, or $\rho = \rho_0\alpha T$ (2.20).
\item Uses: resistance thermometers, thermostats, bolometers.
\end{itemize}

\needspace{5\baselineskip}\section*{[D] Робота струму (p. 121)}
\begin{itemize}
\item $A = IUt$ (2.38) $= I^2Rt = \dfrac{U^2}{R}t$ (2.39)
\item The whole source over the whole loop: $A_{\text{заг}} = \mathcal E It$ (2.41).
\end{itemize}

\needspace{5\baselineskip}\section*{[E] Потужність (pp. 121–123)}
\begin{itemize}
\item \textbf{Потужність} (power) is the work per unit time.
\item \textbf{Корисна} (useful, in the load R): $P = IU = I^2R = \dfrac{U^2}{R}$ (2.43) $= \dfrac{\mathcal E^2R}{(R+r)^2}$ (2.44)
\item \textbf{Загальна} (total): $P_{\text{заг}} = \mathcal EI = \dfrac{\mathcal E^2}{R+r}$ (2.45)
\item \textbf{ККД} (efficiency): $\eta = \dfrac{P}{P_{\text{заг}}} = \dfrac{R}{R+r}$ (2.46)
\item Maximum useful power when \textbf{R = r}: $P_{max} = \dfrac{\mathcal E^2}{4r}$, but then η = 50 \%.
\item \pen{} Graph P(R): zero at R = 0, a peak at R = r, tending to 0 as R → ∞. Overlay η(R) rising towards 1.
\end{itemize}

\needspace{5\baselineskip}\section*{Self-check}
\begin{enumerate}
\item Why can't you get maximum power and maximum efficiency at the same time?
\item Why do metals resist more when hot?
\end{enumerate}

\chapter*{22. Перший та другий закон Кірхгофа}\addcontentsline{toc}{chapter}{22. Перший та другий закон Кірхгофа}\markboth{22. Перший та другий закон Кірхгофа}{22. Перший та другий закон Кірхгофа}


\srcline{\book{} т. 2, \textbf{§ 2.6} «Електричне коло. Правила Кірхгофа» (с. 126–131). \ok{} Formulas checked against the scan.}
The book says \textbf{правила} Кірхгофа; the syllabus says «закон». Both are used.

\needspace{5\baselineskip}\section*{Layout}
\begin{layout}
┌ [A] Прості з'єднання: послідовне / паралельне (табличка)    ┐
├ [B] Терміни: розгалужене коло, вузол, ділянка, контур        ┤
├ [C] I правило (2.54) + рис. 2.12 │ [D] II правило (2.55) + рис. 2.13 ┤
└ [E] Як складати рівняння: знаки, незалежні контури           ┘
\end{layout}

\needspace{5\baselineskip}\section*{[A] Послідовне і паралельне з'єднання опорів (pp. 127–128)}
{\small
\begin{longtable}{>{\raggedright\arraybackslash}p{0.168\linewidth}>{\raggedright\arraybackslash}p{0.392\linewidth}>{\raggedright\arraybackslash}p{0.392\linewidth}}
\toprule
\textbf{} & \textbf{послідовне (рис. 2.7)} & \textbf{паралельне (рис. 2.9)} \\
\midrule
\endhead
струм & I = I₁ = I₂ = … & I = I₁ + I₂ + … \\
напруга & U = U₁ + U₂ + … & U = U₁ = U₂ = … \\
опір & R = R₁ + R₂ + … & 1/R = 1/R₁ + 1/R₂ + … \\
пропорція & U₁/U₂ = R₁/R₂ & I₁/I₂ = R₂/R₁ \\
\bottomrule
\end{longtable}
}
\begin{itemize}
\item \textbf{Дільник напруги (потенціометр)} (voltage divider), p. 127.
\item Optional, batteries of sources (pp. 128–129): in series ℰ = Σℰₖ, r = Σrₖ; n identical in parallel give r = r₁/n.
\end{itemize}

\needspace{5\baselineskip}\section*{[B] Терміни (p. 129)}
\begin{itemize}
\item \textbf{Розгалужене коло} (branched circuit): a network of several loops with shared sections.
\item \textbf{Вузол} (node): «Точку розгалуженого кола, в якій сходяться \textbf{три або більше} провідників із струмами»
\item \textbf{Ділянка} (branch): the part of a loop between two nodes, carrying one current.
\end{itemize}

\needspace{5\baselineskip}\section*{[C] Перше правило (p. 129)}
\begin{itemize}
\item It expresses \textbf{charge conservation} at a node: «алгебрична сума сил струмів, які сходяться у вузлі, дорівнює нулеві»
\item $\sum_{k=1}^n I_k = 0$ (2.54). Incoming currents are +, outgoing are −.
\item Why it holds: otherwise charge would pile up at the node and the potentials would change, so the current would not be steady.
\end{itemize}

\needspace{5\baselineskip}\section*{[D] Друге правило (p. 130)}
\begin{itemize}
\item A generalization of \textbf{Ohm's law} for a loop: «у контурі алгебрична сума спадів напруг (добутків сил струмів на опори відповідних ділянок) дорівнює алгебричній сумі електрорушійних сил, які діють у цьому контурі»
\item $\sum_{k=1}^n I_kR_k = \sum_{k=1}^m \mathcal E_k$ (2.55)
\item Proof idea: write (2.33) for each section of the loop and add them; \textbf{the potentials cancel}.
\end{itemize}

\needspace{5\baselineskip}\section*{[E] Як користуватись (pp. 129–131)}
\begin{enumerate}
\item Pick a direction for each current at random; a negative answer means it flows the other way.
\item Pick a direction to go round each loop. IR counts + if the current runs with that direction. ℰ counts + if going round takes you through the source from − to +.
\item Use only \textbf{independent} loops, ones that can't be built by combining others.
\item \pen{} Do one worked example yourself on рис. 2.13: three branches, two nodes.
\end{enumerate}

\needspace{5\baselineskip}\section*{Self-check}
\begin{enumerate}
\item Which conservation law does each rule express? (Charge; energy, through Ohm's law.)
\item How many independent node equations are there for N nodes? (N − 1.)
\end{enumerate}

\chapter*{23. Закон Джоуля–Ленца}\addcontentsline{toc}{chapter}{23. Закон Джоуля–Ленца}\markboth{23. Закон Джоуля–Ленца}{23. Закон Джоуля–Ленца}


\srcline{\book{} т. 2, \textbf{§ 2.5} «Закон Джоуля — Ленца. Закон збереження енергії електричного поля» (с. 123–126). \ok{} Formulas checked against the scan.}

\needspace{5\baselineskip}\section*{Layout}
\begin{layout}
┌ [A] Формулювання + Q = I²Rt (2.47)            ┐
├ [B] Механізм: зіткнення електронів з ґраткою   ┤
├ [C] Диференціальна форма: w = ρj² = σE² = jE (2.48)–(2.50) ┤
└ [D] Закон збереження енергії в колі (2.51), два наслідки ┘
\end{layout}

\needspace{5\baselineskip}\section*{[A] Закон (p. 123)}
\begin{itemize}
\item Joule (1841) and Lenz (1842), found by experiment.
\item «кількість теплоти, що виділяється в провіднику при проходженні по ньому постійного електричного струму, прямо пропорційна добуткові квадрата сили струму, опору провідника і часу проходження струму»
\item $Q = I^2Rt$ (2.47). It also equals $IUt = \dfrac{U^2}{R}t$ (for a wire at rest with no chemistry).
\item Changing current: $Q = \int_0^t I^2(t)R\,dt$.
\end{itemize}

\needspace{5\baselineskip}\section*{[B] Механізм (pp. 123–124)}
\begin{itemize}
\item Electrons gain energy from the field \textbf{between collisions} and hand it to the lattice ions \textbf{in collisions}. The ordered motion turns into chaotic motion, which is heat.
\item \pen{} A sketch: a sawtooth graph of electron energy over time, rising then dropping at each collision.
\end{itemize}

\needspace{5\baselineskip}\section*{[C] Диференціальна форма (p. 124)}
\begin{itemize}
\item \textbf{питома потужність струму} (specific power), the heat per unit volume per unit time: $w = \dfrac{dQ}{dV\,dt}$
\item $w = \rho j^2$ (2.48), $w = \sigma E^2$ (2.49), $w = \vec j\cdot\vec E$ (2.50)
\item Integrating over volume and time brings back (2.47). Local versus integral form again, as with Ohm's law.
\end{itemize}

\needspace{5\baselineskip}\section*{[D] Закон збереження енергії (pp. 125–126)}
\begin{itemize}
\item $dW = \delta Q + dU + \delta A$ (2.51): the source's energy goes into heat, a change in field energy, and mechanical work.
\item Consequences: (1) nothing moves, so all of the source's work becomes heat; (2) the charges are fixed, so mechanical work comes from the loss of field energy.
\item The details on ponderomotive forces are \textbf{optional}; skip them.
\end{itemize}

\needspace{5\baselineskip}\section*{Pairing with Ohm's law (good for the page footer)}
{\small
\begin{longtable}{>{\raggedright\arraybackslash}p{0.291\linewidth}>{\raggedright\arraybackslash}p{0.291\linewidth}>{\raggedright\arraybackslash}p{0.370\linewidth}}
\toprule
\textbf{} & \textbf{інтегральна} & \textbf{диференціальна} \\
\midrule
\endhead
Ом & I = U/R & j = σE \\
Джоуль–Ленц & Q = I²Rt & w = σE² = jE \\
\bottomrule
\end{longtable}
}

\needspace{5\baselineskip}\section*{Self-check}
\begin{enumerate}
\item Why does a thin wire in a series circuit heat up more? (Same I, bigger R.)
\item Why are heaters made of nichrome and not copper?
\end{enumerate}

\part{Молекулярна фізика і термодинаміка}

\chapter*{Програма МКР №2: 30 питань}\addcontentsline{toc}{chapter}{Програма МКР №2: 30 питань}\markboth{Програма МКР №2: 30 питань}{Програма МКР №2: 30 питань}



Source: \textbf{т. 1, Частина II, розд. 14–19 (с. 298–484)}. Every question below is mapped to a § from the section headings found in the OCR. Questions are copied from the photo.

\needspace{5\baselineskip}\section*{1. Question → textbook map}

{\footnotesize
\begin{longtable}{>{\raggedright\arraybackslash}p{0.023\linewidth}>{\raggedright\arraybackslash}p{0.286\linewidth}>{\raggedright\arraybackslash}p{0.302\linewidth}>{\raggedright\arraybackslash}p{0.092\linewidth}>{\raggedright\arraybackslash}p{0.216\linewidth}}
\toprule
\textbf{№} & \textbf{Питання (syllabus wording)} & \textbf{§ т. 1} & \textbf{с.} & \textbf{Wiki page (planned, in \texttt{topics/molecular/})} \\
\midrule
\endhead
1 & Закони ідеальних газів & 14.6, 14.7 & 311–317 & 01 \\
2 & Основне рівняння молекулярно-кінетичної теорії газу & 14.2, 14.3 (+14.4 температура) & 299–309 & 02 \\
3 & Внутрішня енергія термодинамічної системи & 16.1, 16.2 & 371–377 & 03 \\
4 & Теплота, робота, теплоємність & 16.1, 16.2 & 371–383 & 04 \\
5 & Теплоємність ідеальних газів & 16.5 & 390–397 & 05 \\
6 & Перший закон термодинаміки & 16.2 & 374–383 & 06 \\
7 & Ізотермічний процес & 16.4 (+14.7) & 386–390 & 07-09 \\
8 & Ізохоричний процес & 16.4 (+14.7) & 386–390 & 07-09 \\
9 & Ізобарний процес & 16.4 (+14.7) & 386–390 & 07-09 \\
10 & Адіабатичний процес & 16.3 «Адіабатний процес. Рівняння Пуассона» & 383–386 & 10 \\
11 & Колові процеси & 16.6 «Оборотні та необоротні процеси. Колові процеси (цикли)» & 397–399 & 11 \\
12 & Цикл Карно (необернений) та його ккд & 16.6, 16.7 & 399–403 & 12-13 \\
13 & Цикл Карно (обернений) та його ккд & 16.6, 16.7 (холодильна машина) & 399–403 & 12-13 \\
14 & Ентропія та її властивості & 16.8, 16.10 & 403–409, 411–414 & 14 \\
15 & Теорема Нернста & 16.11 «Третій закон термодинаміки» & 414–415 & 15 \\
16 & Другий закон термодинаміки & 16.8 «Другий закон термодинаміки. Теорема Карно» & 403–409 & 16 \\
17 & Закон розподілу Больцмана & 14.11 & 331–337 & 17,19 \\
18 & Закон розподілу Максвела & 14.8, 14.9, 14.10 & 317–331 & 18 \\
19 & Закон розподілу Максвела-Больцмана & 14.11 «Барометрична формула. Розподіл Максвелла — Больцмана» & 331–337 & 17,19 \\
20 & Закон рівномірного розподілу енергії по ступенях свободи & 16.2 (с. 376), 16.5 & 376, 390–392 & 20-21 \\
21 & Внутрішня енергія ідеальних газів & 16.5 & 390–392 & 20-21 \\
22 & Середня довжина вільного пробігу молекули в газах & 15.1 & 343–348 & 22 \\
23 & Дифузія в газах & 15.4 (+14.5) & 356–364 & 23-25 \\
24 & Внутрішнє тертя в газах & 15.2 «Внутрішнє тертя (в'язкість) у газах» & 348–352 & 23-25 \\
25 & Теплопровідність газів & 15.3 & 352–356 & 23-25 \\
26 & Відмінність реального газу від ідеального & 17.1 & 419–422 & 26 \\
27 & Рівняння Ван-Дер-Ваальса & 17.2, 17.3 (+17.4) & 422–433 & 27 \\
28 & Класифікація кристалів & 18.1 «Кристалічні та аморфні тіла» & 444–450 & 28 \\
29 & Дефекти в кристалах & 18.3 «Дефекти в реальних кристалах» & 453–456 & 29 \\
30 & Рідини (handwritten addition) & 19.1 (+19.2–19.4 optional) & 462–474 & 30 \\
\bottomrule
\end{longtable}
}

That is \textbf{30 questions → 23 pages}. The merges (7–9, 12–13, 17+19, 20–21, 23–25) put questions that share one method or formula on one page.

\needspace{5\baselineskip}\section*{2. Terminology: book vs syllabus}

Use the book's term on the page, and note the syllabus variant once.

{\small
\begin{longtable}{>{\raggedright\arraybackslash}p{0.350\linewidth}>{\raggedright\arraybackslash}p{0.618\linewidth}}
\toprule
\textbf{Syllabus} & \textbf{Book (т. 1)} \\
\midrule
\endhead
ступені свободи & \textbf{ступені вільності} (і) \\
адіабатичний процес & \textbf{адіабатний процес} \\
Максвела, Ван-Дер-Ваальса & \textbf{Максвелла}, \textbf{Ван-дер-Ваальса} \\
необернений / обернений цикл Карно & прямий / \textbf{обернений} цикл; check «необернений» against p. 399 \\
внутрішнє тертя & внутрішнє тертя \textbf{(в'язкість)} \\
теорема Нернста & третій закон термодинаміки (теорема Нернста) \\
\bottomrule
\end{longtable}
}

\needspace{5\baselineskip}\section*{3. Study order (not the syllabus order)}

The syllabus jumps around. Each block below builds on the previous one, so work through them in this order:

{\small
\begin{longtable}{>{\raggedright\arraybackslash}p{0.273\linewidth}>{\raggedright\arraybackslash}p{0.180\linewidth}>{\raggedright\arraybackslash}p{0.431\linewidth}>{\raggedright\arraybackslash}p{0.052\linewidth}}
\toprule
\textbf{Block} & \textbf{Pages} & \textbf{Why here} & \textbf{Size} \\
\midrule
\endhead
\textbf{A. Молекулярно-кінетична теорія} & 2 → 1 → 20–21 & the pressure and temperature of a gas from molecules; i, U = (i/2)νRT & 3 \\
\textbf{B. Статистичні розподіли} & 18 → 17,19 & Maxwell (speeds) → Boltzmann (potential field) → combined & 2 \\
\textbf{C. Явища переносу} & 22 → 23–25 & ⟨l⟩ first; then one shared pattern: D, η, λ ∝ ⅓⟨v⟩⟨l⟩ & 2 \\
\textbf{D. Перший закон} & 3 → 4 → 6 → 5 & U, Q, A → the first law → $C_{V}$, $C_{p}$, Mayer's equation & 4 \\
\textbf{E. Процеси} & 7–9 → 10 & a table: law, A, Q, ΔU, C for each process; then the adiabat & 2 \\
\textbf{F. Цикли і другий закон} & 11 → 12–13 → 16 → 14 → 15 & cycle → Carnot → 2nd law → entropy → 3rd law & 5 \\
\textbf{G. Реальні гази} & 26 → 27 & the deviations → the van der Waals equation and isotherms, critical state & 2 \\
\textbf{H. Конденсовані стани} & 28 → 29 → 30 & crystals → defects → liquids & 3 \\
\bottomrule
\end{longtable}
}

Concept pages to create along the way:
\begin{itemize}
\item \texttt{concepts/isoprocesses-table}: one table of all five processes. It's the most useful single page for the МКР.
\item \texttt{concepts/transport-phenomena}: diffusion, viscosity and heat conduction as one law, "flux = −coefficient × gradient".
\item \texttt{concepts/statistical-vs-thermodynamic}: two ways to describe the same gas (blocks A–C vs D–F).
\end{itemize}



\chapter*{2. Основне рівняння молекулярно-кінетичної теорії газу}\addcontentsline{toc}{chapter}{2. Основне рівняння молекулярно-кінетичної теорії газу}\markboth{2. Основне рівняння молекулярно-кінетичної теорії газу}{2. Основне рівняння молекулярно-кінетичної теорії газу}


\srcline{\book{} т. 1, \textbf{§ 14.2} (с. 299–300), \textbf{§ 14.3} (с. 301–305), \textbf{§ 14.4} (с. 305–309). \ok{} (14.5)–(14.10) checked against the scan.}
Block A, page 1 of 3. Start here: the gas laws (01) and U = (i/2)RT (20–21) both follow from this page.

\needspace{5\baselineskip}\section*{Layout}
\begin{layout}
┌ [A] Модель ідеального газу: 6 положень МКТ (стисло)        ┐
├ [B] Виведення тиску: куб l³, молекула б'ється об стінку, рис. 14.1 ┤
├ [C] Основне рівняння (14.6) у рамці + словами               ┤
├ [D] Температура: θ = ⅔⟨ε⟩, θ = kT → ⟨ε⟩ = 3/2 kT (14.10)     ┤
└ [E] p = n₀kT (14.13) → місток до сторінки 01                 ┘
\end{layout}

\needspace{5\baselineskip}\section*{[A] Ідеальний газ (pp. 299–300)}
\begin{itemize}
\item The book's definition (p. 300): we can neglect the size of the molecules («вважати молекули матеріальними точками») and their interaction at a distance, counting only collisions. «Такий газ називають \textbf{ідеальним}».
\item \textbf{Основні положення МКТ газів} (6 of them, p. 300). Copy them as a short list:
\begin{enumerate}
\item The molecules are identical spheres of negligible size, with no attraction or repulsion between them.
\item Even a small volume holds very many molecules, so statistical laws apply.
\item Collisions are \textbf{elastic}.
\item The time between collisions ≫ the time of a collision, so between collisions the motion is uniform and straight.
\item With no external forces the molecules are spread evenly through the volume.
\item All directions of velocity are equally likely, the \textbf{«молекулярний хаос»}.
\end{enumerate}
\item \textbf{Тепловий рух} (thermal motion): the completely chaotic motion of molecules (p. 300).
\end{itemize}

\needspace{5\baselineskip}\section*{[B] Виведення (pp. 301–303)}
\begin{itemize}
\item Cube of edge l (рис. 14.1). One molecule hits the right wall: Δp = 2m₀$v_{x.}$ It returns after Δt = 2l/$v_{x}$, so the average force is $f = m_0v_x^2/l$.
\item Sum over all n molecules and divide by S = l². The directions are equally likely, so $p_{x}$ = $p_{y}$ = $p_{z}$, which gives
$p = \dfrac23\dfrac1{l^3}\sum_{i=1}^n\dfrac{m_{0i}v_i^2}{2}$ (14.5)
\item \textbf{Середня квадратична швидкість} (root-mean-square speed): $\langle v_{\text{кв}}^2\rangle = \dfrac{v_1^2+v_2^2+\ldots+v_n^2}{n}$
\item \pen{} Sketch: a cube, one molecule bouncing off the right face, with arrows $v_{x}$ → and −$v_{x}$ ←.
\end{itemize}

\needspace{5\baselineskip}\section*{[C] Основне рівняння МКТ (p. 303) — put it in a box}
\[p = \frac23\,n_0\,\frac{m_0\langle v_{\text{кв}}^2\rangle}{2}\qquad(14.6)\]
\begin{itemize}
\item n₀ = n/V is the \textbf{number of molecules per unit volume}; $m_0\langle v_{\text{кв}}^2\rangle/2$ is the \textbf{average kinetic energy of translational motion of one molecule}.
\item In words (p. 303): «тиск ідеального газу дорівнює двом третинам середньої кінетичної енергії поступального руху молекул одиниці об'єму газу»
\item Why it matters: it links \textbf{microscopic} quantities (m₀, v) to a \textbf{macroscopic} one (p) that can be measured.
\item Pressure is a \textbf{статистичний параметр} (statistical parameter): it makes no sense for a single molecule (p. 304).
\item Units: Па = Н/м². Also 1 атм = 101 325 Па; 1 тор = 133,322 Па; 1 бар = 10⁵ Па.
\end{itemize}

\needspace{5\baselineskip}\section*{[D] Температура (pp. 305–308)}
\begin{itemize}
\item \textbf{Термодинамічна рівновага} (thermodynamic equilibrium): all macroscopic changes stop. \textbf{Релаксація} (relaxation) is the spontaneous approach to it; the \textbf{час релаксації} is how long it takes. Its existence is the \textbf{нульове начало термодинаміки} (zeroth law), p. 306.
\item The piston example (рис. 14.2): at equilibrium the average kinetic energies on both sides even out. So temperature "behaves like" ⟨ε⟩, and it is an \textbf{інтенсивний параметр} (it doesn't depend on the mass of the system).
\item Definition in energy units: $\Theta = \dfrac23\dfrac{m_0\langle v_{\text{кв}}^2\rangle}{2}$, which gives $p = n_0\Theta$ (14.8)
\item In kelvins: $\Theta = kT$ (14.9), where k is the \textbf{стала Больцмана} (Boltzmann constant), 1,38·10⁻²³ Дж/К.
\item \[\frac{m_0\langle v_{\text{кв}}^2\rangle}{2} = \frac32kT\qquad(14.10)\]
\item Meaning of k: by how much the energy of translational motion of one molecule changes when T changes by 1 К.
\end{itemize}

\needspace{5\baselineskip}\section*{[E] Наслідок (p. 311)}
\begin{itemize}
\item Combining (14.6) and (14.10): $p = n_0kT$ (14.13). This is the step to 01 gas laws.
\end{itemize}

\needspace{5\baselineskip}\section*{Self-check}
\begin{enumerate}
\item Where in the derivation are "elastic collisions" and "chaos" used?
\item Why can't we speak of the pressure of one molecule?
\item What happens to ⟨v²⟩ if T doubles? (×2, so $v_{\text{кв}}$ grows by √2.)
\end{enumerate}

\chapter*{1. Закони ідеальних газів}\addcontentsline{toc}{chapter}{1. Закони ідеальних газів}\markboth{1. Закони ідеальних газів}{1. Закони ідеальних газів}


\srcline{\book{} т. 1, \textbf{§ 14.6} «Рівняння стану ідеального газу» (с. 311–313), \textbf{§ 14.7} «Основні закони ідеального газу» (с. 313–317). \ok{} (14.13)–(14.18) checked against the scan; (14.22)–(14.31) taken from clear OCR.}
Block A, page 2 of 3. The book's approach: the \textbf{experimental} laws come out \textbf{as consequences} of p = n₀kT.

\needspace{5\baselineskip}\section*{Layout}
\begin{layout}
┌ [A] Рівняння стану: pV = nkT → Клапейрон → Клапейрон–Менделєєв ┐
├ [B] Таблиця трьох ізопроцесів (закон, рівняння, графік)         ┤
│     + три маленькі графіки p–V, V–T, p–T (рис. 14.5–14.7)       │
├ [C] Закон Авогадро                                               ┤
└ [D] Межі застосування моделі                                    ┘
\end{layout}

\needspace{5\baselineskip}\section*{[A] Рівняння стану ідеального газу (pp. 311–313)}
\begin{itemize}
\item \textbf{Термічне рівняння стану} (thermal equation of state): an equation that links p, V, T at thermal equilibrium, $F(p, V, T) = 0$ (14.12).
\item From the MKT: $p = n_0kT$ (14.13) ⟹ $pV = nkT$ (14.14), where n is the number of molecules.
\item For one mole: $pV_M = N_AkT$ (14.15). With \textbf{молярна газова стала} $R = kN_A = 8{,}31$ Дж/(моль·К):
$pV_M = RT$ (14.16), the \textbf{рівняння Клапейрона}.
\item For any amount, with ν = m/M moles:
\[pV = \frac mM RT,\quad\text{\text{або}}\quad pV = \frac{n}{N_A}RT\qquad(14.18)\]
the \textbf{рівняння Клапейрона–Менделєєва}.
\item Meaning of R (p. 313): the work done by 1 mole of ideal gas expanding \textbf{isobarically} when heated by 1 К. \pen{} Sketch the piston of рис. 14.4.
\end{itemize}

\needspace{5\baselineskip}\section*{[B] Ізопроцеси (m = const) (pp. 313–316)}
{\footnotesize
\begin{longtable}{>{\raggedright\arraybackslash}p{0.184\linewidth}>{\raggedright\arraybackslash}p{0.058\linewidth}>{\raggedright\arraybackslash}p{0.207\linewidth}>{\raggedright\arraybackslash}p{0.228\linewidth}>{\raggedright\arraybackslash}p{0.243\linewidth}}
\toprule
\textbf{процес} & \textbf{умова} & \textbf{закон (book's name)} & \textbf{рівняння} & \textbf{графік} \\
\midrule
\endhead
\textbf{ізотермічний} & T = const & закон \textbf{Бойля–Маріотта} & pV = const (14.22) & p–V: hyperbolas (рис. 14.5) \\
\textbf{ізобарний} & p = const & закон \textbf{Гей-Люссака} & V/T = const (14.25); V = V₀(1 + αt) (14.26) & V–T: straight lines to the origin (рис. 14.6) \\
\textbf{ізохорний} & V = const & "рівняння ізохорного процесу" & p/T = const (14.28); p = p₀(1 + βt) (14.29) & p–T: straight lines to the origin (рис. 14.7) \\
\bottomrule
\end{longtable}
}

\begin{itemize}
\item α = β = 0,003661 К⁻¹ ≈ 1/273, the \textbf{коефіцієнт об'ємного розширення} and the \textbf{термічний коефіцієнт тиску}.
\item Optional: the isothermal compressibility $\chi = -\dfrac1V\left(\dfrac{\partial V}{\partial p}\right)_T = \dfrac1p$ (14.24), and the link $\beta = \alpha/(p\chi)$ (14.31).
\item In the textbook the ізохорний law is \textbf{not named after Charles}. If your lecturer uses "закон Шарля", add it yourself.
\item \pen{} Draw the three small graphs with several curves each (T₁ < T₂, p₁ < p₂, V₁ < V₂), and mark which curve has the higher parameter.
\item This page has the laws only. Work, heat and ΔU for each process come in block E (07–09).
\end{itemize}

\needspace{5\baselineskip}\section*{[C] Закон Авогадро (p. 316)}
\begin{itemize}
\item From pV = nkT: «в однакових об'ємах ідеальних газів при однакових температурах і тисках міститься однакове число молекул»
\end{itemize}

\needspace{5\baselineskip}\section*{[D] Межі застосування (p. 313)}
\begin{itemize}
\item The model holds when the molecules' own volume and their interaction can be neglected: \textbf{low pressure and high enough temperature}. Normal atmospheric conditions mostly qualify. Leads to 26 реальний газ.
\end{itemize}

\needspace{5\baselineskip}\section*{Optional (Example 14.1, p. 316)}
\begin{itemize}
\item Derive pV/T = const from experiment alone: isotherm 1→1′, then isochore 1′→2. A good exam trick, and it's half a page.
\end{itemize}

\needspace{5\baselineskip}\section*{Self-check}
\begin{enumerate}
\item Which equation needs the number of molecules, and which needs only the mass?
\item Why do the isobars on a V–T graph all pass through the origin (in kelvins)?
\item What is the physical meaning of R?
\end{enumerate}

\chapter*{20. Закон рівномірного розподілу енергії по ступенях свободи · 21. Внутрішня енергія ідеальних газів}\addcontentsline{toc}{chapter}{20. Закон рівномірного розподілу енергії по ступенях свободи · 21. Внутрішня енергія ідеальних газів}\markboth{20. Закон рівномірного розподілу енергії по ступенях свободи}{20. Закон рівномірного розподілу енергії по ступенях свободи}


\srcline{\book{} т. 1, \textbf{§ 16.2} (с. 375–377, inside «Перший закон термодинаміки»), \textbf{§ 16.5} «Теплоємність ідеального газу» (с. 390–392). \ok{} (16.3), (16.4), (16.46) and the $C_{V}$ values checked against the scan.}
Block A, page 3 of 3. \warn{} Terminology: the book says \textbf{«ступені вільності»}, the syllabus «ступені свободи». Use the book's word and note the variant once.

\needspace{5\baselineskip}\section*{Layout}
\begin{layout}
┌ [A] Ступені вільності: означення, i = i_пост + i_оберт + i_колив (16.4) ┐
├ [B] Закон рівного розподілу (рамка): ½kT на ступінь; коливальний → kT ┤
├ [C] Таблиця: газ → i → U → C_V, C_p + малюнки молекул (1, 2, 3 атоми)  ┤
├ [D] Внутрішня енергія ідеального газу: U = (i/2)RT; тільки від T      ┤
└ [E] Межі класичної теорії (спирт: теорія 3R, дослід ≈ 8R)            ┘
\end{layout}

\needspace{5\baselineskip}\section*{[A] Ступені вільності (p. 376)}
\begin{itemize}
\item «Число таких незалежних координат називають \textbf{кількістю ступенів вільності} й позначають i». These are the coordinates needed to fix the positions of all the points of a system.
\item A free material point has i = 3 (translation). A free rigid body has i = 6 (3 translational + 3 rotational).
\item $i = i_{\text{пост}} + i_{\text{оберт}} + i_{\text{колив}}$ (16.4)
\end{itemize}

\needspace{5\baselineskip}\section*{[B] Закон рівного розподілу (p. 377) — put it in a box}
\begin{itemize}
\item The book's name: \textbf{«закон рівного розподілу енергії за класичними ступенями вільності (закон рівного розподілу Больцмана)»}. On p. 392 it's also called «закон про \textbf{рівномірний} розподіл енергії за ступенями вільності», which is close to the syllabus wording.
\item Wording (p. 377): «у стані теплової рівноваги на кожний класичний ступінь вільності руху частинок системи припадає кінетична енергія, що дорівнює в середньому \textbf{½kT}»
\item Rotational degree of freedom: ½kT. \textbf{Vibrational: kT}, made up of ½kT kinetic plus ½kT potential.
\item General formula: $\langle w\rangle = (3 + r + 2s)\dfrac{kT}{2}$ (16.45), where r is the number of rotational and s the number of vibrational degrees of freedom.
\item It is broken where quantum effects matter (p. 377).
\end{itemize}

\needspace{5\baselineskip}\section*{[C] Таблиця (pp. 391–392), rigid molecules (s = 0)}
{\footnotesize
\begin{longtable}{>{\raggedright\arraybackslash}p{0.199\linewidth}>{\raggedright\arraybackslash}p{0.260\linewidth}>{\raggedright\arraybackslash}p{0.138\linewidth}>{\raggedright\arraybackslash}p{0.153\linewidth}>{\raggedright\arraybackslash}p{0.077\linewidth}>{\raggedright\arraybackslash}p{0.077\linewidth}}
\toprule
\textbf{газ} & \textbf{приклади} & \textbf{i} & \textbf{U (1 моль)} & \textbf{$C_{V}$} & \textbf{$C_{p}$} \\
\midrule
\endhead
одноатомний & He, Ar, Hg vapour & 3 & 3/2 RT & 3/2 R & 5/2 R \\
двохатомний & H₂, O₂, N₂ & 5 (3 + 2) & 5/2 RT & 5/2 R & 7/2 R \\
багатоатомний & H₂O vapour, CH₄ & 6 (3 + 3) & 3RT & 3R & 4R \\
\bottomrule
\end{longtable}
}
\begin{itemize}
\item $C_{p}$ = $C_{V}$ + R. The \textbf{Mayer equation} itself belongs in topic 5; here just use it.
\item \pen{} Draw three small molecules (a dot, a dumbbell, a triangle) with arrows for the axes of rotation. The dumbbell has only 2 rotational axes; rotation about its own axis doesn't count.
\item Optional: CO₂ (linear) gives ⟨w⟩ = 6,5kT, and H₂O (non-linear) gives 6kT when vibrations are counted (16.43–16.44).
\end{itemize}

\needspace{5\baselineskip}\section*{[D] Внутрішня енергія ідеального газу (pp. 376, 390–392)}
\begin{itemize}
\item Generally U is \textbf{кінетична енергія теплового руху + потенціальна енергія взаємодії} (p. 375). Topic 3 covers this in full.
\item In an ideal gas there is \textbf{no interaction}, so U is kinetic energy only: «внутрішня енергія ідеального газу дорівнює сумарній кінетичній енергії його молекул» (p. 390).
\item One mole of a monatomic gas: $U = \dfrac32RT$ (16.2). In general:
\[U = \frac i2RT\qquad(16.3)\]
and for ν moles, $U = \nu\dfrac i2RT$.
\item Key conclusion (p. 392): «внутрішня енергія ідеального газу є функцією \textbf{тільки абсолютної температури}». It does not depend on V. The block E processes rest on this.
\end{itemize}

\needspace{5\baselineskip}\section*{[E] Межі класичної теорії (p. 392)}
\begin{itemize}
\item Simple molecules near room temperature agree with experiment. For complex ones, such as C₂H₅OH, theory gives $C_{V}$ = 3R but experiment gives about 8R.
\item $C_{V}$ = 5/2 R for diatomic gases holds only within a limited range of temperature, because vibrations "switch on" at high T. This is a quantum effect.
\end{itemize}

\needspace{5\baselineskip}\section*{Self-check}
\begin{enumerate}
\item Why does a dumbbell molecule have only 2 rotational degrees of freedom?
\item Why does a vibrational degree of freedom carry kT and not ½kT?
\item Why doesn't an ideal gas's U depend on volume?
\end{enumerate}

\chapter*{18. Закон розподілу Максвела}\addcontentsline{toc}{chapter}{18. Закон розподілу Максвела}\markboth{18. Закон розподілу Максвела}{18. Закон розподілу Максвела}


\srcline{\book{} т. 1, \textbf{§ 14.8} «Швидкості газових молекул та вимірювання їх» (с. 317–320), \textbf{§ 14.9} «Поняття про статистичний розподіл. Функція розподілу. Розподіл молекул за швидкостями (розподіл Максвелла)» (с. 320–330), \textbf{§ 14.10} (с. 331). \ok{} (14.54)–(14.63) checked against the scan.}
Block B, page 1 of 2. Spelling: the book writes \textbf{Максвелла} (double л); the syllabus writes «Максвела».

\needspace{5\baselineskip}\section*{Layout}
\begin{layout}
┌ [A] Постановка: чому «скільки молекул в інтервалі v…v+dv», а не «з швидкістю v» ┐
├ [B] Функція розподілу f(v): означення (14.42), нормування (14.43)               ┤
├ [C] Формула Максвелла (14.57) у рамці + ГРАФІК f(v) для T₁ < T₂ (рис. 14.13)      ┤
├ [D] Три швидкості: v_н, ⟨v⟩, v_кв (14.58), (14.61), (14.62) + співвідношення     ┤
└ [E] Дослідна перевірка: Штерн (рис. 14.9), Ламмерт (рис. 14.15)                  ┘
\end{layout}

\needspace{5\baselineskip}\section*{[A] Постановка задачі (p. 320)}
\begin{itemize}
\item The speed is continuous, so the probability of an \textit{exact} speed is 0. The right question: «скільки з усіх наявних молекул мають швидкості, значення яких лежать у певному інтервалі швидкостей в околі заданої швидкості?»
\item \textbf{Простір швидкостей} (velocity space): each molecule is a point ($v_{x}$, $v_{y}$, $v_{z}$), рис. 14.11.
\end{itemize}

\needspace{5\baselineskip}\section*{[B] Функція розподілу (pp. 322–323)}
\begin{itemize}
\item $dW = \dfrac{dn}{n} = f(\vec v)\,d\omega$ (14.42). f is the \textbf{функція розподілу молекул за швидкостями}, also called the \textbf{густина ймовірності} (probability density).
\item \textbf{Умова нормування}: $\int f\,d\omega = 1$ (14.43): the molecule certainly has \textit{some} speed.
\item Directions are equally likely, so f depends only on |v|. The components are independent, so $f = f(v_x)f(v_y)f(v_z)$ (14.46). Maxwell (1859).
\item Skip the derivation (pp. 323–327); just note the idea: "independence + isotropy ⟹ Gaussian".
\end{itemize}

\needspace{5\baselineskip}\section*{[C] Закон Максвелла (p. 327) — put it in a box}
\[\frac{dn}{n} = \frac{4}{\sqrt\pi}\left(\frac{m_0}{2kT}\right)^{3/2}e^{-\frac{m_0v^2}{2kT}}v^2\,dv\qquad(14.57)\]
\begin{itemize}
\item Words: the \textbf{relative number of molecules} whose speeds lie between v and v + dv.
\item Per component: $f(v_x) = \dfrac{1}{\sqrt\pi\,v_\text{н}}e^{-v_x^2/v_\text{н}^2}$ (14.54), a bell curve centred on 0 (рис. 14.14).
\item Shape of f(v) (pp. 327–328): f = 0 at v = 0 (no molecules at rest) and f → 0 as v → ∞. It is \textbf{asymmetric}: a steep rise and a slow fall. The area under it is 1.
\item \textbf{Heating} (p. 328): the peak shifts right (∝ √T) and gets lower and wider, because the area stays 1.
\item \pen{} The main drawing: two curves for T₁ < T₂, the same area, with $v_{\text{н}}$, ⟨v⟩ and $v_{\text{кв}}$ marked on one of them. Shade one strip dv as dn/n.
\item An example from the book (Table 3, O₂ at 0 °C): about ¾ of the molecules move at 200–600 м/с.
\end{itemize}

\needspace{5\baselineskip}\section*{[D] Характерні швидкості (pp. 328–330)}
{\small
\begin{longtable}{>{\raggedright\arraybackslash}p{0.468\linewidth}>{\raggedright\arraybackslash}p{0.393\linewidth}>{\raggedright\arraybackslash}p{0.091\linewidth}}
\toprule
\textbf{швидкість} & \textbf{формула} & \textbf{№} \\
\midrule
\endhead
\textbf{найбільш імовірна} $v_{\text{н}}$ (peak of f) & $\sqrt{2kT/m_0} = \sqrt{2RT/M}$ & (14.58) \\
\textbf{середня (арифметична)} ⟨v⟩ & $\sqrt{8kT/(\pi m_0)}$ & (14.61) \\
\textbf{середня квадратична} $v_{\text{кв}}$ & $\sqrt{3kT/m_0}$ & (14.62) \\
\bottomrule
\end{longtable}
}
\begin{itemize}
\item $v_{\text{кв}} : \langle v\rangle : v_\text{н} \approx 1{,}2 : 1{,}1 : 1$
\item Universal form with u = v/$v_{\text{н}}$: $\dfrac{dn}{n\,du} = \dfrac4{\sqrt\pi}e^{-u^2}u^2$ (14.63). It doesn't depend on the gas or on T.
\item Scale (p. 318): H₂ ≈ 1,9 км/с, N₂ ≈ 493 м/с, O₂ ≈ 461 м/с at room temperature.
\end{itemize}

\needspace{5\baselineskip}\section*{[E] Дослідна перевірка (pp. 318–320, 331)}
\begin{itemize}
\item \textbf{Stern's experiment (1920)}: a silver-coated wire inside two rotating cylinders. The deposited strip shifts and \textbf{smears}, which shows the atoms have a spread of speeds with a peak. \pen{} Sketch рис. 14.9 from above.
\item \textbf{Lammert (1929)}: two rotating discs with slits offset by an angle φ act as a speed selector. It confirmed Maxwell's law with high accuracy.
\end{itemize}

\needspace{5\baselineskip}\section*{Self-check}
\begin{enumerate}
\item Why is the probability of an exact speed zero?
\item Why does the peak get lower when the gas is heated?
\item Order $v_{\text{н}}$, ⟨v⟩ and $v_{\text{кв}}$ from smallest to largest, and explain why.
\end{enumerate}

\chapter*{17. Закон розподілу Больцмана · 19. Закон розподілу Максвела-Больцмана}\addcontentsline{toc}{chapter}{17. Закон розподілу Больцмана · 19. Закон розподілу Максвела-Больцмана}\markboth{17. Закон розподілу Больцмана}{17. Закон розподілу Больцмана}


\srcline{\book{} т. 1, \textbf{§ 14.11} «Барометрична формула. Розподіл Максвелла — Больцмана» (с. 331–337). \ok{} (14.64)–(14.75) checked against the scan.}
Block B, page 2 of 2. Logic: Maxwell assumed \textbf{no external field}. Gravity added → barometric formula → Boltzmann (any potential field) → combined with Maxwell.

\needspace{5\baselineskip}\section*{Layout}
\begin{layout}
┌ [A] Барометрична формула: стовп повітря, dp = −ρg dh → (14.68) + рис. 14.16 ┐
├ [B] Розподіл Больцмана (14.71) у рамці + графік n(h) (рис. 14.18)          ┤
├ [C] Розподіл Максвелла–Больцмана (14.74): два незалежні множники            ┤
└ [D] Наслідки: атмосфери планет, флуктуації (optional)                     ┘
\end{layout}

\needspace{5\baselineskip}\section*{[A] Барометрична формула (pp. 331–333)}
\begin{itemize}
\item Only the combined action of \textbf{thermal motion + gravity} gives the atmosphere its shape. Without gravity the molecules would spread out; without thermal motion they would fall to the ground.
\item A thin layer dh: $dp = -\rho g\,dh$ (14.64), with ρ = m₀n₀ and n₀ = p/kT ⟹ $\dfrac{dp}{p} = -\dfrac{m_0g}{kT}dh$ (14.66)
\item Integrate:
\[p = p_0e^{-\frac{m_0g(h-h_0)}{kT}}\ (14.67),\qquad p = p_0e^{-\frac{Mg(h-h_0)}{RT}}\ (14.68)\]
the \textbf{барометрична формула} (Laplace, 1821).
\item Assumptions: an ideal gas, \textbf{T = const}, \textbf{g = const}, so it is accurate only up to a few km.
\item The pressure falls exponentially, and faster for heavier gases (larger M) and colder ones (smaller T).
\item \pen{} Sketch рис. 14.16: a column of air, a layer dh, p and p − dp.
\end{itemize}

\needspace{5\baselineskip}\section*{[B] Розподіл Больцмана (p. 335) — put it in a box}
\begin{itemize}
\item Since p = nkT: $n = n_0e^{-\frac{m_0g(h-h_0)}{kT}}$ (14.70)
\item General form:
\[n = n_0e^{-\frac{\Delta U}{kT}}\qquad(14.71)\]
where ΔU is the difference in a molecule's \textbf{potential energy} (ΔU = m₀gh in gravity).
\item Boltzmann showed it is \textbf{universal}: it holds for any particles in thermal equilibrium in \textbf{any potential field}. «Цей статистичний закон називається \textbf{розподілом Больцмана}»
\item Meaning: the concentration is higher where U is lower, which is the tendency towards stable equilibrium.
\item In a layer: $\dfrac{dn}{n} = e^{-\frac{m_0g}{kT}dh}$ (14.73). \warn{} Copied as printed on p. 336, but it can't be right: dh sits inside the exponent. The fraction of molecules in the layer h…h + dh is $\dfrac{dn}{n} = \dfrac{m_0g}{kT}e^{-\frac{m_0gh}{kT}}dh$.
\item \pen{} Graph рис. 14.18: n against h, an exponential decay from n₀.
\end{itemize}

\needspace{5\baselineskip}\section*{[C] Розподіл Максвелла–Больцмана (p. 336)}
\begin{itemize}
\item Substitute n from (14.71) into Maxwell's (14.57):
\[dn = \frac4{\sqrt\pi}n_0\left(\frac{m_0}{2kT}\right)^{3/2}e^{-\frac{m_0\Delta\varphi}{kT}}e^{-\frac{m_0v^2}{2kT}}v^2\,dv\qquad(14.74)\]
where φ = gh is the potential.
\item \textbf{Two independent factors}: the Boltzmann factor $e^{-m_0\Delta\varphi/kT}$ (position) × the Maxwell factor $e^{-m_0v^2/2kT}$ (speed).
\item So the Boltzmann distribution doesn't depend on the speed distribution, and Maxwell's doesn't depend on the potential.
\item \pen{} Memory aid: $e^{-(E_\text{п} + E_\text{к})/kT}$, energy divided by kT in the exponent.
\end{itemize}

\needspace{5\baselineskip}\section*{[D] Optional}
\begin{itemize}
\item Planets' atmospheres (pp. 333–334): only molecules in the "tail" of Maxwell's distribution exceed escape velocity, so light planets (the Moon) lose their atmosphere.
\item Fluctuations (Smoluchowski, p. 335): $\delta = \sqrt{\dfrac2\pi\dfrac1n}$ (14.72), smaller when there are more particles.
\item Perrin's experiment (§ 14.12) used Boltzmann's law to measure $N_{A.}$
\end{itemize}

\needspace{5\baselineskip}\section*{Self-check}
\begin{enumerate}
\item Which two "forces" set the shape of the atmosphere?
\item Why is the barometric formula wrong at large altitudes? (T and g change.)
\item What does it mean that the Maxwell–Boltzmann distribution splits into two factors?
\end{enumerate}

\chapter*{22. Середня довжина вільного пробігу молекули в газах}\addcontentsline{toc}{chapter}{22. Середня довжина вільного пробігу молекули в газах}\markboth{22. Середня довжина вільного пробігу молекули в газах}{22. Середня довжина вільного пробігу молекули в газах}


\srcline{\book{} т. 1, \textbf{§ 15.1} «Середня довжина вільного пробігу молекул. Число зіткнень» (с. 343–348). \ok{} (15.1), (15.2) checked against the scan.}
Block C, page 1 of 2. \warn{} Notation: the book writes the free path as \textbf{⟨l⟩}; λ is used later for thermal conductivity.

\needspace{5\baselineskip}\section*{Layout}
\begin{layout}
┌ [A] Чому важливо: зигзаг, зіткнення → рівновага                  ┐
├ [B] Модель: «циліндр» πd²⟨v⟩ (рис. 15.1) → ⟨z⟩ (15.1) → ⟨l⟩ (15.2) ┤
├ [C] Ефективний переріз σ = πd²; ефективний діаметр (Сезерленд)    ┤
├ [D] Залежності: ⟨l⟩ ~ 1/p; числа для N₂                          ┤
└ [E] Послаблення пучка n = n₀e^(−x/⟨l⟩) (15.5); дослід Борна        ┘
\end{layout}

\needspace{5\baselineskip}\section*{[A] Навіщо (p. 343)}
\begin{itemize}
\item Molecules move at hundreds of м/с, but along a \textbf{zigzag}, so they don't get far. Collisions set up equilibrium and the Maxwell distribution.
\item \textbf{Середня довжина вільного пробігу} ⟨l⟩: «середню відстань, яку проходить молекула між двома послідовними зіткненнями»
\item The model: molecules are identical, perfectly elastic spheres of diameter d; only collisions between \textbf{pairs} count.
\end{itemize}

\needspace{5\baselineskip}\section*{[B] Виведення (pp. 343–344)}
\begin{itemize}
\item Freeze all the molecules but one. In 1 с it sweeps out a cylinder of base πd² and length ⟨v⟩, and hits every molecule whose centre lies inside: ⟨z⟩ = πd²n₀⟨v⟩.
\item The others move too, so use the relative speed $\langle v\rangle_{\text{відн}} = \sqrt2\langle v\rangle$ (from Maxwell's distribution):
\[\langle z\rangle = \sqrt2\,\pi d^2n_0\langle v\rangle\qquad(15.1)\]
\[\langle l\rangle = \frac{\langle v\rangle}{\langle z\rangle} = \frac{1}{\sqrt2\,\pi d^2n_0}\qquad(15.2)\]
\item \pen{} Sketch рис. 15.1: a molecule, a "straightened" cylinder of radius d, and molecules whose centres are inside it.
\end{itemize}

\needspace{5\baselineskip}\section*{[C] Ефективний переріз (pp. 344–345)}
\begin{itemize}
\item \textbf{Повний ефективний поперечний переріз} σ = πd². It is measured in m² but has no geometric meaning: it describes the \textbf{probability} of the collision process.
\item The \textbf{ефективний діаметр} shrinks as T rises (Sutherland): $d^2 = d_0^2(1 + C/T)$ (15.3), so ⟨l⟩ grows with T.
\end{itemize}

\needspace{5\baselineskip}\section*{[D] Порядки величин і залежності (p. 345)}
\begin{itemize}
\item N₂ at normal conditions: ⟨l⟩ ≈ 0,6·10⁻⁷ м, about \textbf{1000 times the size of a molecule}. Collisions are rare even at atmospheric pressure.
\item At T = const, \textbf{⟨l⟩ ∝ 1/p}. When ⟨l⟩ reaches the size of the container the gas is in a \textbf{вакуум} state (§ 15.5).
\end{itemize}

\needspace{5\baselineskip}\section*{[E] Послаблення пучка (pp. 345–347), optional}
\begin{itemize}
\item $-\dfrac{dn}{n} = \dfrac{dx}{\langle l\rangle}$ (15.4) ⟹ $n = n_0e^{-x/\langle l\rangle}$ (15.5). Only about 1/e ≈ ⅓ of molecules travel further than ⟨l⟩.
\item Born's experiment (1921): a beam of silver atoms through four diaphragms gives ⟨l⟩ directly (15.9).
\end{itemize}

\needspace{5\baselineskip}\section*{Self-check}
\begin{enumerate}
\item Why √2 in (15.1)?
\item Why does ⟨l⟩ not depend on T at constant n₀? (Ignoring Sutherland.)
\item What is "vacuum" in terms of ⟨l⟩?
\end{enumerate}

\chapter*{23. Дифузія в газах · 24. Внутрішнє тертя в газах · 25. Теплопровідність газів}\addcontentsline{toc}{chapter}{23. Дифузія в газах · 24. Внутрішнє тертя в газах · 25. Теплопровідність газів}\markboth{23. Дифузія в газах}{23. Дифузія в газах}


\srcline{\book{} т. 1, \textbf{§ 15.2} «Внутрішнє тертя (в'язкість) у газах» (с. 348–352), \textbf{§ 15.3} «Теплопровідність газів» (с. 352–356), \textbf{§ 15.4} «Дифузія газів» (с. 356–363). \ok{} (15.10), (15.14), (15.16)–(15.27) checked against the scan.}
Block C, page 2 of 2. One idea, three times. See the concept page явища переносу.

\needspace{5\baselineskip}\section*{Layout (best as a 3-column table with a shared top)}
\begin{layout}
┌ [A] Що таке явища переносу: перенос маси / імпульсу / енергії; спільна модель ┐
│     (площинка ΔS, 1/6 n₀⟨v⟩ в кожен бік, «пам'ять» на відстані ⟨l⟩) + рис. 15.3  │
├────────────────┬────────────────────┬──────────────────────────────────────┤
│ [B] Дифузія    │ [C] В'язкість      │ [D] Теплопровідність                   │
│ Фік (15.22)    │ Ньютон (15.10)     │ Фур'є (15.16)                          │
│ D = ⅓⟨v⟩⟨l⟩    │ η = ⅓ρ⟨v⟩⟨l⟩       │ λ = ⅓ρ⟨v⟩⟨l⟩c_V                        │
├────────────────┴────────────────────┴──────────────────────────────────────┤
└ [E] Зв'язок: η = Dρ, λ = ηc_V; залежність від p і T                          ┘
\end{layout}

\needspace{5\baselineskip}\section*{[A] Явища переносу (pp. 348, 349, 353, 357)}
\begin{itemize}
\item The book's definition (p. 348): phenomena caused by thermal motion that carry \textbf{mass (diffusion), momentum (internal friction), energy (heat conduction)}. They drive the system towards equilibrium.
\item A shared derivation:
\begin{itemize}
\item Take a small area ΔS ⟂ OX with a gradient along OX.
\item In each direction, $\tfrac16n_0\langle v\rangle$ molecules cross per unit time.
\item Each carries the value of the quantity it had at its \textbf{last collision}, a distance ⟨l⟩ away.
\item The net flux ∝ the difference between x − ⟨l⟩ and x + ⟨l⟩, which is the gradient × 2⟨l⟩.
\end{itemize}
\item \pen{} One shared picture (рис. 15.3/15.5/15.7): the plane ΔS, ⟨l⟩ to the left and right, and the quantity "higher" on the left.
\end{itemize}

\needspace{5\baselineskip}\section*{[B] Дифузія (23) (pp. 356–358)}
\begin{itemize}
\item «Дифузія газів — це процес самовільного вирівнювання густин (концентрацій) у суміші, зумовлений тепловим рухом мікрочастинок». In a single gas it is called \textbf{самодифузія} (self-diffusion). It is \textbf{irreversible}.
\item \textbf{Закон Фіка} (the first): $j_m = -D\dfrac{d\rho}{dx}$ (15.22), where D is the \textbf{коефіцієнт дифузії}, in м²/с.
\item From kinetic theory: $D = \tfrac13\langle v\rangle\langle l\rangle$ (15.27)
\item D ∝ 1/p at T = const, and D grows with T.
\item Optional: unsteady diffusion, $\Delta n = \Delta n_0e^{-t/\tau}$ (15.29); Fick's second law $\partial n/\partial t = D\,\partial^2n/\partial x^2$ (15.30); mutual diffusion and thermodiffusion (pp. 360–363).
\end{itemize}

\needspace{5\baselineskip}\section*{[C] Внутрішнє тертя (в'язкість) (24) (pp. 348–350)}
\begin{itemize}
\item Layers moving at different speeds exchange the \textbf{momentum of ordered motion} carried by thermal motion. The fast layer slows and the slow layer speeds up.
\item \textbf{Закон Ньютона}: $F = -\eta\dfrac{du}{dx}\Delta S$ (15.10), where η is the \textbf{коефіцієнт динамічної в'язкості}, in Па·с. The \textbf{кінематична в'язкість} is ν = η/ρ.
\item From kinetic theory: $\eta = \tfrac13\rho\langle v\rangle\langle l\rangle$ (15.14)
\item \spark{} \textbf{The paradox} (p. 350): η \textbf{does not depend on pressure}, because ρ ∝ p and ⟨l⟩ ∝ 1/p. More molecules carry momentum, but each over a shorter distance. Cutting p by 500 times changes the viscosity of air by only 4 \%.
\item η grows with T (∝ √T, in practice faster). Puluj measured this, p. 350.
\end{itemize}

\needspace{5\baselineskip}\section*{[D] Теплопровідність (25) (pp. 352–354)}
\begin{itemize}
\item «Теплопровідність — це молекулярне перенесення енергії в суцільному середовищі, зумовлене градієнтом температури»
\item \textbf{Закон Фур'є}: $j_q = -\lambda\,\text{grad}\,T$ (15.16), where λ is the \textbf{коефіцієнт теплопровідності}, in Вт/(м·К).
\item Each molecule carries its average energy $\tfrac i2kT$. From kinetic theory:
$\lambda = \tfrac13\rho\langle v\rangle\langle l\rangle c_V$ (15.21), where $c_{V}$ is the specific heat capacity.
\item Also \textbf{independent of p}, except at very low pressure. Light gases conduct better (H₂ is used for cooling).
\end{itemize}

\needspace{5\baselineskip}\section*{[E] Зв'язок між коефіцієнтами (p. 358)}
\[\eta = D\rho,\qquad \lambda = \eta c_V,\qquad \lambda = D\rho c_V\]
{\footnotesize
\begin{longtable}{>{\raggedright\arraybackslash}p{0.230\linewidth}>{\raggedright\arraybackslash}p{0.215\linewidth}>{\raggedright\arraybackslash}p{0.100\linewidth}>{\raggedright\arraybackslash}p{0.215\linewidth}>{\raggedright\arraybackslash}p{0.072\linewidth}>{\raggedright\arraybackslash}p{0.072\linewidth}}
\toprule
\textbf{} & \textbf{що переноситься} & \textbf{закон} & \textbf{коефіцієнт} & \textbf{від p} & \textbf{від T} \\
\midrule
\endhead
дифузія & маса & Фіка & D = ⅓⟨v⟩⟨l⟩ & ∝ 1/p & grows \\
в'язкість & імпульс & Ньютона & η = ⅓ρ⟨v⟩⟨l⟩ & none & grows \\
теплопровідність & енергія & Фур'є & λ = ⅓ρ⟨v⟩⟨l⟩$c_{V}$ & none & grows \\
\bottomrule
\end{longtable}
}

\needspace{5\baselineskip}\section*{Self-check}
\begin{enumerate}
\item Why does D depend on pressure but η and λ don't?
\item What exactly does a molecule "carry" in each of the three phenomena?
\item Where does the factor ⅓ (= 2 × ⅙) come from?
\end{enumerate}

\chapter*{3. Внутрішня енергія термодинамічної системи}\addcontentsline{toc}{chapter}{3. Внутрішня енергія термодинамічної системи}\markboth{3. Внутрішня енергія термодинамічної системи}{3. Внутрішня енергія термодинамічної системи}


\srcline{\book{} т. 1, \textbf{§ 16.1} «Основні поняття і означення» (с. 371–374), \textbf{§ 16.2} «Перший закон термодинаміки», the part «Внутрішня енергія» (с. 375–377). \ok{} Checked against the scan as part of block A (pp. 375–377).}
Block D, page 1 of 4. This page sets up the \textbf{thermodynamic language} that blocks D–F all use. The ideal-gas formula U = (i/2)RT is on 20–21; here we only link to it.

\needspace{5\baselineskip}\section*{Layout}
\begin{layout}
┌ [A] Термодинаміка: предмет; термодинамічна система; параметри стану   ┐
├ [B] Типи систем (ізольована/закрита/відкрита/адіабатна) + рівноважні процеси ┤
├ [C] Внутрішня енергія: означення, що входить / що НЕ входить              ┤
├ [D] Функція стану: ∮dU = 0 (малюнок двох шляхів A→B на p–V)              ┤
└ [E] Адитивність (16.1); калоричне рівняння стану; ідеальний газ → 20–21  ┘
\end{layout}

\needspace{5\baselineskip}\section*{[A] Основні поняття (pp. 371–372)}
\begin{itemize}
\item \textbf{Термодинаміка}: «наука про найбільш загальні властивості макроскопічних фізичних систем, що знаходяться в стані термодинамічної рівноваги, та про процеси переходу між цими станами». It is \textbf{axiomatic}: it rests on laws drawn from experience, not on molecular models (p. 374).
\item \textbf{Термодинамічна система} (thermodynamic system): macroscopic bodies and fields that can exchange \textbf{energy} (as work and heat) and \textbf{matter}.
\item \textbf{Термодинамічні параметри стану} (state parameters): p, T, V, ρ, concentration, field strengths… These are \textbf{внутрішні} (internal) and \textbf{зовнішні} (external), p. 372.
\end{itemize}

\needspace{5\baselineskip}\section*{[B] Системи і процеси (pp. 372–373)}
{\small
\begin{longtable}{>{\raggedright\arraybackslash}p{0.443\linewidth}>{\raggedright\arraybackslash}p{0.246\linewidth}>{\raggedright\arraybackslash}p{0.263\linewidth}}
\toprule
\textbf{система} & \textbf{обмін речовиною} & \textbf{обмін енергією} \\
\midrule
\endhead
ізольована (замкнута) & ні & ні \\
закрита & ні & так \\
відкрита & так & так \\
теплоізольована (адіабатна) & — & no heat exchange \\
\bottomrule
\end{longtable}
}
\begin{itemize}
\item \textbf{Квазістатичний (квазірівноважний) процес}: a continuous sequence of equilibrium states, such as sand added one grain at a time onto a piston. Only these can be drawn on a p–V diagram.
\item Optional: \textbf{фаза} (phase) and \textbf{компонента} (component), p. 373.
\end{itemize}

\needspace{5\baselineskip}\section*{[C] Внутрішня енергія (p. 375)}
\begin{itemize}
\item «Енергію, що залежить від внутрішнього стану фізичної системи, називають \textbf{внутрішньою енергією}»
\item It \textbf{includes} the kinetic energy of all kinds of particle motion plus their interaction energy. In practice that means the \textbf{kinetic energy of thermal motion + the potential energy of interaction} (no intra-atomic or nuclear energy).
\item It \textbf{excludes} the kinetic energy of the system as a whole and its potential energy in external fields; those are "external" energies.
\item Sign: dU > 0 if U increases.
\end{itemize}

\needspace{5\baselineskip}\section*{[D] Функція стану (pp. 375–376) — put it in a box}
\begin{itemize}
\item U is a \textbf{однозначна функція стану} (single-valued function of state): ΔU = U₂ − U₁ depends only on the end states, not the path. For a cycle:
\[\oint dU = 0\]
\item \pen{} Sketch two different paths A→B on a p–V diagram, both with the same ΔU. This sets up the contrast with Q and A (topic 4), which \textbf{do} depend on the path.
\end{itemize}

\needspace{5\baselineskip}\section*{[E] Ще дві властивості (p. 376)}
\begin{itemize}
\item \textbf{Адитивність} (additivity): $U = \sum U_i$ (16.1)
\item \textbf{Калоричне рівняння стану} (caloric equation of state): U = U(V, T), or U(p, V), or U(p, T). For liquids and gases U depends on T (kinetic part) \textbf{and V} (interaction).
\item \textbf{Ideal gas}: there's no interaction, so U = U(T) only. For a monatomic gas $U = \tfrac32RT$ (16.2); the general case is on 20–21.
\end{itemize}

\needspace{5\baselineskip}\section*{Self-check}
\begin{enumerate}
\item Why doesn't the kinetic energy of a flying ball count in its internal energy?
\item What does "function of state" mean for a closed cycle?
\item Why does a real gas's U depend on V but an ideal gas's doesn't?
\end{enumerate}

\chapter*{4. Теплота, робота, теплоємність}\addcontentsline{toc}{chapter}{4. Теплота, робота, теплоємність}\markboth{4. Теплота, робота, теплоємність}{4. Теплота, робота, теплоємність}


\srcline{\book{} т. 1, \textbf{§ 16.2}, the parts «Макроскопічна робота», «Кількість теплоти» (с. 377–381). \ok{} (16.5)–(16.9) checked against the scan.}
Block D, page 2 of 4. The key idea: \textbf{Q and A are functions of the process}; U is a function of state (topic 3).

\needspace{5\baselineskip}\section*{Layout}
\begin{layout}
┌ [A] Робота: поршень (рис. 16.1), d'A = p dV (16.5), A = ∫p dV (16.6)     ┐
├ [B] Робота = площа під кривою; два шляхи A→B (рис. 16.2) → різна A       ┤
├ [C] Кількість теплоти: теплообмін, «мікроскопічна робота»; знаки          ┤
├ [D] Таблиця: робота vs теплота (макро/мікро; обидві — функції процесу)    ┤
└ [E] Теплоємність: C₀ = d'Q/dT (16.7), питома c (16.8), молярна C = Mc (16.9) ┘
\end{layout}

\needspace{5\baselineskip}\section*{[A] Макроскопічна робота (pp. 377–378)}
\begin{itemize}
\item A gas under a piston of area S expands quasistatically by dx: $d\text{'}A = F\,dx = pS\,dx = p\,dV$ (16.5)
\item \textbf{Sign} (the book's convention): work done \textbf{by the gas} when it expands is \textbf{positive}; work done by external forces on the gas is negative.
\item $A = \int_{V_1}^{V_2}p\,dV$ (16.6)
\item The book's definition (p. 378): «робота в термодинаміці — це \textbf{макроскопічна міра зміни внутрішньої енергії} системи у квазістатичних процесах»
\end{itemize}

\needspace{5\baselineskip}\section*{[B] Функція процесу (pp. 378–379)}
\begin{itemize}
\item On a p–V diagram, \textbf{A = the area under the curve} of the process (рис. 16.2).
\item Two different paths A→B have different areas, so different work. Work is a \textbf{функція процесу, а не функція стану}.
\item Hence the notation \textbf{d'A, not dA}: it is not an exact differential.
\item In a cycle the work can be ≠ 0, which is what makes heat engines possible.
\item \pen{} The main drawing: two paths AmB and AnB with the areas under them shaded differently.
\end{itemize}

\needspace{5\baselineskip}\section*{[C] Кількість теплоти (p. 380)}
\begin{itemize}
\item \textbf{Теплообмін} (heat exchange): «Процес передачі енергії від більш нагрітого тіла до менш нагрітого \textbf{без виконання макроскопічної роботи}». The \textbf{кількість теплоти Q} is the measure of it.
\item The mechanism is collisions of molecules, sometimes called "microscopic work".
\item Signs: heat \textbf{given to} the system is +, heat taken from it is −. Unit: Дж.
\item Q is also a function of the process, so it's written \textbf{d'Q}.
\end{itemize}

\needspace{5\baselineskip}\section*{[D] Робота vs теплота (p. 380)}
{\small
\begin{longtable}{>{\raggedright\arraybackslash}p{0.187\linewidth}>{\raggedright\arraybackslash}p{0.400\linewidth}>{\raggedright\arraybackslash}p{0.365\linewidth}}
\toprule
\textbf{} & \textbf{робота A} & \textbf{теплота Q} \\
\midrule
\endhead
форма передачі енергії & \textbf{macrophysical}: energy of ordered motion & \textbf{microphysical}: many molecular processes \\
пов'язана зі зміною & external parameters (V) & none; heat exchange alone \\
функція & процесу (d'A) & процесу (d'Q) \\
спільне & both are measures of \textbf{energy transfer}, in Дж &  \\
\bottomrule
\end{longtable}
}

\needspace{5\baselineskip}\section*{[E] Теплоємність (pp. 380–381)}
\begin{itemize}
\item $C_0 = \dfrac{d\text{'}Q}{dT}$ (16.7), the heat needed to raise the temperature by 1 К.
\item \textbf{Питома теплоємність} (specific heat capacity) c, per unit mass: $C_0 = mc$ (16.8)
\item \textbf{Молярна теплоємність} (molar heat capacity) C, per mole: $C = Mc$ (16.9)
\item Notation: $C_{p}$ (isobaric), $C_{V}$ (isochoric), $C_{T}$ (isothermal), $C_{\text{ад}}$ (adiabatic). Since d'Q depends on the process, \textbf{heat capacity depends on the process too}. Values for an ideal gas are on 05.
\end{itemize}

\needspace{5\baselineskip}\section*{Self-check}
\begin{enumerate}
\item Why is it d'A and not dA?
\item Why does the work around a cycle not have to be zero, while ΔU does?
\item What is the physical difference between heat and work?
\end{enumerate}

\chapter*{6. Перший закон термодинаміки}\addcontentsline{toc}{chapter}{6. Перший закон термодинаміки}\markboth{6. Перший закон термодинаміки}{6. Перший закон термодинаміки}


\srcline{\book{} т. 1, \textbf{§ 16.2} «Перший закон термодинаміки» (с. 374, 381–383). \ok{} (16.10)–(16.17) checked against the scan.}
Block D, page 3 of 4. It combines topics 3 and 4: two ways to change U.

\needspace{5\baselineskip}\section*{Layout}
\begin{layout}
┌ [A] Звідки: механічна енергія не зберігається при терті → теплота = енергія ┐
├ [B] Формулювання + (16.10) dU = d'Q − d'A  ⚠️ знаки!                        ┤
├ [C] Зручна форма (16.11) d'Q = dU + d'A; інтегральна (16.12); цикл (16.13)  ┤
├ [D] Вічний двигун I роду неможливий                                         ┤
└ [E] Для ідеального газу: d'Q = C_V dT + p dV (16.17) → місток до 05, 07–10   ┘
\end{layout}

\needspace{5\baselineskip}\section*{[A] Звідки закон (p. 374)}
\begin{itemize}
\item Mechanical energy is conserved only without \textbf{dissipative forces}. With friction it "disappears" but heat appears. Counting the internal energy of thermal motion restores conservation.
\item «Перший закон термодинаміки — це \textbf{узагальнений закон збереження й перетворення енергії} для тих макроскопічних систем, для яких суттєву роль відіграє тепловий рух»
\end{itemize}

\needspace{5\baselineskip}\section*{[B] Формулювання (p. 381) — put it in a box}
\begin{itemize}
\item «у теплових процесах зміна внутрішньої енергії відбувається за рахунок приросту кількості теплоти і виконаної роботи»
\[dU = d\text{'}Q - d\text{'}A\qquad(16.10)\]
\item \warn{} \textbf{A flaw in the book.} Under (16.10) it says «d'A — робота, виконана \textbf{над} системою», but with the minus sign and its own convention on p. 378 (d'A = p dV > 0 on expansion), d'A must be the work done \textbf{by} the system. Write: \textbf{d'A is the work done by the system against external forces.} The adiabatic case (16.20), d'A = −dU ("the system does work by losing internal energy"), confirms this reading.
\end{itemize}

\needspace{5\baselineskip}\section*{[C] Інші форми (p. 382)}
\begin{itemize}
\item The heat form, convenient for heat engineering: $d\text{'}Q = dU + d\text{'}A$ (16.11). Words: the heat given \textbf{goes into} raising U and doing work.
\item Integral form: $Q = (U_2 - U_1) + A$ (16.12)
\item A cycle (U₂ = U₁): $Q = A$ (16.13), so all the net heat becomes work.
\item If d'Q = 0 the process is \textbf{адіабатний} (topic 10).
\end{itemize}

\needspace{5\baselineskip}\section*{[D] Вічний двигун першого роду (p. 382)}
\begin{itemize}
\item «неможливо побудувати машину, що діє періодично, яка виконувала б роботу без одержання енергії ззовні або виконувала б більшу роботу, ніж одержана ззовні енергія»
\end{itemize}

\needspace{5\baselineskip}\section*{[E] Ідеальний газ (pp. 382–383)}
\begin{itemize}
\item U = U(T, V) (16.14), but for an ideal gas (∂U/∂V)\_T = 0, so $dU = C_VdT$ (16.16), with $C_V = (\partial U/\partial T)_V$.
\item One mole:
\[d\text{'}Q = C_VdT + p\,dV\qquad(16.17)\]
This one line is the source of everything in blocks D–E: Mayer's equation (05), the isoprocesses (07–09) and the adiabat (10).
\end{itemize}

\needspace{5\baselineskip}\section*{Self-check}
\begin{enumerate}
\item State the first law in words without formulas.
\item In a cycle ΔU = 0; where does the heat go?
\item Why is a perpetual motion machine of the first kind impossible?
\end{enumerate}

\chapter*{5. Теплоємність ідеальних газів}\addcontentsline{toc}{chapter}{5. Теплоємність ідеальних газів}\markboth{5. Теплоємність ідеальних газів}{5. Теплоємність ідеальних газів}


\srcline{\book{} т. 1, \textbf{§ 16.2} (Mayer's equation, с. 382–383), \textbf{§ 16.5} «Теплоємність ідеального газу» (с. 390–397). \ok{} (16.16)–(16.19), (16.46), the $C_{V}$/$C_{p}$ table checked against the scan.}
Block D, page 4 of 4. The $C_{V}$ table itself is already on 20–21. Here the focus is \textbf{$C_{V}$ vs $C_{p}$}, \textbf{γ} and \textbf{where classical theory fails}.

\needspace{5\baselineskip}\section*{Layout}
\begin{layout}
┌ [A] C_V = dU/dT; C_p — більше; рівняння Майєра (16.19) у рамці        ┐
├ [B] Класична теорія: C_V = (i/2)R, C_p = ((i+2)/2)R, γ = (i+2)/i — таблиця ┤
├ [C] Графік C_V(T) для двохатомного газу (рис. 16.7): сходинки 3/2 → 5/2 → 7/2 ┤
└ [D] Чому: квантування; характеристичні температури T_оберт, T_колив      ┘
\end{layout}

\needspace{5\baselineskip}\section*{[A] Рівняння Майєра (pp. 382–383) — put it in a box}
\begin{itemize}
\item From the first law for one mole, $d\text{'}Q = C_VdT + p\,dV$ (16.17):
\begin{itemize}
\item V = const: dV = 0, so $C_V = \left(\dfrac{\partial U}{\partial T}\right)_V$ (16.16)
\item p = const: p dV = R dT, so $d\text{'}Q = C_VdT + R\,dT$ (16.18)
\[C_p = C_V + R\qquad(16.19)\]
the \textbf{рівняння Майєра} (Mayer's equation).
\end{itemize}
\item Meaning: at constant p, part of the heat goes into \textbf{work done by expansion}, not only into U. So $C_{p}$ > $C_{V}$, and by exactly R (the meaning of R from 01).
\end{itemize}

\needspace{5\baselineskip}\section*{[B] Класична теорія (pp. 390–392)}
\begin{itemize}
\item From U = (i/2)RT: $C_V = \dfrac i2R$, $C_p = \dfrac{i+2}2R$, $\gamma = \dfrac{C_p}{C_V} = \dfrac{i+2}i$
\end{itemize}
{\footnotesize
\begin{longtable}{>{\raggedright\arraybackslash}p{0.460\linewidth}>{\raggedright\arraybackslash}p{0.077\linewidth}>{\raggedright\arraybackslash}p{0.096\linewidth}>{\raggedright\arraybackslash}p{0.096\linewidth}>{\raggedright\arraybackslash}p{0.192\linewidth}}
\toprule
\textbf{газ} & \textbf{i} & \textbf{$C_{V}$} & \textbf{$C_{p}$} & \textbf{γ} \\
\midrule
\endhead
одноатомний & 3 & 3/2 R & 5/2 R & 5/3 ≈ 1,67 \\
двохатомний (жорсткий) & 5 & 5/2 R & 7/2 R & 7/5 = 1,4 \\
багатоатомний (жорсткий) & 6 & 3R & 4R & 4/3 ≈ 1,33 \\
\bottomrule
\end{longtable}
}
\begin{itemize}
\item γ is the \textbf{показник адіабати} (adiabatic index); it is needed for topic 10.
\item The general formula with vibrations: $U = (3 + r + 2s)\dfrac{RT}2$ (16.46).
\end{itemize}

\needspace{5\baselineskip}\section*{[C] Невдачі класичної теорії (pp. 392–394)}
\begin{enumerate}
\item Complex molecules: for C₂H₅OH the rigid model gives 3R and the elastic one 24R, but experiment gives about 8R. \textbf{Both models fail.}
\item \textbf{The temperature dependence} isn't explained. For a diatomic gas, рис. 16.7:
\begin{itemize}
\item low T: $C_{V}$ ≈ 3/2 R, as if there were translation only (H₂ below 60 К behaves as monatomic);
\item room T: 5/2 R (rotation "switched on");
\item thousands of К: towards 7/2 R (vibrations switch on, but molecules dissociate first).
\end{itemize}
\item \pen{} Draw the "staircase" $C_{V}$(T) with steps at $T_{\text{оберт}}$ and $T_{\text{колив}.}$
\end{enumerate}

\needspace{5\baselineskip}\section*{[D] Квантове пояснення (pp. 394–397), short}
\begin{itemize}
\item The energies of rotation and vibration are \textbf{quantized}: $\varepsilon_{\text{пост}} \ll \varepsilon_{\text{оберт}} \ll \varepsilon_{\text{колив}}$, hence the \textbf{характеристичні температури} (characteristic temperatures) $T_{\text{пост}} \ll T_{\text{оберт}} \ll T_{\text{колив}}$.
\item Below $T_{\text{оберт}}$ or $T_{\text{колив}}$, kT is too small to excite that motion, so its degrees of freedom are "frozen".
\item Values: $T_{\text{оберт}}$ for H₂ is 85 К, for N₂ 2,85 К; $T_{\text{колив}}$ for N₂ is 3340 К, for O₂ 2240 К.
\end{itemize}

\needspace{5\baselineskip}\section*{Self-check}
\begin{enumerate}
\item Why is $C_{p}$ − $C_{V}$ exactly R?
\item Why does diatomic H₂ behave as monatomic at 60 К?
\item Which γ do you use for air, and why?
\end{enumerate}

\chapter*{7. Ізотермічний процес · 8. Ізохоричний процес · 9. Ізобарний процес}\addcontentsline{toc}{chapter}{7. Ізотермічний процес · 8. Ізохоричний процес · 9. Ізобарний процес}\markboth{7. Ізотермічний процес}{7. Ізотермічний процес}


\srcline{\book{} т. 1, \textbf{§ 16.4} «Робота ідеального газу при ізопроцесах. Політропний процес» (с. 386–390); the laws themselves are in \textbf{§ 14.7} (с. 313–316, see 01). \ok{} (16.26)–(16.33) checked against the scan.}
Block E, page 1 of 2. Spelling: the syllabus says «ізохоричний», the book «\textbf{ізохорний}».
Use with the summary table isoprocesses, which is the most useful single page for the МКР.

\needspace{5\baselineskip}\section*{Layout (one block per process, the same 5 rows each time)}
\begin{layout}
┌ [Ізохорний V = const]  ┬ [Ізобарний p = const] ┬ [Ізотермічний T = const] ┐
│ закон: p/T = const     │ V/T = const (Г-Л)     │ pV = const (Б-М)          │
│ графік p–V: вертикаль  │ горизонталь + площа   │ гіпербола + площа         │
│ A = 0                  │ A = p(V₂−V₁)=νRΔT     │ A = νRT ln(V₂/V₁)          │
│ ΔU = νC_VΔT            │ ΔU = νC_VΔT           │ ΔU = 0                     │
│ Q = ΔU                 │ Q = νC_pΔT            │ Q = A                      │
└────────────────────────┴───────────────────────┴────────────────────────────┘
          + footer: політропний процес pVⁿ = const об'єднує всі (16.40)
\end{layout}
(ν = m/M)

\needspace{5\baselineskip}\section*{Ізохорний процес, V = const (pp. 386–387)}
\begin{itemize}
\item Example: heating a gas in a rigid, sealed vessel.
\item $d\text{'}A_V = p\,dV = 0$, $A_V = 0$ (16.26)
\item From the first law, $d\text{'}Q = dU = \dfrac mMC_VdT$ (16.27) and $\Delta U = \dfrac mMC_V(T_2 - T_1)$ (16.28).
\item \textbf{All the heat goes into internal energy.}
\end{itemize}

\needspace{5\baselineskip}\section*{Ізобарний процес, p = const (p. 387)}
\begin{itemize}
\item Example: a gas under a freely moving piston, or water heated in an open vessel.
\item $A_p = p(V_2 - V_1) = \dfrac mMR(T_2 - T_1)$ (16.29). On the p–V diagram this is a \textbf{rectangle} (рис. 16.4).
\item $Q_p = \dfrac mM(C_V + R)(T_2 - T_1) = \dfrac mMC_p(T_2 - T_1)$ (16.30)
\item The heat goes into \textbf{both ΔU and work} (Mayer).
\end{itemize}

\needspace{5\baselineskip}\section*{Ізотермічний процес, T = const (pp. 387–388)}
\begin{itemize}
\item It needs ideal thermal contact with a \textbf{термостат} (a body at constant temperature) and must be very slow. All first-order phase transitions (boiling, for example) are isothermal.
\item dU = 0, so $d\text{'}Q = d\text{'}A$ (16.31): \textbf{all the heat becomes work}.
\item \[A_T = \frac mMRT\ln\frac{V_2}{V_1}\qquad(16.32)\]
\item Also: $Q_T = A_T = \dfrac mMRT\ln\dfrac{p_1}{p_2} = p_1V_1\ln\dfrac{V_2}{V_1}$ (16.33)
\item \pen{} The area under the hyperbola. Remember where the logarithm comes from: ∫dV/V.
\end{itemize}

\needspace{5\baselineskip}\section*{Footer: політропний процес (pp. 388–389), optional but elegant}
\begin{itemize}
\item A process with \textbf{constant heat capacity C}: $pV^n = \text{const}$ (16.40), with $n = \dfrac{C_p - C}{C_V - C}$ (16.37).
\end{itemize}
{\small
\begin{longtable}{>{\raggedright\arraybackslash}p{0.571\linewidth}>{\raggedright\arraybackslash}p{0.190\linewidth}>{\raggedright\arraybackslash}p{0.190\linewidth}}
\toprule
\textbf{процес} & \textbf{C} & \textbf{n} \\
\midrule
\endhead
ізобарний & $C_{p}$ & 0 \\
ізотермічний & ∞ & 1 \\
адіабатний & 0 & γ \\
ізохорний & $C_{V}$ & ±∞ \\
\bottomrule
\end{longtable}
}
\begin{itemize}
\item All four processes are special cases of one equation.
\end{itemize}

\needspace{5\baselineskip}\section*{Self-check}
\begin{enumerate}
\item In which process does all the heat go into work, and in which all into U?
\item Why is $C_{T}$ = ∞? (Heat is added but ΔT = 0.)
\item Draw all three on one p–V diagram from the same starting point.
\end{enumerate}

\chapter*{10. Адіабатичний процес}\addcontentsline{toc}{chapter}{10. Адіабатичний процес}\markboth{10. Адіабатичний процес}{10. Адіабатичний процес}


\srcline{\book{} т. 1, \textbf{§ 16.3} «Адіабатний процес. Рівняння Пуассона» (с. 383–386). \ok{} (16.20)–(16.25) checked against the scan.}
Block E, page 2 of 2. Terminology: the syllabus says «адіабатичний», the book «\textbf{адіабатний}». Use the book's word.

\needspace{5\baselineskip}\section*{Layout}
\begin{layout}
┌ [A] Означення + умови (теплоізоляція, «достатня повільність»)        ┐
├ [B] Перший закон: dU + d'A = 0 (16.20) → pdV = −C_V dT (16.21)        ┤
├ [C] Виведення → рівняння Пуассона TV^(γ−1) = const, pV^γ = const      ┤
├ [D] ГРАФІК: адіабата крутіша за ізотерму (рис. 16.3) + чому           ┤
└ [E] Робота (16.24)–(16.25); приклади з життя                          ┘
\end{layout}

\needspace{5\baselineskip}\section*{[A] Означення (pp. 383–384)}
\begin{itemize}
\item \textbf{Адіабатний процес}: «що відбувається \textbf{без теплообміну} з навколишнім середовищем» (d'Q = 0).
\item Two conditions: \textbf{heat insulation} (non-conducting walls) and \textbf{enough slowness} to stay quasistatic. In practice the process is made fast, so heat has no time to leave, but slower than the speed of sound.
\end{itemize}

\needspace{5\baselineskip}\section*{[B] Перший закон (p. 384)}
\begin{itemize}
\item $dU + d\text{'}A = 0$ (16.20) ⟹ $d\text{'}A = -dU$. «система виконує роботу внаслідок \textbf{зменшення внутрішньої енергії}»
\item One mole: $p\,dV = -C_VdT$ (16.21). \textbf{Expansion cools} the gas; compression heats it.
\end{itemize}

\needspace{5\baselineskip}\section*{[C] Рівняння Пуассона (p. 384) — put it in a box}
\begin{itemize}
\item Substitute p = RT/V: $\dfrac{dT}T = -\dfrac R{C_V}\dfrac{dV}V$ (16.22). Integrate, using R/$C_{V}$ = γ − 1:
\[TV^{\gamma-1} = \text{const}\ (16.23),\qquad pV^\gamma = \text{const},\quad pT^{\frac\gamma{1-\gamma}} = \text{const}\ (16.23a)\]
\item \textbf{γ = $C_{p}$/$C_{V}$}, the \textbf{показник адіабати} or \textbf{коефіцієнт Пуассона} (values on 05).
\end{itemize}

\needspace{5\baselineskip}\section*{[D] Адіабата vs ізотерма (pp. 384–385)}
\begin{itemize}
\item Isotherm: p ∝ 1/V. Adiabat: p ∝ 1/$V^{γ}$ with γ > 1, so the \textbf{adiabat is steeper}.
\item Physically: during adiabatic expansion p falls both because \textbf{V grows and because T falls}.
\item \pen{} The main drawing, рис. 16.3: an isotherm and an adiabat through one point, the adiabat steeper, with the work area V₁→V₂ shaded.
\end{itemize}

\needspace{5\baselineskip}\section*{[E] Робота (p. 385)}
\begin{itemize}
\item $A = -\dfrac mMC_V\int_{T_1}^{T_2}dT = \dfrac mMC_V(T_1 - T_2)$ (16.24)
\item \[A = \frac mM\frac{RT_1}{\gamma - 1}\left[1 - \left(\frac{V_1}{V_2}\right)^{\gamma-1}\right]\qquad(16.25)\]
\item Examples from life:
\begin{itemize}
\item air rising in the troposphere expands adiabatically and cools, which is why it is colder higher up;
\item internal combustion engines;
\item producing low temperatures.
\end{itemize}
\end{itemize}

\needspace{5\baselineskip}\section*{Self-check}
\begin{enumerate}
\item Why does a gas cool when it expands adiabatically, when no heat is taken away?
\item Why is the adiabat steeper than the isotherm? Give two reasons in one sentence.
\item Air: γ = 1,4. How does T change if the volume halves? (×2\textasciicircum{}0,4 ≈ 1,32.)
\end{enumerate}

\chapter*{11. Колові процеси}\addcontentsline{toc}{chapter}{11. Колові процеси}\markboth{11. Колові процеси}{11. Колові процеси}


\srcline{\book{} т. 1, \textbf{§ 16.6} «Оборотні та необоротні процеси. Колові процеси (цикли)» (с. 397–399). \ok{} Checked against the scan (pp. 397–399; (16.13) on p. 382).}
Block F, page 1 of 5.

\needspace{5\baselineskip}\section*{Layout}
\begin{layout}
┌ [A] Оборотні vs необоротні процеси (таблиця + приклади)          ┐
├ [B] Цикл: означення; ΔU = 0 → Q = A (16.13)                       ┤
├ [C] ГРАФІК: цикл 1a2b1 на p–V (рис. 16.8): A = площа всередині     ┤
└ [D] Прямий (теплова машина) vs обернений (холодильна) цикл        ┘
\end{layout}

\needspace{5\baselineskip}\section*{[A] Оборотні та необоротні процеси (p. 397)}
\begin{itemize}
\item \textbf{Оборотний процес} (reversible): «такий процес, при якому система, проведена в зворотному напрямі, повертається в початковий стан \textbf{через ті самі проміжні стани}, що й у прямому процесі, а стан тіл поза системою залишається незмінним». All quasistatic processes are reversible.
\item \textbf{Необоротний процес} (irreversible): one that can't be reversed without changes in the system or its surroundings. \textbf{All real processes are irreversible.}
\item Typical irreversible processes: heat exchange across a finite ΔT, diffusion, expansion into vacuum, friction. The main sign of irreversibility is \textbf{dissipation heat} (p. 398).
\item Why study ideal reversible processes: they set the limit that real processes approach.
\end{itemize}

\needspace{5\baselineskip}\section*{[B] Цикл (p. 398)}
\begin{itemize}
\item \textbf{Цикл (коловий процес)}: «термодинамічний процес, у результаті якого система повертається в початковий стан»
\item U is a function of state, so ΔU = 0 over a cycle, and $Q = A$ (16.13): the net heat received becomes work (06).
\item Why cycles: one expansion gives only limited work. To repeat it you must compress the gas back \textbf{along a different path}, at lower pressure, otherwise the net work is 0.
\end{itemize}

\needspace{5\baselineskip}\section*{[C] Робота циклу (pp. 398–399)}
\begin{itemize}
\item The work of a cycle equals the \textbf{area enclosed by the loop} on the p–V diagram (рис. 16.8): the expansion work minus the compression work.
\item \pen{} The main drawing: a loop 1a2b1, horizontal hatching under the upper branch (expansion), vertical hatching under the lower one (compression), and the difference shaded.
\end{itemize}

\needspace{5\baselineskip}\section*{[D] Прямий і обернений цикл (p. 399)}
{\small
\begin{longtable}{>{\raggedright\arraybackslash}p{0.114\linewidth}>{\raggedright\arraybackslash}p{0.533\linewidth}>{\raggedright\arraybackslash}p{0.305\linewidth}}
\toprule
\textbf{} & \textbf{\textbf{прямий}} & \textbf{\textbf{обернений}} \\
\midrule
\endhead
машина & теплова машина (двигун) & холодильна машина \\
обхід на p–V & clockwise & anticlockwise \\
робота & $A_{\text{розш}}$ > $A_{\text{стиск}}$, A > 0 & $A_{\text{стиск}}$ > $A_{\text{розш}}$ \\
теплота & part of the heat → work, the rest → холодильник (cooler) & work moves heat from cold to hot \\
\bottomrule
\end{longtable}
}
\begin{itemize}
\item \textbf{Оборотний цикл} (reversible cycle): all its parts are reversible. \textbf{Необоротний} (irreversible): at least one part is not.
\end{itemize}

\needspace{5\baselineskip}\section*{Self-check}
\begin{enumerate}
\item Why can't you compress the gas back along the same curve?
\item Which way round does a heat engine's cycle go on a p–V diagram?
\item Why is every real process irreversible?
\end{enumerate}

\chapter*{12. Цикл Карно (необернений) та його ккд · 13. Цикл Карно (обернений) та його ккд}\addcontentsline{toc}{chapter}{12. Цикл Карно (необернений) та його ккд · 13. Цикл Карно (обернений) та його ккд}\markboth{12. Цикл Карно (необернений) та його ккд}{12. Цикл Карно (необернений) та його ккд}


\srcline{\book{} т. 1, \textbf{§ 16.6} «… Цикл Карно та його коефіцієнт корисної дії» (с. 399–401), \textbf{§ 16.7} «Теплові машини» (с. 402–403). \ok{} (16.53)–(16.68) checked against the scan.}
Block F, page 2 of 5. Terminology: the syllabus says «\textbf{необернений}» meaning the \textbf{direct} cycle. The book never uses this word; it says \textbf{прямий} / \textbf{обернений} (direct / reverse). Write «прямий (необернений)» once on the page.

\needspace{5\baselineskip}\section*{Layout}
\begin{layout}
┌ [A] Що таке цикл Карно: 2 ізотерми + 2 адіабати; ідеальна машина     ┐
├ [B] ГРАФІК 1→2→3→4→1 (рис. 16.9) з T₁, T₂, Q₁, Q₂ + таблиця 4 стадій ┤
├ [C] ККД прямого циклу: η = 1 − Q₂/Q₁ (16.53) = 1 − T₂/T₁ (16.64) — рамка ┤
├ [D] Обернений цикл: холодильна машина ε (16.68), тепловий насос φ (16.67) ┤
└ [E] Схеми трьох машин (рис. 16.10–16.12) — стрілки потоків енергії     ┘
\end{layout}

\needspace{5\baselineskip}\section*{[A] Цикл Карно (p. 399)}
\begin{itemize}
\item Carnot's 1824 questions: is there a limit to improving heat engines, and does the working substance matter?
\item «Він складається з \textbf{двох адіабатних та двох ізотермічних} процесів». It is \textbf{reversible}: the cycle of an \textbf{ideal} heat engine, with no losses to conduction, radiation or friction.
\item It is also the route to a quantitative \textbf{second law} (16).
\end{itemize}

\needspace{5\baselineskip}\section*{[B] Прямий цикл з ідеальним газом (pp. 400–401), 1 mole, clockwise}
{\footnotesize
\begin{longtable}{>{\raggedright\arraybackslash}p{0.069\linewidth}>{\raggedright\arraybackslash}p{0.193\linewidth}>{\raggedright\arraybackslash}p{0.183\linewidth}>{\raggedright\arraybackslash}p{0.362\linewidth}>{\raggedright\arraybackslash}p{0.113\linewidth}}
\toprule
\textbf{стадія} & \textbf{процес} & \textbf{контакт} & \textbf{робота} & \textbf{теплота} \\
\midrule
\endhead
1→2 & ізотермічне розширення, T₁ & \textbf{нагрівник} (heater) & $A_{12} = RT_1\ln\frac{V_2}{V_1} = Q_1$ (16.55) & Q₁ > 0 received \\
2→3 & адіабатне розширення & isolated & $A_{23} = C_V(T_1 - T_2)$ (16.56) & 0 \\
3→4 & ізотермічне стискання, T₂ & \textbf{холодильник} (cooler) & $A_{34} = -RT_2\ln\frac{V_3}{V_4} = -Q_2$ (16.57) & Q₂ given away \\
4→1 & адіабатне стискання & isolated & $A_{41} = C_V(T_2 - T_1)$ (16.58) & 0 \\
\bottomrule
\end{longtable}
}
\begin{itemize}
\item The adiabatic works cancel. From the adiabats: $\dfrac{V_2}{V_1} = \dfrac{V_3}{V_4}$ (16.63), the \textbf{condition for the cycle to close}.
\item \pen{} The main drawing, рис. 16.9: two isotherms (T₁ upper, T₂ lower) and two steeper adiabats, with arrows 1→2→3→4 and Q₁ in, Q₂ out.
\end{itemize}

\needspace{5\baselineskip}\section*{[C] ККД (pp. 400–401) — put it in a box}
\[\eta = \frac A{Q_1} = \frac{Q_1 - Q_2}{Q_1} = 1 - \frac{Q_2}{Q_1}\ (16.53)\qquad \eta = \frac{T_1 - T_2}{T_1} = 1 - \frac{T_2}{T_1}\ (16.64)\]
\begin{itemize}
\item The efficiency of a reversible Carnot cycle depends \textbf{only on the temperatures of the heater and the cooler}.
\item \textbf{η < 1 always}: η = 1 would need T₁ → ∞ or T₂ → 0, which is impossible.
\item If ΔT → 0 then η → 0, so a heat engine \textbf{must have two thermostats}.
\end{itemize}

\needspace{5\baselineskip}\section*{[D] Обернений цикл (pp. 399, 402–403)}
\begin{itemize}
\item It runs anticlockwise: the work of \textbf{external forces} carries heat from the \textbf{colder} body to the \textbf{hotter} one.
\item \textbf{Холодильна машина} (refrigerator): $\varepsilon = \dfrac{Q_2}{A}$ (16.68), the \textbf{холодильний коефіцієнт}, heat taken from the cold body per unit of work.
\begin{itemize}
\item For a reverse Carnot cycle it follows from (16.64) that $\varepsilon = \dfrac{T_2}{T_1 - T_2}$. \warn{} This is derived here; the book doesn't write it out.
\end{itemize}
\item \textbf{Тепловий насос} (heat pump): $\varphi = \dfrac{Q_1}{A} = \dfrac{T_1}{T_1 - T_2}$ (16.67), the \textbf{коефіцієнт перетворення}, heat delivered to the warm body per unit of work. It can be > 1.
\end{itemize}

\needspace{5\baselineskip}\section*{[E] Три типи теплових машин (p. 402)}
\begin{itemize}
\item \pen{} Three small block diagrams (рис. 16.10–16.12):
\begin{itemize}
\item двигун (engine): heater → Q₁ → working body → A out, and Q₂ → cooler;
\item холодильник (refrigerator): A in, cold body → Q₂ → working body → Q₁ → surroundings;
\item тепловий насос (heat pump): the same arrows, but the goal is Q₁ going into the warm room.
\end{itemize}
\end{itemize}

\needspace{5\baselineskip}\section*{Self-check}
\begin{enumerate}
\item Why doesn't the Carnot efficiency depend on the gas?
\item Where does the work of the two adiabats go?
\item Why can a heat pump's φ be greater than 1 without breaking energy conservation?
\end{enumerate}

\chapter*{16. Другий закон термодинаміки}\addcontentsline{toc}{chapter}{16. Другий закон термодинаміки}\markboth{16. Другий закон термодинаміки}{16. Другий закон термодинаміки}


\srcline{\book{} т. 1, \textbf{§ 16.8} «Другий закон термодинаміки. Теорема Карно. Поняття про ентропію. Характеристичні функції» (с. 403–409). \ok{} (16.69)–(16.74) checked against the scan.}
Block F, page 3 of 5. Entropy itself (the same §, from p. 406) is on 14.

\needspace{5\baselineskip}\section*{Layout}
\begin{layout}
┌ [A] Що визначає закон: НАПРЯМ процесів + межу перетворення Q → A      ┐
├ [B] Три формулювання: Клаузіус · Томсон · Оствальд (у рамках)          ┤
├ [C] Теорема Карно: два твердження; η ≤ η₀ (16.69)                     ┤
├ [D] Кількісно: нерівність Клаузіуса Q₁/T₁ + Q₂/T₂ ≤ 0 (16.71); зведена теплота ┤
└ [E] Узагальнення на довільний цикл ∮d'Q/T ≤ 0 (16.74) + рис. 16.13     ┘
\end{layout}

\needspace{5\baselineskip}\section*{[A] Що визначає (p. 403)}
\begin{itemize}
\item The first law allows any process that conserves energy. The \textbf{second law} «визначає \textbf{напрям} термодинамічних процесів та встановлює максимально можливі границі перетворення теплоти в роботу». It sums up experiment.
\end{itemize}

\needspace{5\baselineskip}\section*{[B] Формулювання (pp. 403–404)}
\begin{itemize}
\item \textbf{Клаузіус (1850)}: «теплота не може \textbf{сама собою} переходити від менш нагрітого тіла до більш нагрітого»
\item \textbf{Томсон (Кельвін, 1851)}: «не можна побудувати машину, що діє періодично, яка б неперервно перетворювала теплоту в роботу \textbf{тільки внаслідок охолодження одного тіла}, без того, щоб у навколишніх тілах не відбулося одночасно певних залишкових змін»
\item \textbf{Оствальд}: «\textbf{вічний двигун другого роду} неможливий». Such an engine would run by cooling the ocean; it wouldn't break energy conservation, but the second law forbids it.
\item \pen{} Contrast it with the perpetual motion machine of the \textbf{first} kind (06), which is forbidden by energy conservation.
\end{itemize}

\needspace{5\baselineskip}\section*{[C] Теорема Карно (p. 404)}
\begin{enumerate}
\item The efficiency of a \textbf{reversible} Carnot cycle doesn't depend on the working substance or the design; it equals that of the ideal-gas cycle at the same temperatures.
\item The efficiency of an \textbf{irreversible} cycle is always lower.
\item $\eta \le \eta_0$ (16.69): equality for a reversible cycle, < for an irreversible one.
\end{enumerate}

\needspace{5\baselineskip}\section*{[D] Нерівність Клаузіуса (p. 404)}
\begin{itemize}
\item Put in $\eta = \dfrac{Q_1 - |Q_2|}{Q_1}$ and $\eta_0 = \dfrac{T_1 - T_2}{T_1}$ (16.70). With Q₂ < 0:
\[\frac{Q_1}{T_1} + \frac{Q_2}{T_2} \le 0\qquad(16.71)\]
\item \textbf{Зведена теплота} (reduced heat) is the heat exchanged with a reservoir divided by that reservoir's absolute temperature.
\item In words: the sum of reduced heats in a reversible Carnot cycle is 0; in an irreversible one it is < 0.
\end{itemize}

\needspace{5\baselineskip}\section*{[E] Довільний цикл (p. 405)}
\begin{itemize}
\item Any reversible cycle can be cut by a grid of isotherms and adiabats into elementary Carnot cycles (рис. 16.13); the inner pieces cancel.
\item \[\oint\frac{d\text{'}Q}T = 0\ \text{(\text{оборотний})}\ (16.73),\qquad \oint\frac{d\text{'}Q}T \le 0\ \text{(\text{загалом})}\ (16.74)\]
\item (16.74) is the second law for any cycle. Because (16.73) equals 0, d'Q/T must be the differential of a \textbf{function of state}, which is \textbf{entropy}. Continue on 14.
\item \pen{} Sketch рис. 16.13: an oval cut by a grid of isotherms and adiabats.
\end{itemize}

\needspace{5\baselineskip}\section*{Self-check}
\begin{enumerate}
\item Clausius's and Thomson's wordings sound different. Why are they equivalent?
\item Why doesn't a perpetual motion machine of the second kind break the first law?
\item What is "reduced heat", and why is it summed and not plain heat?
\end{enumerate}

\chapter*{14. Ентропія та її властивості}\addcontentsline{toc}{chapter}{14. Ентропія та її властивості}\markboth{14. Ентропія та її властивості}{14. Ентропія та її властивості}


\srcline{\book{} т. 1, \textbf{§ 16.8} (from p. 406, «Поняття про ентропію»), \textbf{§ 16.10} «Зв'язок ентропії з імовірністю стану системи» (с. 411–414); Example 16.4 (с. 410–411). \ok{} (16.75)–(16.82) checked against the scan; (16.84), (16.91) from clear OCR.}
Block F, page 4 of 5. Two views: \textbf{thermodynamic} (dS = d'Q/T) and \textbf{statistical} (S = k ln w).

\needspace{5\baselineskip}\section*{Layout}
\begin{layout}
┌ [A] Означення: dS = d'Q/T (16.75) — рамка; Клаузіус 1854          ┐
├ [B] Властивості: функція стану · адитивна · до сталої (16.76)      ┤
├ [C] Закон зростання ентропії: dS ≥ d'Q/T (16.81), ізольована dS ≥ 0 (16.82) ┤
├ [D] ΔS для процесів ідеального газу (приклад 16.4) — табличка     ┤
├ [E] Статистичний зміст: S = k ln w (16.91); мікро-/макростан; рис. 16.15 ┤
└ footer: основне рівняння TdS = dU + pdV (16.84); «теплова смерть» — чому ні ┘
\end{layout}

\needspace{5\baselineskip}\section*{[A] Означення (p. 406) — put it in a box}
\[dS = \frac{d\text{'}Q}T\qquad(16.75)\]
\begin{itemize}
\item From (16.73): d'Q/T is an exact differential, so a function of state S exists. Clausius introduced it in 1854; the name is Greek for "turning, change".
\item The book's definition: «Ентропія — функція стану системи, яка визначається тим, що її повний диференціал dS при елементарному \textbf{оборотному} процесі дорівнює відношенню нескінченно малої кількості теплоти d'Q, що її дістала система, до термодинамічної температури T системи»
\item $d\text{'}Q = T\,dS$, the same structure as $d\text{'}A = p\,dV$: T is the "generalized force" of heat exchange and S its "generalized coordinate".
\end{itemize}

\needspace{5\baselineskip}\section*{[B] Властивості (p. 406)}
\begin{enumerate}
\item A \textbf{function of state}: it doesn't depend on the path.
\item \textbf{Additive}: S = ΣSᵢ.
\item Defined only \textbf{up to a constant}: $S = S_0 + \int\dfrac{d\text{'}Q}T$ (16.76). The third law fixes S₀ (15).
\end{enumerate}

\needspace{5\baselineskip}\section*{[C] Закон зростання ентропії (pp. 406–408)}
\begin{itemize}
\item An irreversible path A→B, returned reversibly (рис. 16.14), gives $S_2 - S_1 > \int_1^2\dfrac{d\text{'}Q}T$ (16.80). In general:
\[dS \ge \frac{d\text{'}Q}T\qquad(16.81)\]
\item An \textbf{isolated} system (d'Q = 0): $dS \ge 0$ (16.82). Equality holds for reversible processes in equilibrium (16.83).
\item «ентропія замкненої системи \textbf{не може зменшуватись}; вона або зростає, якщо в системі відбуваються необоротні процеси, або залишається незмінною …». Entropy is a \textbf{measure of irreversibility}.
\end{itemize}

\needspace{5\baselineskip}\section*{[D] ΔS для ідеального газу, 1 моль (pp. 410–411, Example 16.4)}
{\small
\begin{longtable}{>{\raggedright\arraybackslash}p{0.380\linewidth}>{\raggedright\arraybackslash}p{0.588\linewidth}}
\toprule
\textbf{процес} & \textbf{ΔS} \\
\midrule
\endhead
адіабатний (оборотний) & 0, \textbf{ізоентропійний} (isentropic) \\
ізотермічний & Q/T = R ln(V₂/V₁) \\
ізохорний & $C_{V}$ ln(T₂/T₁) \\
ізобарний & $C_{p}$ ln(T₂/T₁) \\
\bottomrule
\end{longtable}
}
\begin{itemize}
\item In general (Example 16.5): $\Delta S = \nu C_V\ln\dfrac{T_2}{T_1} + \nu R\ln\dfrac{V_2}{V_1}$
\end{itemize}

\needspace{5\baselineskip}\section*{[E] Статистичний зміст (pp. 411–413)}
\begin{itemize}
\item \textbf{Макростан} (macrostate): given p, T, ρ… \textbf{Мікростан} (microstate): a particular distribution of the particles over coordinates and momenta.
\item \textbf{Термодинамічна ймовірність w} (thermodynamic probability): the \textbf{number of microstates} that realize a given macrostate. It is ≥ 1, usually a huge number, unlike mathematical probability. $w = \dfrac{n!}{n_1!\,n_2!\cdots n_N!}$ (16.90)
\item The example on p. 412: 6 molecules in 3 cells. 2-2-2 can happen in 90 ways, 6-0-0 in only 1. The \textbf{even} distribution is the most probable, which is equilibrium.
\item \textbf{Принцип Больцмана} (1886):
\[S = k\ln w\qquad(16.91)\]
\item Meaning: «довільний природний процес протікає в напрямі \textbf{збільшення ймовірності стану}». Entropy measures \textbf{disorder} and how close the system is to equilibrium. A decrease is not forbidden, only very improbable (fluctuations).
\item \pen{} Sketch рис. 16.15: two molecules in two cells, with all the arrangements.
\end{itemize}

\needspace{5\baselineskip}\section*{Footer}
\begin{itemize}
\item The fundamental equation (16.84): $T\,dS = dU + p\,dV$, the first and second laws together.
\item \textbf{"Heat death of the Universe"} (pp. 413–414): this is wrong, because the Universe is not a thermodynamic system.
\item Optional: characteristic functions H = U + pV, F = U − TS, G = U − TS + pV (pp. 408–409).
\end{itemize}

\needspace{5\baselineskip}\section*{Self-check}
\begin{enumerate}
\item Why is d'Q not an exact differential while d'Q/T is?
\item In an isolated system, can entropy decrease? Answer strictly (thermodynamics) and statistically.
\item Why is the even distribution of molecules the equilibrium?
\end{enumerate}

\chapter*{15. Теорема Нернста}\addcontentsline{toc}{chapter}{15. Теорема Нернста}\markboth{15. Теорема Нернста}{15. Теорема Нернста}


\srcline{\book{} т. 1, \textbf{§ 16.11} «Третій закон термодинаміки. Недосяжність абсолютного нуля температур» (с. 414–415). The book calls this the \textbf{third law}; «теорема Нернста» is its historical name. \ok{} (16.93) and the consequences checked against the scan (pp. 414–415).}
Block F, page 5 of 5. A short topic: half a page.

\needspace{5\baselineskip}\section*{Layout}
\begin{layout}
┌ [A] Навіщо: ентропія відома лише до сталої (16.76) → потрібен новий постулат ┐
├ [B] Формулювання Нернста (1906) — рамка                                     ┤
├ [C] Формулювання Планка: lim S = 0 при T → 0 (16.93)                        ┤
└ [D] Наслідки: недосяжність 0 К · C_p, C_V → 0 · коефіцієнти розширення → 0  ┘
\end{layout}

\needspace{5\baselineskip}\section*{[A] Навіщо (p. 414)}
\begin{itemize}
\item The first and second laws give only \textbf{changes} in entropy; S₀ in (16.76) stays undetermined. That's enough until the \textbf{chemical composition} changes, which is the problem of chemical affinity.
\item So a \textbf{new postulate} was needed, generalizing new experiments at low temperatures.
\end{itemize}

\needspace{5\baselineskip}\section*{[B] Формулювання Нернста (1906) (pp. 414–415)}
\begin{itemize}
\item «при наближенні температури до абсолютного нуля ентропія будь-якої рівноважної системи при ізотермічних процесах \textbf{перестає залежати від термодинамічних параметрів стану} і при T → 0 набуває \textbf{ту саму для всіх систем сталу величину}, яку можна вважати такою, що дорівнює нулю»
\item Nernst considered only \textbf{condensed} systems: solids and liquids.
\end{itemize}

\needspace{5\baselineskip}\section*{[C] Формулювання Планка (p. 415)}
\begin{itemize}
\item Planck went further: entropy itself goes to zero, for \textbf{any} system:
\[\lim_{T\to0}S = 0\qquad(16.93)\]
the most general statement of the third law.
\end{itemize}

\needspace{5\baselineskip}\section*{[D] Наслідки (p. 415)}
\begin{enumerate}
\item \textbf{Absolute zero can't be reached}. All cooling methods use adiabatic processes. Near 0 К the difference between adiabatic and isothermal processes vanishes, so cooling becomes ever less effective. You can only approach 0 К asymptotically. The law is sometimes stated as the «принцип недосяжності абсолютного нуля».
\item $C_p, C_V \to 0$ as T → 0. This is confirmed by experiment and by the quantum theories of Einstein and Debye; compare 05, where classical theory fails.
\item The thermal coefficients of expansion and pressure → 0.
\end{enumerate}

\needspace{5\baselineskip}\section*{The three laws together (good for the page footer)}
{\small
\begin{longtable}{>{\raggedright\arraybackslash}p{0.086\linewidth}>{\raggedright\arraybackslash}p{0.882\linewidth}}
\toprule
\textbf{закон} & \textbf{суть} \\
\midrule
\endhead
I & energy is conserved: Q = ΔU + A \\
II & direction: S of an isolated system doesn't decrease \\
III & S → 0 as T → 0; 0 К can't be reached \\
\bottomrule
\end{longtable}
}

\needspace{5\baselineskip}\section*{Self-check}
\begin{enumerate}
\item Why couldn't the first two laws fix the value of entropy?
\item Why does adiabatic cooling stop working near 0 К?
\end{enumerate}

\chapter*{26. Відмінність реального газу від ідеального}\addcontentsline{toc}{chapter}{26. Відмінність реального газу від ідеального}\markboth{26. Відмінність реального газу від ідеального}{26. Відмінність реального газу від ідеального}


\srcline{\book{} т. 1, \textbf{§ 17.1} «Реальні гази та відхилення їхніх властивостей від законів ідеального газу. Сили міжмолекулярної взаємодії в газах» (с. 419–422). \ok{} Checked against the scan (pp. 419–421).}
Block G, page 1 of 2.

\needspace{5\baselineskip}\section*{Layout}
\begin{layout}
┌ [A] Означення реального газу; дві причини відхилень (власний об'єм, притягання) ┐
├ [B] Кількісні vs якісні відхилення; коли модель ідеального газу ще працює    ┤
├ [C] ГРАФІК сили F(r): відштовхування/притягання, r₀ (рис. 17.2)              ┤
├ [D] Сили Ван дер Ваальса: Кеєзом · Дебай · Лондон (таблиця) + водневий зв'язок ┤
└ [E] Ізотерми Амага, крива і температура Бойля (optional)                      ┘
\end{layout}

\needspace{5\baselineskip}\section*{[A] Реальний газ (p. 419)}
\begin{itemize}
\item «\textbf{Реальними} називають гази, для яких суттєву роль відіграють міжмолекулярні сили взаємодії»
\item Compared with the ideal model (non-interacting point molecules), two things are added:
\begin{enumerate}
\item \textbf{the molecules' own volume}. At normal conditions it is about 10⁻⁴ of the gas volume, but at about 10⁸ Па it is significant;
\item \textbf{intermolecular attraction}, which adds extra pressure.
\end{enumerate}
\end{itemize}

\needspace{5\baselineskip}\section*{[B] Характер відхилень (pp. 419–420)}
\begin{itemize}
\item \textbf{Кількісні} (quantitative): the ideal-gas laws, pV = RT included, hold only approximately.
\item \textbf{Якісні} (qualitative): interaction lets a gas turn into a \textbf{liquid and a solid}, which is impossible for an ideal gas.
\item The ideal-gas laws hold at \textbf{low pressure and high temperature}, when the average potential energy of interaction ≪ the average kinetic energy.
\end{itemize}

\needspace{5\baselineskip}\section*{[C] Сили міжмолекулярної взаємодії (pp. 420–421)}
\begin{itemize}
\item Repulsion $F_1 \propto 1/r^n$ (n = 9–15) acts at short range; attraction $F_2 \propto 1/r^m$ (m = 7) acts further out.
\item At r = r₀, F = 0, the \textbf{рівноважна відстань} (equilibrium distance). For r > r₀ attraction wins; for r < r₀ repulsion wins.
\item \pen{} The main drawing, рис. 17.2: F(r), with the repulsion curve above, the attraction curve below, their sum and the zero at r₀. Optionally add the potential energy curve with a well at r₀.
\end{itemize}

\needspace{5\baselineskip}\section*{[D] Сили Ван дер Ваальса (pp. 421–422)}
\begin{itemize}
\item Weak forces of attraction at distances of about 10⁻⁹ м.
\end{itemize}
{\footnotesize
\begin{longtable}{>{\raggedright\arraybackslash}p{0.138\linewidth}>{\raggedright\arraybackslash}p{0.069\linewidth}>{\raggedright\arraybackslash}p{0.219\linewidth}>{\raggedright\arraybackslash}p{0.305\linewidth}>{\raggedright\arraybackslash}p{0.188\linewidth}}
\toprule
\textbf{тип} & \textbf{автор} & \textbf{у яких молекул} & \textbf{природа} & \textbf{особливості} \\
\midrule
\endhead
орієнтаційні & Кеєзом & polar (H₂O, NH₃) & permanent dipoles & depend on T, not additive \\
індукційні & Дебай & polar + induced & induced dipole & small, don't depend on T \\
дисперсійні & Лондон & \textbf{non-polar} (the main ones) & instantaneous dipoles from electron motion & additive \\
\bottomrule
\end{longtable}
}
\begin{itemize}
\item \textbf{Водневий зв'язок} (hydrogen bond), p. 422: specific attraction in water and alcohols through OH groups.
\end{itemize}

\needspace{5\baselineskip}\section*{[E] Optional: ізотерми Амага (p. 420)}
\begin{itemize}
\item In pV–p coordinates an ideal gas gives straight lines; a real gas's curves first dip (attraction) and then rise (own volume). The minimum lies on the \textbf{крива Бойля}; $T_{\text{Б}}$ is the \textbf{температура Бойля}.
\end{itemize}

\needspace{5\baselineskip}\section*{Self-check}
\begin{enumerate}
\item Name the two causes of deviation from the ideal gas, and say when each one dominates.
\item Why can't an ideal gas be liquefied?
\item Which van der Waals forces act between N₂ molecules?
\end{enumerate}

\chapter*{27. Рівняння Ван-Дер-Ваальса}\addcontentsline{toc}{chapter}{27. Рівняння Ван-Дер-Ваальса}\markboth{27. Рівняння Ван-Дер-Ваальса}{27. Рівняння Ван-Дер-Ваальса}


\srcline{\book{} т. 1, \textbf{§ 17.2} «Рівняння Ван-дер-Ваальса» (с. 422–424), \textbf{§ 17.3} «Ізотерми газу Ван-дер-Ваальса. Критичний стан речовини…» (с. 425–430), \textbf{§ 17.4} (с. 430–432). \ok{} (17.4), (17.9)–(17.12) checked against the scan.}
Block G, page 2 of 2. Spelling: the book writes \textbf{Ван-дер-Ваальса}.

\needspace{5\baselineskip}\section*{Layout}
\begin{layout}
┌ [A] Рівняння (17.4) у рамці + що означає кожна поправка          ┐
├ [B] Звідки a/V² (внутрішній тиск) і b (4 × об'єм молекул)        ┤
├ [C] ГРАФІК ізотерм (рис. 17.4/17.5): хвиля, горизонталь, критична точка K ┤
├ [D] Критичні параметри (17.12); критичний стан                  ┤
└ [E] Порівняння з дослідом; зведене рівняння (17.18)              ┘
\end{layout}

\needspace{5\baselineskip}\section*{[A] Рівняння (p. 423) — put it in a box}
\[\left(p + \frac a{V_M^2}\right)(V_M - b) = RT\qquad(17.4)\]
\begin{itemize}
\item For one mole. $V_{M}$ is the molar volume; \textbf{a, b} are experimental constants of the gas.
\item \textbf{a/$V_{M}$²} accounts for the \textbf{attraction} of molecules. It has the dimension of pressure and is called the \textbf{внутрішній тиск} (internal pressure), $p_{i.}$
\item \textbf{b} corrects for the molecules' \textbf{own volume} (repulsion at short range).
\item For mass m (17.7): $\left(p + \dfrac{m^2}{M^2}\dfrac a{V^2}\right)\left(V - \dfrac mMb\right) = \dfrac mMRT$
\item It is only approximate: quantitatively good at high T and low p, but \textbf{qualitatively} valid everywhere.
\end{itemize}

\needspace{5\baselineskip}\section*{[B] Звідки поправки (pp. 423–424)}
\begin{itemize}
\item \textbf{b}: the centres of two hard spheres of radius r can't come closer than 2r, so the "forbidden" volume per molecule is ½·(4/3)π(2r)³ = 4·(4/3)πr³. So \textbf{b = 4 × the volume of the molecules} in a mole.
\item \textbf{a/V²}: a molecule near the wall is pulled back into the gas. That force ∝ n₀ and the number of molecules hitting the wall ∝ n₀, so $p_{i}$ ∝ n₀² ∝ 1/V².
\item \pen{} Sketch a molecule near the wall with the forces from its neighbours pulling it inwards.
\end{itemize}

\needspace{5\baselineskip}\section*{[C] Ізотерми (pp. 425–430)}
\begin{itemize}
\item Multiplying out gives a \textbf{cubic in $V_{M}$} (17.8): for given p and T there can be \textbf{three} volumes.
\item Below $T_{\text{к}}$: an S-shaped "wave". On the \textbf{experimental} isotherm it is replaced by a \textbf{horizontal} line (рис. 17.5), where liquid and vapour coexist at constant p.
\begin{itemize}
\item The right branch is gas; the left, steep branch is liquid (it barely compresses).
\item \textbf{Метастабільні стани} (metastable): the section bc is \textbf{пересичена пара} (supersaturated vapour); ef is \textbf{перегріта рідина} (superheated liquid).
\item The section cde is absolutely unstable (∂p/∂V > 0), so it can't be realized.
\end{itemize}
\item \textbf{Бінодаль} (binodal) bounds the two-phase region; \textbf{спінодаль} (spinodal) bounds the unstable states.
\item Above $T_{\text{к}}$: monotonic hyperbolas, so a gas \textbf{can't be liquefied by pressure alone} (CO₂ above 31 °C).
\item \pen{} The main drawing: a family of isotherms T₁ < $T_{\text{к}}$ < T₂, the wave and the horizontal line on T₁, the inflection point K, and the bell-shaped binodal with labelled regions "рідина", "рідина + пара", "газ".
\end{itemize}

\needspace{5\baselineskip}\section*{[D] Критичний стан (pp. 426–427)}
\begin{itemize}
\item At T = $T_{\text{к}}$ the three roots merge and the isotherm has an \textbf{inflection point K}.
\[p_\text{к} = \frac a{27b^2},\qquad V_\text{к} = 3b,\qquad T_\text{к} = \frac{8a}{27Rb}\qquad(17.12)\]
\item In the critical state the difference between liquid and gas disappears, and the compressibility → ∞.
\item You can go from gas to liquid \textbf{without a phase transition} by going round the two-phase region via T > $T_{\text{к}}$ (p. 430).
\end{itemize}

\needspace{5\baselineskip}\section*{[E] Порівняння з дослідом (pp. 430–432), short}
\begin{itemize}
\item a and b depend on T. The \textbf{критичний коефіцієнт} $K_\text{к} = \dfrac{RT_\text{к}}{p_\text{к}V_\text{к}} = \dfrac83 \approx 2{,}67$ in theory, but for H₂ it is 3,12 and for water 4,46. In reality $V_{\text{к}}$ ≈ 2b.
\item \textbf{Зведене рівняння} (reduced equation), with π = p/$p_{\text{к}}$, ω = V/$V_{\text{к}}$, τ = T/$T_{\text{к}}$: $\left(\pi + \dfrac3{\omega^2}\right)(3\omega - 1) = 8\tau$ (17.18), the same for all substances. This is the \textbf{закон відповідних станів} (law of corresponding states).
\end{itemize}

\needspace{5\baselineskip}\section*{Self-check}
\begin{enumerate}
\item Which constant stands for attraction, and which for repulsion?
\item Why is b four times the molecules' own volume, not equal to it?
\item Why can't you liquefy a gas above $T_{\text{к}}$ however hard you compress it?
\end{enumerate}

\chapter*{28. Класифікація кристалів}\addcontentsline{toc}{chapter}{28. Класифікація кристалів}\markboth{28. Класифікація кристалів}{28. Класифікація кристалів}


\srcline{\book{} т. 1, \textbf{§ 18.1} «Кристалічні та аморфні тіла. Основні характеристики кристалів» (с. 444–450); classification by type of bond, \textbf{§ 18.2} «Сили зв'язку у твердих тілах» (с. 450–453). \ok{} Checked against the scan (pp. 446–448).}
Block H, page 1 of 3. The book classifies crystals in \textbf{two ways}: by \textbf{symmetry} (syngonies, Bravais lattices) and by \textbf{type of bond}. Cover both.

\needspace{5\baselineskip}\section*{Layout}
\begin{layout}
┌ [A] Тверде тіло: кристалічне vs аморфне (дальній/ближній порядок, анізотропія) ┐
├ [B] Моно- і полікристали; закони кристалографії; елементи симетрії          ┤
├ [C] Гратка, вузли, елементарна комірка, rₙ = ma₁ + na₂ + pa₃ (18.1)          ┤
├ [D] ТАБЛИЦЯ 7 сингоній + 4 типи граток Браве (14 разом) + малюнки кубічних   ┤
└ [E] Класифікація за типом зв'язку: іонні · ковалентні · металічні · молекулярні ┘
\end{layout}

\needspace{5\baselineskip}\section*{[A] Кристалічні та аморфні тіла (pp. 444–445)}
\begin{itemize}
\item \textbf{Тверде тіло} (solid): «агрегатний стан речовини, який характеризується стабільністю форми та об'єму», plus ordered positions of atoms that vibrate slightly about equilibrium.
\end{itemize}
{\small
\begin{longtable}{>{\raggedright\arraybackslash}p{0.143\linewidth}>{\raggedright\arraybackslash}p{0.535\linewidth}>{\raggedright\arraybackslash}p{0.274\linewidth}}
\toprule
\textbf{} & \textbf{кристалічні} & \textbf{аморфні} \\
\midrule
\endhead
порядок & \textbf{дальній} (long-range): strict periodic repetition in 3D & only \textbf{ближній} (short-range) \\
властивості & \textbf{анізотропія} (anisotropy): properties depend on direction & \textbf{ізотропні} (isotropic) \\
приклади & metals, salt, quartz & glass, amber, resin, polymers \\
\bottomrule
\end{longtable}
}
\begin{itemize}
\item \textbf{Монокристал} (single crystal, one large regular crystal) vs \textbf{полікристал} (polycrystal, many small grains fused together at random; most metals).
\end{itemize}

\needspace{5\baselineskip}\section*{[B] Закони кристалографії і симетрія (pp. 445–446)}
\begin{itemize}
\item Flat faces and straight edges; \textbf{constant angles between faces}; rational ratios of the intercepts on the axes.
\item \textbf{Закон симетрії}: only axes of order \textbf{1, 2, 3, 4 and 6} are possible (no 5th order).
\item Elements of symmetry: the \textbf{вісь симетрії} (axis, rotation by 2π/n), the \textbf{площина симетрії} P (plane, mirror) and the \textbf{центр інверсії} C (centre of inversion).
\end{itemize}

\needspace{5\baselineskip}\section*{[C] Просторова гратка (pp. 446–447)}
\begin{itemize}
\item Put a point at the centre of mass of each structural unit to get the \textbf{просторові гратки} (space lattice); its points are the \textbf{вузли} (nodes).
\item \textbf{Елементарна комірка} (unit cell): a parallelepiped built on the vectors a₁, a₂, a₃, whose lengths are the \textbf{періоди ідентичності} (lattice periods). Any node is at $\vec r_n = m\vec a_1 + n\vec a_2 + p\vec a_3$ (18.1).
\item \textbf{Гратки Браве} (Bravais lattices): unit cells of minimal area that keep the full symmetry. There are \textbf{14} of them in \textbf{7 сингоній} (crystal systems).
\end{itemize}

\needspace{5\baselineskip}\section*{[D] 7 сингоній (pp. 447–449, Table 18.4)}
{\small
\begin{longtable}{>{\raggedright\arraybackslash}p{0.320\linewidth}>{\raggedright\arraybackslash}p{0.111\linewidth}>{\raggedright\arraybackslash}p{0.259\linewidth}>{\raggedright\arraybackslash}p{0.246\linewidth}}
\toprule
\textbf{сингонія} & \textbf{ребра} & \textbf{кути} & \textbf{приклади (book)} \\
\midrule
\endhead
триклінна & a ≠ b ≠ c & α ≠ β ≠ γ ≠ 90° & feldspars \\
моноклінна & a ≠ b ≠ c & α = γ = 90° ≠ β & gypsum, mica, talc \\
ромбічна & a ≠ b ≠ c & α = β = γ = 90° & sulphur, topaz \\
ромбоедрична (тригональна) & a = b = c & α = β = γ ≠ 90° & calcite, tellurium \\
тетрагональна & a = b ≠ c & α = β = γ = 90° & rutile, cassiterite \\
гексагональна & a = b ≠ c & α = β = 90°, γ = 120° & ice, graphite, zinc \\
кубічна & a = b = c & α = β = γ = 90° & most metals, diamond \\
\bottomrule
\end{longtable}
}
\begin{itemize}
\item Lattice types: \textbf{примітивна} (primitive, nodes at the corners), \textbf{базоцентрована} (base-centred), \textbf{об'ємноцентрована} (body-centred: Fe-α, W, Na), \textbf{гранецентрована} (face-centred: Al, Cu, Ag, Au).
\item \pen{} Draw the three cubic lattices: simple, body-centred and face-centred.
\item Optional: \textbf{Miller indices} (h k l) for planes, p. 449; \textbf{230 space groups} (Fedorov, 1890); \textbf{32 classes} of symmetry.
\end{itemize}

\needspace{5\baselineskip}\section*{[E] За типом зв'язку (§ 18.2, pp. 450–453)}
{\footnotesize
\begin{longtable}{>{\raggedright\arraybackslash}p{0.138\linewidth}>{\raggedright\arraybackslash}p{0.103\linewidth}>{\raggedright\arraybackslash}p{0.269\linewidth}>{\raggedright\arraybackslash}p{0.289\linewidth}>{\raggedright\arraybackslash}p{0.120\linewidth}}
\toprule
\textbf{тип кристала} & \textbf{у вузлах} & \textbf{зв'язок} & \textbf{властивості} & \textbf{приклади} \\
\midrule
\endhead
іонний & ± ions & electrostatic & strong, high $T_{melt}$, poor conductor & NaCl \\
ковалентний (атомний) & neutral atoms & shared electron pairs, \textbf{directional, saturated} & very hard, high $T_{melt}$, dielectrics or semiconductors & diamond, Si \\
металічний & + ions & \textbf{electron gas} & good electrical and thermal conductor, shiny, ductile & Cu, Au \\
молекулярний & molecules & van der Waals, hydrogen & low $T_{melt}$, insulators & organic crystals, ice \\
\bottomrule
\end{longtable}
}

\needspace{5\baselineskip}\section*{Self-check}
\begin{enumerate}
\item What is the difference between long-range and short-range order?
\item Why can't a crystal have a 5th-order symmetry axis? (Pentagons can't tile a plane.)
\item Why do metals conduct electricity while covalent crystals don't?
\end{enumerate}

\chapter*{29. Дефекти в кристалах}\addcontentsline{toc}{chapter}{29. Дефекти в кристалах}\markboth{29. Дефекти в кристалах}{29. Дефекти в кристалах}


\srcline{\book{} т. 1, \textbf{§ 18.3} «Дефекти в реальних кристалах» (с. 453–456). \ok{} Checked against the scan (pp. 453–456).}
Block H, page 2 of 3. The structure is simple: \textbf{classification by dimension (0D–3D)}.

\needspace{5\baselineskip}\section*{Layout}
\begin{layout}
┌ [A] Навіщо: дефекти визначають структурно-чутливі властивості           ┐
├ [B] Класифікація за вимірністю — головна таблиця 0D/1D/2D/3D            ┤
├ [C] Точкові: Шотткі, Френкель, домішки (рис. 18.5 а/б/в) + n_в (18.4)    ┤
├ [D] Лінійні: дислокації — крайова (рис. 18.6) і гвинтова (рис. 18.7); вектор Бюргерса ┤
└ [E] Поверхневі та об'ємні (коротко)                                     ┘
\end{layout}

\needspace{5\baselineskip}\section*{[A] Навіщо (p. 453)}
\begin{itemize}
\item A perfect structure exists only in small volumes. Defects largely determine the \textbf{структурно-чутливі властивості} (structure-sensitive properties): electrical and thermal conductivity, strength…
\end{itemize}

\needspace{5\baselineskip}\section*{[B] Класифікація (p. 453) — the main table}
{\small
\begin{longtable}{>{\raggedright\arraybackslash}p{0.169\linewidth}>{\raggedright\arraybackslash}p{0.391\linewidth}>{\raggedright\arraybackslash}p{0.392\linewidth}}
\toprule
\textbf{вимірність} & \textbf{тип} & \textbf{приклади} \\
\midrule
\endhead
0 (нульвимірні) & \textbf{точкові} (point): size \textasciitilde{} one or a few interatomic distances in all directions & vacancies (Schottky), Frenkel pairs, impurity atoms \\
1 (одновимірні) & \textbf{лінійні} (line) & \textbf{дислокації}: edge and screw \\
2 (двовимірні) & \textbf{поверхневі} (planar) & grain boundaries, mosaic blocks, stacking faults, twins, the surface \\
3 (тривимірні) & \textbf{об'ємні} (volume) & cracks, inclusions, pores \\
\bottomrule
\end{longtable}
}

\needspace{5\baselineskip}\section*{[C] Точкові дефекти (pp. 453–455)}
\begin{itemize}
\item \textbf{Atomic} ones:
\begin{itemize}
\item \textbf{дефект за Шотткі} (Schottky defect): an empty node, a vacancy (рис. 18.5 а);
\item \textbf{дефект за Френкелем} (Frenkel defect): a vacancy \textbf{plus} a nearby interstitial atom (рис. 18.5 в). It needs a little less energy than a Schottky defect, 2–4 еВ in both cases;
\item \textbf{домішкові атоми} (impurity atoms): interstitial or substituting a host atom (рис. 18.5 б).
\end{itemize}
\item Electronic ones (free electrons, holes, the \textbf{exciton} and the \textbf{polaron}, Pekar's theory) and energy ones (\textbf{phonons}, quanta of lattice vibration): mention them in one line.
\item Equilibrium concentration of vacancies: $n_{\text{в}} = n_0e^{-W_{\text{в}}/kT}$ (18.4). It grows quickly with T; near melting $n_{\text{в}}$/n₀ ≈ 10⁻⁴. Compare Boltzmann's distribution on 17, 19.
\item \pen{} Draw three small square grids: one with a missing atom, one with an atom moved from its node into a gap, one with a foreign atom.
\end{itemize}

\needspace{5\baselineskip}\section*{[D] Дислокації (pp. 455–456)}
\begin{itemize}
\item \textbf{Крайова} (edge dislocation, рис. 18.6): one atomic plane ends inside the crystal, like a book with one page torn halfway.
\item \textbf{Гвинтова} (screw dislocation, рис. 18.7): two half-planes slip by one period relative to each other.
\item Dislocations can't end inside a crystal; they form loops or reach the surface.
\item \textbf{Вектор Бюргерса} (Burgers vector): go round a contour by a whole number of steps. Around a dislocation the contour \textbf{doesn't close}, and the gap a′a″ is the Burgers vector, a measure of the lattice distortion.
\item \textbf{Густина дислокацій} (dislocation density): the number of dislocation lines crossing a unit area.
\item \pen{} Sketch the "extra half-plane" with a Burgers contour round it.
\end{itemize}

\needspace{5\baselineskip}\section*{[E] Поверхневі та об'ємні (p. 456)}
\begin{itemize}
\item \textbf{Двійники} (twins): two crystals of one substance grown together along a crystallographic plane that is a mirror plane.
\end{itemize}

\needspace{5\baselineskip}\section*{Self-check}
\begin{enumerate}
\item Schottky vs Frenkel: what's the difference?
\item Why does the number of vacancies grow with temperature, and by what law?
\item How do you detect a dislocation with a Burgers contour?
\end{enumerate}

\chapter*{30. Рідини}\addcontentsline{toc}{chapter}{30. Рідини}\markboth{30. Рідини}{30. Рідини}


\srcline{\book{} т. 1, \textbf{§ 19.1} «Загальні властивості та будова рідин» (с. 462–466), \textbf{§ 19.2} «В'язкість рідин» (с. 467). \ok{} Checked against the scan (pp. 462–467).}
Block H, page 3 of 3. \warn{} This question was \textbf{handwritten} on the МКР sheet («30. Рідини»), so its scope is unclear. This page covers the structure and properties of liquids (§19.1–19.2). If your lecturer means surface tension and capillarity, see the "possible extension" at the end.

\needspace{5\baselineskip}\section*{Layout}
\begin{layout}
┌ [A] Рідина між газом і твердим тілом: ближній порядок без дальнього     ┐
├ [B] ГРАФІК радіальної функції розподілу N(r): тверде/рідина/газ (рис. 19.1) ┤
├ [C] Тепловий рух за Френкелем: коливання + стрибки; час «осілого життя» τ ┤
├ [D] Властивості: мала стисливість β_T, теплове розширення α_p, аномалія води ┤
└ [E] Перенос у рідинах: D, λ, η — і їх залежність від T (протилежна газам!) ┘
\end{layout}

\needspace{5\baselineskip}\section*{[A] Місце рідин (pp. 462–463)}
\begin{itemize}
\item The three states differ in how strong \textbf{the interaction of molecules is compared with their thermal motion}.
\item Liquids sit \textbf{between gases and solids}. Molecules interact strongly, as in a solid, but sit on no periodic lattice.
\item Frenkel's idea: a liquid is a "\textbf{disordered solid}", with \textbf{short-range order} and \textbf{no long-range order} (28).
\end{itemize}

\needspace{5\baselineskip}\section*{[B] Радіальна функція розподілу (pp. 462–463)}
\begin{itemize}
\item N(r) is the number of neighbours at a distance r from a chosen molecule.
\begin{itemize}
\item Solid: sharp discrete peaks to large r (fixed positions).
\item Liquid: a few blurred maxima, damped towards a constant (short-range order).
\item Gas: smooth, with no structure.
\end{itemize}
\item \textbf{Bernal's model}: the structure is described by the radii of the coordination spheres, the number of molecules in them, and how blurred they are.
\item \pen{} Draw three small graphs N(r) side by side, as in рис. 19.1 а/б/в.
\end{itemize}

\needspace{5\baselineskip}\section*{[C] Тепловий рух (Френкель) (pp. 463–464)}
\begin{itemize}
\item \textbf{Oscillation} about a temporary equilibrium (period about 10⁻¹² с), then a \textbf{jump} to a neighbouring position. These jumps are \textbf{трансляційні рухи} and need vacancies, or "holes".
\item The \textbf{час осілого життя} (residence time): $\tau = Ae^{\Delta W/kT}$ (19.1), where ΔW is the potential barrier. It shrinks as T rises.
\item The jumps explain \textbf{fluidity} and \textbf{diffusion} in liquids.
\end{itemize}

\needspace{5\baselineskip}\section*{[D] Властивості (pp. 464–466)}
\begin{itemize}
\item They keep their \textbf{volume} (not their shape) and have a \textbf{free surface}.
\item \textbf{Мала стисливість} (low compressibility): $\beta_T = -\dfrac1V\left(\dfrac{\partial V}{\partial p}\right)_T$ (19.2), about 10⁻⁹–10⁻¹⁰ Па⁻¹, tens to hundreds of thousands of times smaller than for gases. This is strong repulsion at work.
\item \textbf{Теплове розширення} (thermal expansion): $\alpha_p = \dfrac1V\left(\dfrac{\partial V}{\partial T}\right)_p$ (19.3), usually > 0.
\begin{itemize}
\item \spark{} \textbf{The water anomaly}: between 0 and 3,98 °C water \textbf{shrinks} when heated, so its density is greatest at 3,98 °C. The cause is its open, tetrahedral, hydrogen-bonded structure.
\end{itemize}
\item Heat capacity is close to that of solids and much larger than that of gases, because heat goes into both kinetic and potential energy.
\end{itemize}

\needspace{5\baselineskip}\section*{[E] Перенос у рідинах (pp. 466–467) — contrast with gases (23–25)}
{\small
\begin{longtable}{>{\raggedright\arraybackslash}p{0.104\linewidth}>{\raggedright\arraybackslash}p{0.229\linewidth}>{\raggedright\arraybackslash}p{0.619\linewidth}}
\toprule
\textbf{} & \textbf{гази} & \textbf{рідини} \\
\midrule
\endhead
механізм & free flights over ⟨l⟩ & jumps between equilibrium positions \\
D & ⅓⟨v⟩⟨l⟩; 0,1–0,9 см²/с & ⟨δ²⟩/(6τ) (19.4); 4–5 orders of magnitude smaller; $D = Be^{-\Delta W/kT}$ (19.5) \\
η від T & \textbf{grows} with T & \textbf{falls} with T: $\eta = Ae^{\Delta W/kT}$, \textbf{Frenkel's formula} (19.7) \\
λ від T & grows & falls roughly linearly \\
\bottomrule
\end{longtable}
}
\begin{itemize}
\item Water's viscosity falls by 20 \% from 20 to 30 °C; glycerine's falls 2,4 times.
\end{itemize}

\needspace{5\baselineskip}\section*{Possible extension, if "рідини" includes surface phenomena}
\begin{itemize}
\item § 19.3 «Поверхневий шар рідини. Поверхневий натяг» (с. 468–471), § 19.4 «Формула Лапласа. Змочування та капілярні явища» (с. 471–474). Not yet in the wiki. Ask your lecturer, and tell me to add a page if needed.
\end{itemize}

\needspace{5\baselineskip}\section*{Self-check}
\begin{enumerate}
\item Why is a liquid closer to a solid than to a gas?
\item Why does a liquid's viscosity fall with T while a gas's grows?
\item Why is water densest at 4 °C?
\end{enumerate}

\appendix

\part{Додатки}

\chapter*{Інтегральна ↔ диференціальна форма законів}\addcontentsline{toc}{chapter}{Інтегральна ↔ диференціальна форма законів}\markboth{Інтегральна ↔ диференціальна форма законів}{Інтегральна ↔ диференціальна форма законів}


The same pattern comes up in at least four syllabus topics. The integral form links \textbf{different places} (a surface and a volume, or a loop). The differential form links quantities \textbf{at one point}, which is what close-range action needs. The bridges are two mathematical theorems.

{\footnotesize
\begin{longtable}{>{\raggedright\arraybackslash}p{0.209\linewidth}>{\raggedright\arraybackslash}p{0.172\linewidth}>{\raggedright\arraybackslash}p{0.172\linewidth}>{\raggedright\arraybackslash}p{0.270\linewidth}>{\raggedright\arraybackslash}p{0.098\linewidth}}
\toprule
\textbf{Закон} & \textbf{інтегральна} & \textbf{диференціальна} & \textbf{міст} & \textbf{topics} \\
\midrule
\endhead
Гаусс (джерела E) & ∮E dS = q/ε₀ & div E = ρ/ε₀ & Остроградського–Гаусса & 6, 7, 13 \\
потенціальність E & ∮E dl = 0 & rot E = 0 & Стокса & 4, 9 \\
збереження заряду & ∮j dS = −∂q/∂t & ∂ρ/∂t = −div j & Остроградського–Гаусса & 1, 20 \\
Ом & I = U/R & j = σE & (local cylinder dl·dS) & 19 \\
Джоуль–Ленц & Q = I²Rt & w = σE² = jE & (local cylinder) & 23 \\
\bottomrule
\end{longtable}
}

Book's framing (т. 2, p. 41): «Теорема Гаусса дає можливість звести основні рівняння електростатики до диференціальної форми й узгодити їх, таким чином, з теорією близькодії.»

\chapter*{Принцип суперпозиції}\addcontentsline{toc}{chapter}{Принцип суперпозиції}\markboth{Принцип суперпозиції}{Принцип суперпозиції}


Fields from several sources add \textbf{without affecting one another}. This is an experimental fact, not something that follows from Coulomb's law (т. 2, p. 25).

{\small
\begin{longtable}{>{\raggedright\arraybackslash}p{0.171\linewidth}>{\raggedright\arraybackslash}p{0.531\linewidth}>{\raggedright\arraybackslash}p{0.169\linewidth}>{\raggedright\arraybackslash}p{0.065\linewidth}}
\toprule
\textbf{Величина} & \textbf{як додається} & \textbf{формула} & \textbf{topic} \\
\midrule
\endhead
сила & vector sum & F = ΣFᵢ (1.10) & 2 \\
напруженість & vector sum & E = ΣEᵢ (1.13) & 3 \\
потенціал & \textbf{algebraic} sum & φ = Σφᵢ (1.56) & 4 \\
поле двох площин & the fields add between the plates and cancel outside & σ/ε₀ & 11 \\
поле диполя & the fields of +q and −q add & (1.20)–(1.22) & 5 \\
\bottomrule
\end{longtable}
}

The book's wording: «суперпозиція — це здатність фізичних полів до накладання, при якому векторні характеристики їх додаються геометрично, а після розходження їх ці поля існують незалежно одне від одного».

\chapter*{Процеси в ідеальному газі: одна таблиця}\addcontentsline{toc}{chapter}{Процеси в ідеальному газі: одна таблиця}\markboth{Процеси в ідеальному газі: одна таблиця}{Процеси в ідеальному газі: одна таблиця}


ν = m/M moles. The first law: Q = ΔU + A. For an ideal gas always \textbf{ΔU = $νC_{V}$ΔT}. Sources: т. 1, § 14.7, § 16.3–16.4.

{\scriptsize
\begin{longtable}{>{\raggedright\arraybackslash}p{0.114\linewidth}>{\raggedright\arraybackslash}p{0.058\linewidth}>{\raggedright\arraybackslash}p{0.167\linewidth}>{\raggedright\arraybackslash}p{0.133\linewidth}>{\raggedright\arraybackslash}p{0.106\linewidth}>{\raggedright\arraybackslash}p{0.057\linewidth}>{\raggedright\arraybackslash}p{0.057\linewidth}>{\raggedright\arraybackslash}p{0.071\linewidth}>{\raggedright\arraybackslash}p{0.093\linewidth}}
\toprule
\textbf{процес} & \textbf{умова} & \textbf{рівняння} & \textbf{закон} & \textbf{A (газ)} & \textbf{ΔU} & \textbf{Q} & \textbf{C} & \textbf{n (pVⁿ)} \\
\midrule
\endhead
ізохорний & V = const & p/T = const & (14.28) & 0 & $νC_{V}$ΔT & = ΔU & $C_{V}$ & ±∞ \\
ізобарний & p = const & V/T = const & Гей-Люссака & p(V₂−V₁) = νRΔT & $νC_{V}$ΔT & $νC_{p}$ΔT & $C_{p}$ = $C_{V}$ + R & 0 \\
ізотермічний & T = const & pV = const & Бойля–Маріотта & νRT ln(V₂/V₁) & 0 & = A & ∞ & 1 \\
адіабатний & d'Q = 0 & $pV^{γ}$ = const, $TV^{γ−1}$ = const & Пуассона & $νC_{V}$(T₁−T₂) & −A & 0 & 0 & γ \\
політропний & C = const & pVⁿ = const & (16.40) & — & $νC_{V}$ΔT & νCΔT & C & ($C_{p−C}$)/($C_{V−C}$) \\
\bottomrule
\end{longtable}
}

\pen{} One p–V diagram from a common point with all four curves: an isochore (vertical), an isobar (horizontal), an isotherm (hyperbola) and an adiabat (a steeper hyperbola).

Values: one-atom γ = 5/3, two-atom 7/5, many-atom 4/3 (see 05).

\chapter*{Явища переносу: один закон, три величини}\addcontentsline{toc}{chapter}{Явища переносу: один закон, три величини}\markboth{Явища переносу: один закон, три величини}{Явища переносу: один закон, три величини}


All three transport phenomena have the shape \textbf{потік = −коефіцієнт × градієнт} (flux = −coefficient × gradient). In a gas each coefficient is \textbf{⅓ × (what one molecule carries) × ⟨v⟩ × ⟨l⟩}.

{\footnotesize
\begin{longtable}{>{\raggedright\arraybackslash}p{0.184\linewidth}>{\raggedright\arraybackslash}p{0.104\linewidth}>{\raggedright\arraybackslash}p{0.239\linewidth}>{\raggedright\arraybackslash}p{0.155\linewidth}>{\raggedright\arraybackslash}p{0.165\linewidth}>{\raggedright\arraybackslash}p{0.057\linewidth}}
\toprule
\textbf{} & \textbf{потік чого} & \textbf{закон} & \textbf{градієнт} & \textbf{коефіцієнт (газ)} & \textbf{topic} \\
\midrule
\endhead
дифузія & маси & $j_{m}$ = −D dρ/dx (Фік) & густини & D = ⅓⟨v⟩⟨l⟩ & 23–25 \\
в'язкість & імпульсу & F/ΔS = −η du/dx (Ньютон) & швидкості течії & η = ⅓ρ⟨v⟩⟨l⟩ & 23–25 \\
теплопровідність & енергії & $j_{q}$ = −λ dT/dx (Фур'є) & температури & λ = ⅓ρ⟨v⟩⟨l⟩$c_{V}$ & 23–25 \\
\bottomrule
\end{longtable}
}

The same "flux ∝ −gradient" shape appears in electricity: \textbf{Ohm's law j = σE = −σ grad φ} (electrostatics 19). All four laws are local, differential laws; see інтегральна ↔ диференціальна форма.

The microscopic ingredient they share is the mean free path ⟨l⟩ (22).

\chapter*{Глосарій}\addcontentsline{toc}{chapter}{Глосарій}\markboth{Глосарій}{Глосарій}
\needspace{5\baselineskip}\section*{1. Електростатика}

\needspace{5\baselineskip}\subsection*{Тема 1. Електричний заряд та його властивості (§ 1.1)}

\begin{itemize}
\item \textbf{електричний заряд} — внутрішня характеристика деяких елементарних матеріальних частинок, яка проявляється в їхній взаємодії на відстані або під час дії на них зовнішніх електромагнітних полів; є як джерелом, так і об'єктом дії електромагнітного поля · $F_1/F_2 = q_1/q_2$ (1.1) · [Кл] · (EN) electric charge · \book{} т. 2, с. 14–15
\item \textbf{електростатика} — розділ фізики, в якому вивчають властивості й взаємодію нерухомих електричних зарядів та властивості їхніх електричних полів · (EN) electrostatics · \book{} т. 2, с. 14
\item \textbf{електростатичне поле} — поле нерухомих зарядів, характеристики якого в кожній точці не змінюються з часом (функції лише координат) · (EN) electrostatic field · \book{} т. 2, с. 14–15
\item \textbf{позитивні / негативні заряди} — два види зарядів; позитивні виникають на скляній паличці при натиранні шовком, негативні — на ебонітовій (янтарній) при натиранні хутром; однойменні відштовхуються, різнойменні притягуються · (EN) positive / negative charges · \book{} т. 2, с. 15
\item \textbf{електризація тертям} — перерозподіл електронів між тілами при щільному контакті; тіла набувають однакових за модулем і протилежних за знаком зарядів, сумарний заряд не змінюється · (EN) triboelectric charging · \book{} т. 2, с. 16
\item \textbf{елементарний заряд} — найменша порція електрики, що дорівнює за модулем заряду електрона · $|e| = 1{,}602\,189\,2(46)\cdot10^{-19}$ Кл · (EN) elementary charge · \book{} т. 2, с. 16, 18
\item \textbf{закон квантування електричного заряду} — усі заряди, які існують у природі у вільному стані, складаються з окремих дискретних зарядів, однакових за абсолютним значенням і таких, що дорівнюють зарядові електрона; приймається як постулат · $q = n|e|$, $n = \pm1, \pm2, \ldots$ · (EN) charge quantization · \book{} т. 2, с. 18–19
\item \textbf{питомий заряд електрона} — відношення заряду частинки до її маси · $e/m = 1{,}7588\cdot10^{11}$ Кл/кг · (EN) specific charge · \book{} т. 2, с. 17
\item \textbf{маса спокою електрона} · $m = 9{,}109\,534(47)\cdot10^{-31}$ кг · (EN) electron rest mass · \book{} т. 2, с. 18
\item \textbf{інваріантність заряду} — незалежність числового значення заряду від стану інерціальних систем відліку (від швидкості руху носія) · (EN) charge invariance · \book{} т. 2, с. 20
\item \textbf{замкнена система} — система, в яку з оточуючого середовища крізь її межі не проникають електрично заряджені частинки · (EN) closed (isolated) system · \book{} т. 2, с. 20
\item \textbf{закон збереження заряду} — повний електричний заряд замкненої системи є величиною постійною і дорівнює алгебричній сумі її позитивних та негативних зарядів · $q = \sum_{i=1}^{n} q_i = \text{const}$ (1.2) · (EN) conservation of charge · \book{} т. 2, с. 20–21
\item \textbf{об'ємна густина заряду} · $\rho = \lim_{\Delta V\to0} \Delta q/\Delta V = dq/dV$ (1.3) · [Кл/м³] · (EN) volume charge density · \book{} т. 2, с. 21
\item \textbf{поверхнева густина заряду} · $\sigma = \lim_{\Delta S\to0} \Delta q/\Delta S = dq/dS$ (1.4) · [Кл/м²] · (EN) surface charge density · \book{} т. 2, с. 21
\item \textbf{лінійна густина заряду} · $\tau = \lim_{\Delta l\to0} \Delta q/\Delta l = dq/dl$ (1.5) · [Кл/м] · (EN) linear charge density · \book{} т. 2, с. 21
\end{itemize}

\needspace{5\baselineskip}\section*{2. Постійний струм}

\needspace{5\baselineskip}\section*{3. Молекулярна фізика та термодинаміка}

\needspace{5\baselineskip}\subsection*{Блок A. Основне рівняння МКТ · Закони ідеальних газів · Ступені вільності, U}

\begin{itemize}
\item \textbf{ідеальний газ} — газ, у якому можна знехтувати розмірами молекул (вважати їх матеріальними точками) та їхньою взаємодією на відстані, враховуючи лише взаємодію під час зіткнень · (EN) ideal gas · \book{} т. 1, с. 300
\item \textbf{тепловий рух} — абсолютно хаотичний рух молекул, коли всі напрями швидкостей рівноймовірні · (EN) thermal motion · \book{} т. 1, с. 300
\item \textbf{середня квадратична швидкість} · $\langle v_{\text{кв}}^2\rangle = (v_1^2+\ldots+v_n^2)/n$ · (EN) root-mean-square speed · \book{} т. 1, с. 303
\item \textbf{основне рівняння молекулярно-кінетичної теорії газів} — тиск ідеального газу дорівнює двом третинам середньої кінетичної енергії поступального руху молекул одиниці об'єму газу · $p = \tfrac23 n_0 \tfrac{m_0\langle v_{\text{кв}}^2\rangle}{2}$ (14.6) · (EN) basic equation of kinetic theory · \book{} т. 1, с. 303
\item \textbf{тиск} (газу) — усереднена за часом результуюча сила дії молекул на одиницю площі стінки; статистичний параметр · [Па = Н/м²] · (EN) pressure · \book{} т. 1, с. 301, 304
\item \textbf{термодинамічна рівновага} — стан, у якому припиняються всі макроскопічні зміни; самовільний перехід до неї — \textbf{релаксація}, час — \textbf{час релаксації} · (EN) thermodynamic equilibrium, relaxation · \book{} т. 1, с. 305
\item \textbf{нульове начало термодинаміки} — кожна система при фіксованих зовнішніх умовах з часом самовільно сягає стану термодинамічної рівноваги · (EN) zeroth law · \book{} т. 1, с. 306
\item \textbf{температура} — інтенсивний параметр, що характеризує стан внутрішнього руху частинок рівноважної системи; для ідеального газу $\Theta = \tfrac23\langle\varepsilon\rangle$, $\Theta = kT$ (14.9) · [К] · (EN) temperature · \book{} т. 1, с. 307–308
\item \textbf{стала Больцмана} — множник, що виражає співвідношення між одиницями енергії і термодинамічної температури · $k = 1{,}38\cdot10^{-23}$ Дж/К; $\tfrac{m_0\langle v_{\text{кв}}^2\rangle}{2} = \tfrac32kT$ (14.10) · (EN) Boltzmann constant · \book{} т. 1, с. 308
\item \textbf{термічне рівняння стану} — рівняння, яке в стані теплової рівноваги пов'язує p, V, T · $F(p,V,T)=0$ (14.12); для ідеального газу $pV = nkT$ (14.14) · (EN) thermal equation of state · \book{} т. 1, с. 311–312
\item \textbf{молярна (універсальна) газова стала} — $R = kN_A = 8{,}31$ Дж/(моль·К); чисельно дорівнює роботі розширення одного моля ідеального газу при ізобарному нагріванні на 1 К · (EN) molar gas constant · \book{} т. 1, с. 312–313
\item \textbf{рівняння Клапейрона} · $pV_M = RT$ (14.16) · (EN) Clapeyron equation · \book{} т. 1, с. 312
\item \textbf{рівняння Клапейрона–Менделєєва} · $pV = \tfrac mM RT$ (14.18) · (EN) ideal gas law · \book{} т. 1, с. 312–313
\item \textbf{закон Бойля–Маріотта} (ізотермічний процес, T = const) · $pV = \text{const}$ (14.22) · (EN) Boyle–Mariotte law · \book{} т. 1, с. 314
\item \textbf{закон Гей-Люссака} (ізобарний процес, p = const) · $V/T = \text{const}$ (14.25), $V = V_0(1+\alpha t)$ (14.26), α = 0,003661 К⁻¹ · (EN) Gay-Lussac's law · \book{} т. 1, с. 315
\item \textbf{рівняння ізохорного процесу} (V = const) · $p/T = \text{const}$ (14.28), $p = p_0(1+\beta t)$ (14.29) · (EN) isochoric law (Charles's law) · \book{} т. 1, с. 315–316
\item \textbf{закон Авогадро} — в однакових об'ємах ідеальних газів при однакових температурах і тисках міститься однакове число молекул · (EN) Avogadro's law · \book{} т. 1, с. 316
\item \textbf{внутрішня енергія} — енергія, що залежить від внутрішнього стану фізичної системи; однозначна функція стану; адитивна · $U = \sum U_i$ (16.1) · (EN) internal energy · \book{} т. 1, с. 375–376
\item \textbf{ступені вільності} (syllabus: «ступені свободи») — кількість незалежних координат, потрібних для визначення положення всіх точок системи · $i = i_{\text{пост}}+i_{\text{оберт}}+i_{\text{колив}}$ (16.4) · (EN) degrees of freedom · \book{} т. 1, с. 376–377
\item \textbf{закон рівного розподілу енергії за класичними ступенями вільності} (закон рівного розподілу Больцмана) — у стані теплової рівноваги на кожний класичний ступінь вільності припадає в середньому кінетична енергія ½kT (коливальний — kT) · (EN) equipartition theorem · \book{} т. 1, с. 377
\item \textbf{внутрішня енергія ідеального газу} — сумарна кінетична енергія молекул; функція тільки абсолютної температури · $U = \tfrac i2RT$ (16.3) на 1 моль · (EN) internal energy of an ideal gas · \book{} т. 1, с. 377, 390–392
\end{itemize}

\needspace{5\baselineskip}\subsection*{Блок B. Розподіл Максвелла · Больцман, Максвелл–Больцман}

\begin{itemize}
\item \textbf{функція розподілу молекул за швидкостями} (густина ймовірності) — $dW = dn/n = f(\vec v)\,d\omega$ (14.42); умова нормування $\int f\,d\omega = 1$ (14.43) · (EN) velocity distribution function · \book{} т. 1, с. 322
\item \textbf{закон розподілу Максвелла} (syllabus: «Максвела») — відносна кількість молекул зі швидкостями від v до v + dv · $\frac{dn}{n} = \frac4{\sqrt\pi}\left(\frac{m_0}{2kT}\right)^{3/2}e^{-m_0v^2/2kT}v^2dv$ (14.57) · (EN) Maxwell distribution · \book{} т. 1, с. 327
\item \textbf{найбільш імовірна швидкість} · $v_\text{н} = \sqrt{2kT/m_0}$ (14.58) · (EN) most probable speed · \book{} т. 1, с. 328
\item \textbf{середня (арифметична) швидкість} · $\langle v\rangle = \sqrt{8kT/\pi m_0}$ (14.61) · (EN) mean speed · \book{} т. 1, с. 330
\item \textbf{середня квадратична (теплова) швидкість} · $v_{\text{кв}} = \sqrt{3kT/m_0}$ (14.62); $v_{\text{кв}} : \langle v\rangle : v_\text{н} \approx 1{,}2 : 1{,}1 : 1$ · (EN) rms speed · \book{} т. 1, с. 318, 330
\item \textbf{барометрична формула} · $p = p_0e^{-Mg(h-h_0)/RT}$ (14.68); T = const, g = const · (EN) barometric formula · \book{} т. 1, с. 332
\item \textbf{розподіл Больцмана} — універсальний закон розподілу частинок у будь-якому потенціальному полі в стані теплової рівноваги · $n = n_0e^{-\Delta U/kT}$ (14.71) · (EN) Boltzmann distribution · \book{} т. 1, с. 335
\item \textbf{розподіл Максвелла–Больцмана} — добуток больцманівського й максвеллівського множників · (14.74) · (EN) Maxwell–Boltzmann distribution · \book{} т. 1, с. 336
\end{itemize}

\needspace{5\baselineskip}\subsection*{Блок C. Довжина вільного пробігу · Явища переносу}

\begin{itemize}
\item \textbf{середня довжина вільного пробігу} — середня відстань, яку проходить молекула між двома послідовними зіткненнями · $\langle l\rangle = 1/(\sqrt2\pi d^2n_0)$ (15.2) · (EN) mean free path · \book{} т. 1, с. 343–344
\item \textbf{середнє число зіткнень} · $\langle z\rangle = \sqrt2\pi d^2n_0\langle v\rangle$ (15.1) · (EN) collision frequency · \book{} т. 1, с. 344
\item \textbf{ефективний поперечний переріз} (повний) · $\sigma = \pi d^2$; характеризує ймовірність процесу, а не геометрію · (EN) effective cross-section · \book{} т. 1, с. 344
\item \textbf{явища переносу} — зумовлені тепловим рухом молекул явища перенесення маси (дифузія), імпульсу (внутрішнє тертя) і енергії (теплопровідність) · (EN) transport phenomena · \book{} т. 1, с. 348
\item \textbf{дифузія газів} — процес самовільного вирівнювання густин (концентрацій) у суміші, зумовлений тепловим рухом мікрочастинок; в однокомпонентному середовищі — \textbf{самодифузія} · (EN) diffusion, self-diffusion · \book{} т. 1, с. 356
\item \textbf{закон Фіка} · $j_m = -D\,d\rho/dx$ (15.22); $D = \tfrac13\langle v\rangle\langle l\rangle$ (15.27) · [м²/с] · (EN) Fick's law · \book{} т. 1, с. 356–357
\item \textbf{внутрішнє тертя (в'язкість)} — гальмівні сили між шарами газу, зумовлені обміном імпульсом впорядкованого руху · закон Ньютона $F = -\eta\frac{du}{dx}\Delta S$ (15.10); $\eta = \tfrac13\rho\langle v\rangle\langle l\rangle$ (15.14) · [Па·с] · (EN) viscosity · \book{} т. 1, с. 348–350
\item \textbf{теплопровідність} — молекулярне перенесення енергії в суцільному середовищі, зумовлене градієнтом температури · закон Фур'є $j_q = -\lambda\,\text{grad}\,T$ (15.16); $\lambda = \tfrac13\rho\langle v\rangle\langle l\rangle c_V$ (15.21) · [Вт/(м·К)] · (EN) thermal conductivity · \book{} т. 1, с. 352–354
\end{itemize}

\needspace{5\baselineskip}\subsection*{Блок D–E. U · Q, A, C · I закон · C газів · Ізопроцеси · Адіабата}

\begin{itemize}
\item \textbf{термодинаміка} — наука про найбільш загальні властивості макроскопічних систем у стані термодинамічної рівноваги та про процеси переходу між цими станами · (EN) thermodynamics · \book{} т. 1, с. 371
\item \textbf{термодинамічна система}; ізольована / закрита / відкрита / адіабатна · (EN) thermodynamic system · \book{} т. 1, с. 371, 373
\item \textbf{квазістатичний (квазірівноважний) процес} — ідеалізований процес із неперервної послідовності рівноважних станів · (EN) quasistatic process · \book{} т. 1, с. 372
\item \textbf{функція стану / функція процесу} — U, S — функції стану (dU, dS); A, Q — функції процесу (d'A, d'Q) · (EN) state function / path function · \book{} т. 1, с. 375, 378–380
\item \textbf{калоричне рівняння стану} — залежність внутрішньої енергії від параметрів стану · U = U(V, T) · (EN) caloric equation of state · \book{} т. 1, с. 376
\item \textbf{робота газу} — макроскопічна міра зміни внутрішньої енергії у квазістатичних процесах · $d\text{'}A = p\,dV$ (16.5), $A = \int p\,dV$ (16.6) · (EN) work · \book{} т. 1, с. 378
\item \textbf{теплообмін, кількість теплоти} — передача енергії від більш нагрітого тіла до менш нагрітого без виконання макроскопічної роботи; Q — її міра · [Дж] · (EN) heat transfer, heat · \book{} т. 1, с. 380
\item \textbf{теплоємність} · $C_0 = d\text{'}Q/dT$ (16.7); \textbf{питома} c: $C_0 = mc$ (16.8); \textbf{молярна} C = Mc (16.9) · (EN) heat capacity · \book{} т. 1, с. 380–381
\item \textbf{перший закон термодинаміки} — у теплових процесах зміна внутрішньої енергії відбувається за рахунок приросту кількості теплоти і виконаної роботи · $dU = d\text{'}Q - d\text{'}A$ (16.10), $Q = \Delta U + A$ (16.12) · (EN) first law of thermodynamics · \book{} т. 1, с. 381–382
\item \textbf{вічний двигун першого роду} — машина, що виконувала б роботу без одержання енергії ззовні; неможливий · (EN) perpetual motion machine of the first kind · \book{} т. 1, с. 382
\item \textbf{рівняння Майєра} · $C_p = C_V + R$ (16.19) · (EN) Mayer's relation · \book{} т. 1, с. 383
\item \textbf{показник адіабати (коефіцієнт Пуассона)} · $\gamma = C_p/C_V$ · (EN) adiabatic index · \book{} т. 1, с. 384
\item \textbf{ізохорний процес} (syllabus: «ізохоричний») · $A = 0$, $Q = \Delta U$ (16.26)–(16.28) · (EN) isochoric process · \book{} т. 1, с. 386–387
\item \textbf{ізобарний процес} · $A = p(V_2-V_1) = \nu R\Delta T$ (16.29), $Q = \nu C_p\Delta T$ (16.30) · (EN) isobaric process · \book{} т. 1, с. 387
\item \textbf{ізотермічний процес} · $\Delta U = 0$, $Q = A = \nu RT\ln(V_2/V_1)$ (16.32) · (EN) isothermal process · \book{} т. 1, с. 387–388
\item \textbf{адіабатний процес} (syllabus: «адіабатичний») — процес без теплообміну з навколишнім середовищем · $d\text{'}A = -dU$ (16.20) · (EN) adiabatic process · \book{} т. 1, с. 383–384
\item \textbf{рівняння Пуассона} · $TV^{\gamma-1} = \text{const}$ (16.23), $pV^\gamma = \text{const}$ (16.23a) · (EN) Poisson's (adiabat) equation · \book{} т. 1, с. 384
\item \textbf{політропний процес} — процес в ідеальному газі зі сталою теплоємністю · $pV^n = \text{const}$ (16.40), $n = (C_p - C)/(C_V - C)$ (16.37) · (EN) polytropic process · \book{} т. 1, с. 388–389
\end{itemize}

\needspace{5\baselineskip}\subsection*{Блок F. Колові процеси · Карно · II закон · Ентропія · Нернст}

\begin{itemize}
\item \textbf{оборотний процес} — процес, при якому система, проведена в зворотному напрямі, повертається в початковий стан через ті самі проміжні стани, а стан тіл поза системою залишається незмінним · (EN) reversible process · \book{} т. 1, с. 397
\item \textbf{необоротний процес} — процес, який не можна здійснити у зворотному напрямі без відповідних змін у системі або оточенні; всі реальні процеси — необоротні · (EN) irreversible process · \book{} т. 1, с. 397
\item \textbf{цикл (коловий процес)} — процес, у результаті якого система повертається в початковий стан; робота = площа циклу на p–V · (EN) cycle · \book{} т. 1, с. 398
\item \textbf{прямий / обернений цикл} — цикл теплової машини (за годинниковою стрілкою) / холодильної машини (проти) · (EN) direct / reverse cycle · \book{} т. 1, с. 399
\item \textbf{цикл Карно} — оборотний цикл з двох ізотерм і двох адіабат; цикл ідеальної теплової машини · $\eta = 1 - Q_2/Q_1$ (16.53) $= 1 - T_2/T_1$ (16.64) · (EN) Carnot cycle · \book{} т. 1, с. 399–401
\item \textbf{холодильний коефіцієнт} · $\varepsilon = Q_2/A$ (16.68) · (EN) coefficient of performance (refrigerator) · \book{} т. 1, с. 403
\item \textbf{коефіцієнт перетворення (теплового насоса)} · $\varphi = Q_1/A = T_1/(T_1-T_2)$ (16.67) · (EN) heat pump COP · \book{} т. 1, с. 402–403
\item \textbf{другий закон термодинаміки} — Клаузіус: теплота не може сама собою переходити від менш нагрітого тіла до більш нагрітого; Томсон: неможлива періодична машина, що перетворює теплоту в роботу лише охолодженням одного тіла · (EN) second law of thermodynamics · \book{} т. 1, с. 403–404
\item \textbf{вічний двигун другого роду} — неможливий (Оствальд) · (EN) perpetual motion machine of the second kind · \book{} т. 1, с. 404
\item \textbf{теорема Карно} — ККД оборотного циклу Карно не залежить від робочого тіла; ККД необоротного циклу завжди менший · $\eta \le \eta_0$ (16.69) · (EN) Carnot's theorem · \book{} т. 1, с. 404
\item \textbf{зведена теплота}; \textbf{нерівність Клаузіуса} · $Q_1/T_1 + Q_2/T_2 \le 0$ (16.71), $\oint d\text{'}Q/T \le 0$ (16.74) · (EN) reduced heat, Clausius inequality · \book{} т. 1, с. 404–405
\item \textbf{ентропія} — функція стану, повний диференціал якої при оборотному процесі $dS = d\text{'}Q/T$ (16.75); адитивна; закон зростання $dS \ge 0$ для ізольованої системи (16.82) · [Дж/К] · (EN) entropy · \book{} т. 1, с. 406–408
\item \textbf{термодинамічна ймовірність} — число мікростанів, які реалізують даний макростан · (16.90) · (EN) thermodynamic probability (statistical weight) · \book{} т. 1, с. 411–412
\item \textbf{принцип Больцмана} · $S = k\ln w$ (16.91) · (EN) Boltzmann's entropy formula · \book{} т. 1, с. 413
\item \textbf{третій закон термодинаміки (теорема Нернста)} — при T → 0 ентропія рівноважної системи перестає залежати від параметрів і прямує до сталої (нуля); $\lim_{T\to0}S = 0$ (16.93, Планк) · (EN) third law, Nernst theorem · \book{} т. 1, с. 414–415
\end{itemize}

\needspace{5\baselineskip}\subsection*{Блок G–H. Реальний газ · Ван-дер-Ваальс · Кристали · Дефекти · Рідини}

\begin{itemize}
\item \textbf{реальний газ} — газ, для якого суттєву роль відіграють міжмолекулярні сили взаємодії · (EN) real gas · \book{} т. 1, с. 419
\item \textbf{сили Ван дер Ваальса} — орієнтаційні (Кеєзом), індукційні (Дебай), дисперсійні (Лондон) · (EN) van der Waals forces · \book{} т. 1, с. 421–422
\item \textbf{рівняння Ван-дер-Ваальса} · $(p + a/V_M^2)(V_M - b) = RT$ (17.4); a/V² — \textbf{внутрішній тиск}, b — поправка на власний об'єм · (EN) van der Waals equation · \book{} т. 1, с. 423
\item \textbf{критичний стан}; \textbf{критичні параметри} · $p_\text{к} = a/27b^2$, $V_\text{к} = 3b$, $T_\text{к} = 8a/27Rb$ (17.12) · (EN) critical state · \book{} т. 1, с. 426
\item \textbf{метастабільні стани} — пересичена пара, перегріта рідина · (EN) metastable states · \book{} т. 1, с. 429
\item \textbf{бінодаль / спінодаль} — межа двофазної області / межа нестійких станів · (EN) binodal / spinodal · \book{} т. 1, с. 429–430
\item \textbf{тверде тіло} — агрегатний стан зі стабільністю форми та об'єму і впорядкованим розміщенням атомів · (EN) solid · \book{} т. 1, с. 444
\item \textbf{дальній / ближній порядок} — строга повторюваність структури на великих відстанях (кристали) / узгодження лише сусідніх частинок (аморфні тіла, рідини) · (EN) long-range / short-range order · \book{} т. 1, с. 444, 463
\item \textbf{анізотропія} — залежність фізичних властивостей від напряму · (EN) anisotropy · \book{} т. 1, с. 445
\item \textbf{елементарна комірка; гратки Браве; сингонія} — 14 граток Браве, 7 сингоній · $\vec r_n = m\vec a_1 + n\vec a_2 + p\vec a_3$ (18.1) · (EN) unit cell, Bravais lattices, crystal system · \book{} т. 1, с. 446–447
\item \textbf{точкові дефекти} — за Шотткі (вакансія), за Френкелем (вакансія + міжвузловий атом), домішки · $n_{\text{в}} = n_0e^{-W_{\text{в}}/kT}$ (18.4) · (EN) point defects · \book{} т. 1, с. 453–455
\item \textbf{дислокація} (крайова, гвинтова); \textbf{вектор Бюргерса} · (EN) dislocation, Burgers vector · \book{} т. 1, с. 455–456
\item \textbf{рідина} — «розупорядковане тверде тіло» з ближнім порядком; тепловий рух за Френкелем: коливання + стрибки · $\tau = Ae^{\Delta W/kT}$ (19.1) · (EN) liquid · \book{} т. 1, с. 462–464
\item \textbf{формула Френкеля (в'язкість рідин)} · $\eta = Ae^{\Delta W/kT}$ (19.7), спадає з T · (EN) Frenkel viscosity formula · \book{} т. 1, с. 467
\end{itemize}


\end{document}
